% main.tex -- the paper. Compile with: pdflatex main (three passes).
\documentclass[11pt]{article}
% preamble.tex -- packages, theorem environments, and every macro of the paper.

\usepackage[T1]{fontenc}
\usepackage[utf8]{inputenc}
\usepackage{lmodern}
\usepackage[margin=1.05in]{geometry}
\usepackage{amsmath,amssymb,amsthm,mathtools}
\usepackage{booktabs,array,longtable}
\usepackage[shortlabels]{enumitem}
\usepackage[dvipsnames]{xcolor}
\usepackage{microtype}
\usepackage[colorlinks=true,linkcolor=blue!50!black,citecolor=green!35!black,urlcolor=blue!50!black]{hyperref}

% ---------------------------------------------------------------------------
% Theorem environments: one shared counter, numbered within sections.
% ---------------------------------------------------------------------------
\theoremstyle{plain}
\newtheorem{theorem}{Theorem}[section]
\newtheorem{lemma}[theorem]{Lemma}
\newtheorem{proposition}[theorem]{Proposition}
\newtheorem{corollary}[theorem]{Corollary}
\newtheorem*{theorem*}{Theorem}
\theoremstyle{definition}
\newtheorem{definition}[theorem]{Definition}
\newtheorem*{definition*}{Definition}
\newtheorem{assumption}[theorem]{Assumption}
\theoremstyle{remark}
\newtheorem{remark}[theorem]{Remark}
\numberwithin{equation}{section}

% ---------------------------------------------------------------------------
% General
% ---------------------------------------------------------------------------
\newcommand{\N}{\mathbb{N}}
\newcommand{\Z}{\mathbb{Z}}
\newcommand{\R}{\mathbb{R}}
\newcommand{\C}{\mathbb{C}}
\newcommand{\one}{\mathbf{1}}
\newcommand{\dd}{\mathrm{d}}
\newcommand{\iu}{\mathrm{i}}
\DeclareMathOperator{\Tr}{Tr}
\DeclareMathOperator{\tr}{tr}
\DeclareMathOperator{\ran}{ran}
\DeclareMathOperator{\spec}{spec}
\DeclareMathOperator{\sgn}{sgn}
\DeclarePairedDelimiter{\norm}{\lVert}{\rVert}
\DeclarePairedDelimiter{\abs}{\lvert}{\rvert}
\DeclarePairedDelimiterX{\ip}[2]{\langle}{\rangle}{#1,#2}
\newcommand{\ket}[1]{\lvert #1\rangle}
\newcommand{\bra}[1]{\langle #1\rvert}
\newcommand{\kb}[2]{\lvert #1\rangle\langle #2\rvert}

% ---------------------------------------------------------------------------
% The torus and its lowest Landau level (Section 2)
% ---------------------------------------------------------------------------
\newcommand{\Tor}{\mathbb{T}}          % the torus R^2 / Lambda
\newcommand{\Nphi}{N_{\phi}}           % number of flux quanta
\newcommand{\h}{\mathfrak{h}}          % one-particle lowest Landau level of the torus
\newcommand{\hpl}{\mathfrak{h}_{\mathrm{pl}}} % planar lowest Landau level
\newcommand{\Fock}{\mathcal{F}}        % spinless Fock space over \h
\newcommand{\Ff}{\mathcal{F}}          % second quantization F(u)
\newcommand{\thx}{\theta_x}            % psi_j -> psi_{j+1}
\newcommand{\thy}{\theta_y}            % psi_j -> e^{2 pi i j/N_phi} psi_j
\newcommand{\Lam}{\Lambda}             % lattice of periods
\newcommand{\cell}{\mathsf{C}}         % a period cell of Lambda
\newcommand{\LL}{\mathcal{L}}          % lifts L_k
\newcommand{\Pp}{\mathfrak{P}}         % periodized image P_k
\newcommand{\Hb}{\mathcal{B}}          % Bargmann space
\newcommand{\Dis}{\mathcal{D}}         % displacement operator D(beta)
\newcommand{\Ec}{\mathcal{E}}          % lattice sums E_m
\newcommand{\U}{\mathfrak{u}}          % centre-of-mass torus lowest Landau level u_k
\newcommand{\Pk}{\Pi}                  % periodization of planar functions
\newcommand{\Fib}{\mathcal{I}}         % fibering maps I_a (Lemma 2.9(c))
\newcommand{\Ppl}{P_{\mathrm{pl}}}     % kernel of the orthogonal projection onto the planar LLL
\newcommand{\tl}{\tilde t}             % lattice magnetic translations, a representation of Lambda

% the Hamiltonian, pairs, constants
\newcommand{\HT}{H_{\Tor}}             % the torus V_1 Hamiltonian
\newcommand{\KT}{K_{\Tor}}             % the comparison operator
\newcommand{\gT}{\gamma_{\Tor}}        % the torus constant
\newcommand{\gs}{\gamma_{*}}           % the planar constant
\newcommand{\gl}{\underline{\gamma}}   % the uniform lower bound
\newcommand{\etas}{\eta_{*}}           % the certificate constant sum_nu lambda_nu^2
\newcommand{\VT}{V_{\Tor}}             % range of h
\newcommand{\bcen}[1]{b_{#1}}          % centre pair operators (orbital form)
\newcommand{\Dlev}{\Delta}             % 2(L^2/4-1)e^{-L^2/8}
\newcommand{\Ups}{\Upsilon}            % remainder of the pair Gram levels
\newcommand{\dG}{\delta_G}             % deviation of the pair Gram matrix G from the identity
\newcommand{\dtail}{\delta_{\mathrm{tail}}}
\newcommand{\Cth}{C_{3}}               % 72 sum ||alpha||^2
\newcommand{\ut}{\tilde u}             % periodized pair states
\newcommand{\eps}{\varepsilon}         % the allowance of the certificate, 3/10^6
\newcommand{\etr}{\epsilon_{\mathrm{tr}}} % the error of the transfer principle


\title{An explicit uniform spectral gap for the Laughlin state parent Hamiltonian on rectangular tori}
\author{A multi-agent AI team,\\ in collaboration with Barry Bradlyn}
\date{Preprint \quad---\quad 23 October 2026}

\begin{document}
\maketitle

\begin{abstract}
We consider spinless fermions in the lowest Landau level of a rectangular torus with periods $L_x$, $L_y$,
interacting through the untruncated $V_1$ pseudopotential, normalized so that every pair of relative angular
momentum one pays energy one. We prove a transfer principle: every planar certificate of the form used by OpenAI on the
sphere, with constant $\gamma$, gives $\HT^2\ge(\gamma-\etr(\Tor))\HT$ for the Hamiltonian $\HT$ on the whole Fock
space of every rectangular torus, where $\etr(\Tor)$ is bounded by an explicit function of the shorter period $L$
that tends to zero at the Gaussian rate $O(L^4\,2^{-L^2/20})$. For OpenAI's planar certificate the bounds are explicit: $\HT^2\ge\gT\HT$, with a
constant $\gT$ that depends on the torus only through its lattice of periods, and for every rectangular torus whose
shorter period is at least $L$, with $L\ge16$ (in units of the magnetic length), one has $\gT\ge\gl(L)$, where $\gl$ is
nondecreasing and tends at the same rate to $\gs=4616733319001/10^{14}$, the constant of OpenAI's certificate;
numerically $\gl(L)>0$ for $L\ge16.457$, $\gl(L)>1/25$ for $L\ge16.857$ and $\gl(L)\ge\gs-10^{-5}$ for $L\ge18.685$.
The bound is uniform in the particle number and in the aspect ratio. Hence the spectrum of $\HT$ in every particle
sector lies in $\{0\}\cup[\gT,\infty)$. At $\Nphi=3N$ flux quanta the kernel on $N$ particles is spanned by the three
centre-of-mass copies of the Laughlin state, and no sector with more particles has a zero mode; the same
three-dimensional count holds, for every modular parameter, for the explicit theta-function pair operators
$B_n^{(\tau)}$, whose identification with the $V_1$ Hamiltonian of an oblique torus is not proved here. The proof
transcribes OpenAI's spherical argument to the torus and consumes the planar certificate as data. On the torus the
averaging over rotations is replaced by an average over magnetic translations, and torus Moyal identities make the
pair-level and four-body parts of the comparison operator exact. The three-body coefficient of the normal-ordered
square of $\HT$ is exactly the periodized planar three-body target, so the three-body target transfers without error;
the comparison operator's three-body part carries one non-orthonormality term, bounded by $\Cth\Ec_8(1+\Ec_8)\HT$. The
remaining errors are Gaussian in the shorter period.
\end{abstract}

\setcounter{tocdepth}{2}
\tableofcontents

% sections/01-introduction.tex
\section{Introduction}\label{sec:intro}

We prove that spectral-gap certificates for the Laughlin $V_1$ model pass from the plane to rectangular tori with explicit
errors. The main result is a transfer principle (Theorem~\ref{thm:transferintro}): every planar certificate of the form introduced
by OpenAI~\cite{OAI:A-Fock-space-inequality-and-the-Laughlin-spectral-gap-September-24-2026}, with constant $\gamma$, gives
$\HT^2\ge(\gamma-\etr(\Tor))\HT$ on the whole Fock space for the untruncated $V_1$ Hamiltonian $\HT$ of every rectangular
torus, where $\etr(\Tor)$ is bounded by an explicit function of the shorter period that tends to zero at a Gaussian rate. For
the certificate of~\cite{OAI:A-Fock-space-inequality-and-the-Laughlin-spectral-gap-September-24-2026} we evaluate these bounds.
On every rectangular torus whose shorter period is at least $L$, with $L\ge16$, the inequality $\HT^2\ge\gT\HT$ holds with a
constant $\gT\ge\gl(L)$, uniformly in the particle number and in the aspect ratio, where $\gl$ is an explicit nondecreasing
function that tends to the planar constant $\gs$
of~\cite{OAI:A-Fock-space-inequality-and-the-Laughlin-spectral-gap-September-24-2026} at a Gaussian rate in $L$; both
$\gs-\gl(L)$ and the bound on $\etr(\Tor)$ are $O(L^4\,2^{-L^2/20})$ (Remark~\ref{rem:rate}).
The thresholds are explicit: a numerical evaluation of $\gl$ (Remark~\ref{rem:evaluation}) gives $\gl(L)>0$ for $L\ge16.457$,
$\gl(L)>1/25$ for $L\ge16.857$ and $\gl(L)\ge\gs-10^{-5}$ for $L\ge18.685$, in units of the magnetic length
(Theorem~\ref{thm:openai}).

Laughlin's wave function~\cite{Laughlin} is the exact zero-energy ground state of Haldane's $V_1$ pseudopotential
Hamiltonian~\cite{Haldane}, in which a pair of spinless fermions in the lowest Landau level pays energy one if its relative
angular momentum is one and nothing otherwise. Whether this Hamiltonian has a spectral gap above its zero modes that stays
bounded below as the system grows is a long-standing question~\cite[Appendix~A]{Rougerie}. For the untruncated interaction it
was answered on the sphere and in the plane by OpenAI~\cite{OAI:A-Fock-space-inequality-and-the-Laughlin-spectral-gap-September-24-2026}:
with $\gs=4616733319001/10^{14}>1/25$, the spherical Hamiltonian $H_Q$ satisfies $H_Q^2\ge\gamma H_Q$ on the whole Fock space for every
$\gamma<\gs$ and all sufficiently large flux $Q$, and the planar Hamiltonian satisfies $(H^{\mathrm{pl}}_N)^2\ge\gs H^{\mathrm{pl}}_N$ for every
particle number. The proof compares $H_Q^2$ with a positive operator built from seven explicit rows, averaged over rotations; the
rows satisfy finitely many planar inequalities, verified by exact rational arithmetic. Uniform gaps for truncated versions of the
pseudopotential were obtained earlier, on thin cylinders and thin tori~\cite{NWY,WY}. Two further works treat the untruncated
interaction on the sphere: Gao~\cite{Gao} raised the constant to about $0.219$, on the sphere and in the plane, by a semidefinite
search, and Scheer and Khalaf~\cite{ScheerKhalaf} constructed a semidefinite hierarchy of lower bounds on the sphere, with exact
rational certificates and a rigorous lower bound on the limit inferior of the gap as the flux grows, and leave the torus to
future work. In all three spherical results the flux threshold is existential: it is shown to exist, not computed.

A uniform bound of this kind, for every flat torus whose shortest lattice vector exceeds an existential threshold, has also
been claimed in a recent candidate proof prepared by Hou with advice from Wu~\cite{HouWu}; the present argument, obtained
independently, is restricted to rectangular tori but makes the threshold, the constant $\gl(L)$ and its Gaussian rate of
convergence to $\gs$ explicit. That candidate proof, which is unrefereed and not on arXiv, states the bound $H^2\ge\gamma H$ for
every $\gamma<\gs$, uniformly in the particle number, on tori of every lattice shape and with every boundary twist, and also
claims that the kernel at filling one third is three-dimensional for every shape. Where the two overlap --- the uniform
bound on rectangular tori and the three-dimensional kernel --- the statements agree. It transfers fixed certificates and states
no principle for a class of certificates; Theorem~\ref{thm:transferintro} below is such a principle, with its error explicit.

\subsection{The transfer principle}\label{ssec:transfer}

Let $\Tor$ be a rectangular torus with periods $L_x,L_y$ (in units of the magnetic length) and $\Nphi=L_xL_y/2\pi$ flux quanta,
let $\h$ be its lowest Landau level and $\Fock$ the fermionic Fock space over $\h$, and let $\HT$ be the untruncated $V_1$
Hamiltonian: the second quantization of the periodization, over the lattice of periods, of the planar projector onto pairs of
relative angular momentum one (Definition~\ref{def:HT}). Let $L:=\min\{L_x,L_y\}$ be the shorter period.

The certificate of~\cite{OAI:A-Fock-space-inequality-and-the-Laughlin-spectral-gap-September-24-2026} is a finite family of rows
together with planar inequalities that they satisfy. On the torus, at a base point $R$, row $\nu$ is the operator
\[
F_\nu(R)=\lambda_\nu\,B(\ut_{t_\nu}(R))+\sum_{(p,j)\in\mathcal S_\nu}\alpha^\nu_{pj}\,c(\chi^R_{p+j-t_\nu})^*\,c(\chi_j^R)\,B(\ut_p(R))
\]
of Definition~\ref{def:KT}, where $B(\ut_p(R))$ removes a pair in the periodized pair state of label $p$ about $R$ and $c(\chi^R_m)$ a
particle in the periodized orbital of degree $m$ about $R$ (Section~\ref{sec:local}). The proof uses the certificate only through
the data of the rows and the inequalities, which define the following class.

\begin{definition*}[planar certificates of the same form]
A \emph{planar certificate of the same form} as that
of~\cite{OAI:A-Fock-space-inequality-and-the-Laughlin-spectral-gap-September-24-2026} consists of finitely many rows $\nu$, each with an
integer $t_\nu\ge0$, a real $\lambda_\nu$ and real coefficients $\alpha^\nu_{pj}$ supported in a finite set
$\mathcal S_\nu\subset\{(p,j):0\le p\le P_c,\ 0\le j\le M_c,\ 0\le p+j-t_\nu\le M_c\}$, where $P_c,M_c\ge0$ are integers, and an
allowance $\eps\ge0$, such that, with the cutoffs $Z_c:=\max\{p+j:(p,j)\in\bigcup_\nu\mathcal S_\nu\}$ and
$D_c:=\max\{p+j+(q+l-t_\nu)\}$ over $(p,j),(q,l)\in\mathcal S_\nu$, the inequalities
\begin{equation}\label{eq:certgeneral}
e_z\le\tfrac32q_z\ \ (z\le Z_c,\ z\notin\{0,1,3\}),\qquad G_D(2\kb{e_1}{e_1}-E_D+\eps\one)G_D\ge0\ \ (1\le D\le D_c)
\end{equation}
hold. Here $q_z=(3z-1)(-\frac12)^z$ are the planar three-body coefficients (Lemma~\ref{lem:qz}); $e_z$ and $E_D$ are the three-body
scalars and the four-body copy matrices of the rows, defined by \eqref{eq:ez} and \eqref{eq:ED} from the planar coupling
coefficients (the sum over $T$ in \eqref{eq:ez} running to $Z_c$, and \eqref{eq:ED} taken for $1\le D\le D_c$); $G_D$ is the Gram
matrix of the planar four-body relative states $W_4\omega_{D,r}$, $r\in\mathcal R_D$, defined with them in Section~\ref{sec:moyal}; $\mathcal R_D$
is the set of odd integers in $[1,D]$; and $\kb{e_1}{e_1}$ is the matrix with $(1,1)$ entry $1$ and all other entries zero. The
\emph{constant} of the certificate is
\[
\gamma:=1-\etas-\delta_\infty(Z_c)-\eps\,N_4,\qquad \etas=\sum_\nu\lambda_\nu^2,
\]
\[
\delta_\infty(Z_c):=\frac32\sum_{\substack{z>Z_c\\ z\ \mathrm{odd}}}\frac{(3z-1)(z+1)}{2^z},\qquad N_4:=\sum_{D=1}^{D_c}\Bigl\lfloor\frac{(D+1)^2}4\Bigr\rfloor .
\]
We require that some $\mathcal S_\nu$ be nonempty, that $M_c\ge1$ and that $D_c\ge2$, so that $Z_c$ is defined and the bounds of
Lemma~\ref{lem:Ebounds}(b) are used only for $m\ge1$; the certificate of~\cite{OAI:A-Fock-space-inequality-and-the-Laughlin-spectral-gap-September-24-2026}
has $M_c=8$ and $D_c=23$.
\end{definition*}

Every row of this form ends in a pair annihilator and has a definite weight $t_\nu+1$ (each of its summands removes two particles
and lowers the total orbital degree by $t_\nu+1$), and these two features are all that the proof uses of a row; that is why the
principle is a property of the form, not of the numbers in the certificate
of~\cite{OAI:A-Fock-space-inequality-and-the-Laughlin-spectral-gap-September-24-2026}.

\begin{theorem}[the transfer principle]\label{thm:transferintro}
Given a planar certificate of the same form with constant $\gamma$, for every rectangular torus, on $\Fock$,
\[
\HT^2\ge(\gamma-\etr(\Tor))\,\HT,
\]
where, with $\Ec_m=\Ec_m(\Tor)$ and $\sigma_z=\sigma_z(\Tor)$,
\begin{multline*}
\etr(\Tor):=(1-\lambda_{\min})+\frac32\sum_{\substack{z>Z_c\\ z\ \mathrm{odd}}}\frac{(3z-1)(z+1)}{2^z}\sigma_z
+(M_c+1)(P_c+1)\sum_\nu\norm{\alpha^\nu}^2\,\Ec_{M_c}(1+\Ec_{M_c})\\
+\eps N_4\bigl((1+\Ec_{D_c})^2-1\bigr)+2N_4\,\Ec_{D_c-1}(1+\Ec_{D_c})^2 .
\end{multline*}
Here $\lambda_{\min}$ is the smallest eigenvalue of the Gram matrix of the torus pair functions (Lemma~\ref{lem:levels}),
$\sigma_z(\Tor)=\sup_R\max_{m\le z}\bigl(\norm{\chi^R_m}^2-1\bigr)$ (Lemma~\ref{lem:Ebounds}(c)), $\Ec_m(\Tor)$ is the lattice sum, over
the nonzero periods, of the norms of the truncated displacement matrices $\Dis^{(m)}$ (Section~\ref{sec:local}), and
$\norm{\alpha^\nu}^2=\sum_{p,j}(\alpha^\nu_{pj})^2$. For every rectangular torus whose shorter period is at least $L$, $\etr(\Tor)\le\etr(L)$
with $\etr(L)\to0$ as $L\to\infty$. In particular, if $\gamma>0$, then for every $\gamma'<\gamma$ there is $L_0$ such that $\HT^2\ge\gamma'\HT$
on the whole Fock space of every rectangular torus whose shorter period is at least $L_0$, uniformly in the particle number and the
aspect ratio.
\end{theorem}

The function $\etr(L)$, obtained in the proof from explicit bounds on the terms of $\etr(\Tor)$, is $O(L^4\,2^{-L^2/20})$ as
$L\to\infty$, with a constant that depends on the certificate (Remark~\ref{rem:rate}). The theorem does not require the planar
certificate to come from a spherical proof: any finite family of rows of this form satisfying \eqref{eq:certgeneral} with $\gamma>0$
yields a gap on rectangular tori, uniform above a threshold period, with that rate.

For the certificate of~\cite{OAI:A-Fock-space-inequality-and-the-Laughlin-spectral-gap-September-24-2026}
(Section~\ref{sec:moyal} and Appendix~\ref{app:certificate}), $P_c=7$, $M_c=8$, $Z_c=15$, $D_c=23$, $\delta_\infty(15)=61/4096$ and
$N_4=1222$, with $\etas=93527408868499/10^{14}$ and $\eps=3/10^6$; its constant is $\gamma=\gs=4616733319001/10^{14}>1/25$, and
$\gs-\etr(\Tor)$ is, term by term, the constant $\gT$ of \eqref{eq:gT}. Theorem~\ref{thm:transferintro} for this certificate, together
with an evaluation of $\etr$ for it, gives the following explicit bounds.

\begin{theorem}[OpenAI's certificate on rectangular tori]\label{thm:openai}
For every rectangular torus $\Tor$, on the whole Fock space,
\[
\HT^2\ge\gT\,\HT,
\]
with the explicit constant $\gT$ of \eqref{eq:gT}, which depends on the torus only through its lattice of periods. For $L\ge16$ let
$\gl(L)$ be the explicit function of Corollary~\ref{cor:uniform}. Then $\gT\ge\gl(L)$ for every rectangular torus whose shorter period
is at least $L$, so that $\HT^2\ge\gl(L)\HT$ uniformly in the particle number and in the aspect ratio; $\gl$ is nondecreasing on
$[16,\infty)$, and $\gl(L)\to\gs$ as $L\to\infty$, with $\gs-\gl(L)=O(L^4\,2^{-L^2/20})$. Numerically, $\gl(L)>0$ for $L\ge16.457$,
$\gl(L)>1/25$ for $L\ge16.857$ and $\gl(L)\ge\gs-10^{-5}$ for $L\ge18.685$, in units of the magnetic length; these values are
floating-point evaluations of rigorous closed-form bounds, not interval-arithmetic certificates (Remark~\ref{rem:evaluation}).
\end{theorem}

Theorem~\ref{thm:openai} is thus the case of Theorem~\ref{thm:transferintro} for the certificate
of~\cite{OAI:A-Fock-space-inequality-and-the-Laughlin-spectral-gap-September-24-2026}, with $\etr$ evaluated. Section~\ref{sec:theorem}
proves the two in the opposite order: Theorem~\ref{thm:main} and Corollary~\ref{cor:uniform}, with Remark~\ref{rem:rate} and the
thresholds evaluated in Section~\ref{ssec:constants}, prove Theorem~\ref{thm:openai} directly, and Section~\ref{ssec:transferproof} then
proves Theorem~\ref{thm:transferintro}, restated as Theorem~\ref{thm:transfer}, by the same argument with the data of the certificate
left general.

\begin{remark}[scope]\label{rem:scope}
Theorem~\ref{thm:transferintro} is proved for the form of the definition: rows each of whose summands is a pair annihilator, alone
or dressed by one creation and one annihilation operator; the two-body $V_1$ interaction of a single species of spinless fermions;
and rectangular tori without boundary twists. Three natural extensions are not covered. Rows of higher degree in the fermion
operators have squares with parts on five or more particles; they would need the torus Moyal identity on that many particles,
which Theorem~\ref{thm:moyal} states and proves for every number of particles, and a comparison term for those parts, which is not
constructed here. Two species, as in the spinful hollow-core model of Section~\ref{sec:outlook}, and general lattices of periods
with boundary twists are not treated. Whether these extensions hold is not claimed here.
\end{remark}

By Theorem~\ref{thm:openai}, the spectrum of $\HT$ in every particle sector lies in $\{0\}\cup[\gT,\infty)$. At $\Nphi=3N$ the kernel on
$N$ particles is three-dimensional, spanned by the three centre-of-mass copies of the Laughlin state, and no sector with more than
$N$ particles has a zero mode (Theorem~\ref{thm:kernel}, Proposition~\ref{prop:overfill}); the same count holds, for every modular
parameter, for the explicit theta-function pair operators $B_n^{(\tau)}$ of \eqref{eq:Bntau}, whose identification with the $V_1$
Hamiltonian of an oblique torus is not proved here. So $\HT\ge\gT(\one-P_N)$ on $N$ particles, with $P_N$ of rank three, and
$\HT\ge\gT$ on more particles (Corollary~\ref{cor:spectrum}).

At the square aspect ratio $L=\sqrt{2\pi\Nphi}$, so the bound is positive for $\Nphi\ge44$ and exceeds $1/25$ for $\Nphi\ge46$. At
filling one third ($\Nphi=3N$) this means a gap for $N\ge15$ ($\gl(L)\ge0.0385$ at $N=15$) and a gap larger than $1/25$ for $N\ge16$
($\gl(L)\ge0.0456$ at $N=16$). For elongated tori the condition is on the shorter period, so larger systems are needed. All error terms
decay like $e^{-cL^2}$ in the shorter period; the binding one at the thresholds is a lattice sum of norms of truncated
displacement matrices, bounded by a multiple of $L^{44}e^{-L^2/4}$ (Lemma~\ref{lem:Ebounds}), and the thresholds are where it
becomes small. The slowest-decaying bound is that on the three-body tail beyond the blocks covered by the certificate,
$O(L^4\,2^{-L^2/20})$ (Proposition~\ref{prop:three}); it sets the rate of Remark~\ref{rem:rate}.

\subsection{From the sphere to the torus}

The torus proof follows the architecture
of~\cite{OAI:A-Fock-space-inequality-and-the-Laughlin-spectral-gap-September-24-2026}. One expands $\HT^2$ in normal order into
two-, three- and four-body terms, and compares them with the corresponding terms of a positive comparison operator
$\KT=\frac1\pi\int_\Tor\sum_\nu F_\nu(R)^*F_\nu(R)\,\dd R$, where the rows $F_\nu(R)$ of Section~\ref{ssec:transfer} are the planar
rows of the certificate, those of~\cite{OAI:A-Fock-space-inequality-and-the-Laughlin-spectral-gap-September-24-2026} for
Theorem~\ref{thm:openai}, built at the base point $R$ from periodized orbitals and pair states. The planar certificate is used as
data, without change: for that of~\cite{OAI:A-Fock-space-inequality-and-the-Laughlin-spectral-gap-September-24-2026}, the row
coefficients, the thirteen three-body inequalities and the twenty-three four-body matrix inequalities. Section, lemma and equation numbers
of~\cite{OAI:A-Fock-space-inequality-and-the-Laughlin-spectral-gap-September-24-2026} refer to its version of 24 September 2026,
and its lemmas are also cited by name. On the sphere the comparison operator is
averaged over rotations and the finite-flux corrections are controlled by a relative-error transfer lemma; on the torus the
average over the base point $R$ replaces the rotations, and the corrections take a different and simpler form. Four facts,
proved here, make the transcription work:
\begin{enumerate}[(1)]
\item \emph{The projector replacement.} On two particles $\HT$ is not a projection. The Gram matrix of the torus pair functions is
computed in closed form as a lattice sum of image overlaps (Lemma~\ref{lem:gram}); its eigenvalues are
$1\pm\Dlev(L_y)\pm\Dlev(L_x)$ for even $\Nphi$ and $1\pm(\Dlev(L_x)^2+\Dlev(L_y)^2)^{1/2}$ for odd $\Nphi$, up to a smaller remainder, with
$\Dlev(\ell)=2(\ell^2/4-1)e^{-\ell^2/8}$ (Lemma~\ref{lem:levels}). Only the smallest eigenvalue enters the budget.
\item \emph{Periodized local modes.} The periodized orbitals about a point have an explicit Gram matrix, a lattice sum of
displacement matrices with phases (Lemma~\ref{lem:localgram}), whose deviation from the identity is bounded uniformly by lattice
sums of norms of truncated displacement matrices (Lemma~\ref{lem:Ebounds}).
\item \emph{Torus Moyal identities.} Averaging a projector onto periodized $k$-particle states with a fixed relative part and a
centre-of-mass orbital about $R$ over $R\in\Tor$ gives exactly the periodized image of the relative projector, and different
centre-of-mass degrees average to zero (Theorem~\ref{thm:moyal}). This replaces Schur's lemma for $SU(2)$; the proof uses Schur's
lemma for the magnetic translations of the centre-of-mass torus. With it, the pair-level part of $\KT$ is exactly $\etas\HT$, and
its three- and four-body parts are exactly the periodized images of the planar partial traces that define the certificate,
up to one small non-orthonormality term (Proposition~\ref{prop:parts}).
\item \emph{The square of $\HT$.} The three-body coefficient of $\HT^2$ is a sum of products of two pair operators sharing exactly
one particle, and periodization is multiplicative on such products (Lemma~\ref{lem:product}). Hence it is \emph{exactly} the
periodized planar three-body target, with the planar coefficients $q_z=(3z-1)(-\frac12)^z$ on every block
(Theorem~\ref{thm:square}). The three-body target therefore transfers without error, and the tail beyond the blocks covered by
the certificate is bounded with an explicit rate (Proposition~\ref{prop:three}). The four-body coefficient involves products on
disjoint pairs, and a Bessel-type bound for the localized channels of the relative coordinate of the two pair centres closes
the budget (Lemma~\ref{lem:bessel}).
\end{enumerate}
The resulting error terms are all relative to $\HT$ (they vanish on its kernel and carry no factor of the particle number), and all
are Gaussian in the shorter period. They are collected in the error $\etr(\Tor)$ of Theorem~\ref{thm:transferintro}, and for the
certificate of~\cite{OAI:A-Fock-space-inequality-and-the-Laughlin-spectral-gap-September-24-2026} in the constant $\gT$
of \eqref{eq:gT}.

Section~\ref{sec:kernel} determines the kernel at $\Nphi=3N$, for the pair operators $B_n^{(\tau)}$ of \eqref{eq:Bntau} at every
modular parameter and hence for $\HT$ on every rectangular torus. The threefold degeneracy of the torus Laughlin states is
classical~\cite{HaldaneRezayiTorus}; that these states are the only zero modes of the torus $V_1$ Hamiltonian is stated for
rectangular tori in~\cite[Section~II.2]{ONDS}, and a proof for every lattice shape and boundary twist is claimed
in~\cite{HouWu}. The argument of Section~\ref{sec:kernel} uses third-order vanishing at coincidences, division by the cube of
the theta-function Jastrow factor, and Liouville's theorem.

\subsection{Structure of the paper}

Section~\ref{sec:setting} sets up the torus, the planar lowest Landau level, the periodization and the periodized image of planar
operators, and defines $\HT$; Proposition~\ref{prop:orbital} gives its orbital form. Section~\ref{sec:kernel} determines the
kernel. Section~\ref{sec:projector} computes the pair Gram matrix and records the normal-ordered square of a lifted pair
operator. Section~\ref{sec:local} introduces the periodized local modes and bounds their Gram matrix. Section~\ref{sec:moyal}
proves the torus Moyal identities, states the certificate, and decomposes the comparison operator. Section~\ref{sec:square}
proves the exact three-body identity and the three-body comparison, and Section~\ref{sec:bessel} the four-body Bessel bound.
Section~\ref{sec:theorem} proves Theorem~\ref{thm:transferintro} and evaluates it for the certificate
of~\cite{OAI:A-Fock-space-inequality-and-the-Laughlin-spectral-gap-September-24-2026}: it assembles the budget and evaluates the
constants for that certificate, which proves Theorem~\ref{thm:openai}, derives the spectral consequences, and then proves
Theorem~\ref{thm:transferintro} by the same argument with the data of the certificate left general.
Section~\ref{sec:outlook} discusses the transfer principle as a basis for further work and lists what is not proved.
Appendix~\ref{app:certificate} records the certificate data, and Appendix~\ref{app:notation} lists the notation.

% sections/02-setting.tex
\section{Setting}\label{sec:setting}

We use units in which the magnetic field and the magnetic length are equal to one ($\hbar=e=1$). For
$a=(a_x,a_y)$ and $b=(b_x,b_y)$ in $\R^2$ we write $a\wedge b:=a_xb_y-a_yb_x$, and for $r=(x,y)\in\R^2$ we write
$w:=x+\iu y$ for the corresponding complex coordinate. Inner products are linear in the second argument.

\subsection{The torus and its lowest Landau level}

\begin{definition}[the torus and its lowest Landau level]\label{def:torus}
Let $\Nphi\in\N$ and $L_x,L_y>0$ with $L_xL_y=2\pi\Nphi$; put $\kappa:=2\pi/L_y$, so that $L_x=\kappa\Nphi$. Let
$\Lam:=L_x\Z\times L_y\Z$, $\Tor:=\R^2/\Lam$, $\cell:=[0,L_x)\times[0,L_y)$, and let $L:=\min\{L_x,L_y\}$ be the
\emph{shorter period}. Let
\[
D_x:=\partial_x,\qquad D_y:=\partial_y-\iu x
\]
be the covariant derivatives of the vector potential $(0,x)$. A smooth function $f$ on $\R^2$ is an \emph{LLL
function} if $(D_x+\iu D_y)f=0$. The \emph{lowest Landau level of the torus}, $\h$, is the space of LLL functions with
\begin{equation}\label{eq:bc}
f(x,y+L_y)=f(x,y),\qquad f(x+L_x,y)=e^{\iu L_xy}f(x,y),
\end{equation}
with the inner product $\ip fg:=\int_\cell\overline f\,g\,\dd r$.
\end{definition}

An LLL function is exactly a function of the form $f=e^{-x^2/2}g(w)$ with $g$ entire, since
$(D_x+\iu D_y)(e^{-x^2/2}g)=e^{-x^2/2}(\partial_x+\iu\partial_y)g$. The functions $\abs f^2$ and $\overline fg$ are
$\Lam$-periodic for $f,g\in\h$, so the inner product is an integral over $\Tor$. All results concern these
boundary conditions, without Aharonov--Bohm phases through the cycles.

\begin{lemma}[Landau-gauge orbitals]\label{lem:orbitals}
For $n\in\Z$ let $\phi_n(x,y):=(L_y\sqrt\pi)^{-1/2}e^{\iu\kappa ny}e^{-(x-\kappa n)^2/2}$ and for $j\in\Z_{\Nphi}$ let
$\psi_j:=\sum_{m\in\Z}\phi_{\tilde\jmath+m\Nphi}$, where $\tilde\jmath\in\Z$ is any representative of $j$. Write $[n]$
for the class of $n\in\Z$ in $\Z_{\Nphi}$.
\begin{enumerate}[(i)]
\item The series converges absolutely and locally uniformly, $\psi_j$ does not depend on $\tilde\jmath$, and
$\psi_j\in\h$.
\item $(\psi_j)_{j\in\Z_{\Nphi}}$ is an orthonormal basis of $\h$; in particular $\dim\h=\Nphi$, and every $f\in\h$ is
bounded.
\end{enumerate}
\end{lemma}

\begin{proof}
(i) $\abs{\phi_n(x,y)}=(L_y\sqrt\pi)^{-1/2}e^{-(x-\kappa n)^2/2}$, and these Gaussians are summable over
$n\in\tilde\jmath+\Nphi\Z$ locally uniformly in $x$. Each $\phi_n$ is $e^{-x^2/2}$ times an entire function of $w$
(namely $(L_y\sqrt\pi)^{-1/2}e^{\kappa nw-\kappa^2n^2/2}$), so $e^{x^2/2}\psi_j$ is a locally uniform limit of
entire functions, hence entire, and $\psi_j$ is an LLL function. Replacing $\tilde\jmath$ by $\tilde\jmath+\Nphi$
relabels the sum. Since $\kappa L_y=2\pi$ each $\phi_n$ is $L_y$-periodic in $y$, and since $\kappa\Nphi=L_x$,
$\phi_n(x+L_x,y)=e^{\iu L_xy}\phi_{n-\Nphi}(x,y)$; summing over the class gives \eqref{eq:bc}.

(ii) For $n,n'\in\Z$ the $y$-integral of $\overline{\phi_n}\phi_{n'}$ over $[0,L_y)$ is
$\delta_{nn'}\pi^{-1/2}e^{-(x-\kappa n)^2}$. Hence $\ip{\psi_i}{\psi_j}=0$ for $i\neq j$, and
$\ip{\psi_j}{\psi_j}=\sum_m\int_0^{L_x}\pi^{-1/2}e^{-(x-\kappa\tilde\jmath-mL_x)^2}\dd x=1$. Let $f\in\h$. By the first
condition in \eqref{eq:bc}, $f(x,y)=\sum_nf_n(x)e^{\iu\kappa ny}$ with smooth $f_n$; the LLL condition reads
$f_n'+(x-\kappa n)f_n=0$, so $f_n=a_ne^{-(x-\kappa n)^2/2}$, and the second condition in \eqref{eq:bc} gives
$a_n=a_{n-\Nphi}$. Hence $f=(L_y\sqrt\pi)^{1/2}\sum_{j=0}^{\Nphi-1}a_j\psi_j$. Each $\psi_j$ is bounded by (i).
\end{proof}

\subsection{The planar lowest Landau level and magnetic translations}

Let $\hpl$ be the Hilbert space of square-integrable LLL functions on $\R^2$. For $a\in\R^2$ the \emph{magnetic
translation} $t_a$ acts on functions on $\R^2$ by
\begin{equation}\label{eq:ta}
(t_af)(r):=e^{-\iu a_xa_y/2}\,e^{\iu a_xy}\,f(r-a).
\end{equation}
For $\gamma=(jL_x,kL_y)\in\Lam$ put $\tl_\gamma:=(-1)^{jk\Nphi}t_\gamma$, and on the torus put
$\thx:=t_{(\kappa,0)}$ and $\thy:=t_{(0,-L_y/\Nphi)}$, that is $(\thx f)(x,y)=e^{\iu\kappa y}f(x-\kappa,y)$ and
$(\thy f)(x,y)=f(x,y+L_y/\Nphi)$.

\begin{lemma}[magnetic translations]\label{lem:magtrans}
\begin{enumerate}[(i)]
\item $t_a$ maps LLL functions to LLL functions, $\abs{(t_af)(r)}=\abs{f(r-a)}$, and $t_a$ is unitary on $\hpl$
with $t_a^*=t_{-a}$.
\item $t_at_b=e^{\iu(a\wedge b)/2}\,t_{a+b}$, hence $t_at_b=e^{\iu\,a\wedge b}\,t_bt_a$.
\item $\gamma\mapsto\tl_\gamma$ is a representation of $\Lam$, and $\h$ is the space of LLL functions $f$ with
$\tl_\gamma f=f$ for all $\gamma\in\Lam$.
\item For $\delta\in\Lam/\Nphi:=(L_x/\Nphi)\Z\times(L_y/\Nphi)\Z$, $t_\delta$ maps $\h$ onto $\h$ unitarily. In particular
$\thx\psi_j=\psi_{j+1}$, $\thy\psi_j=e^{2\pi\iu j/\Nphi}\psi_j$, $\thy\thx=e^{2\pi\iu/\Nphi}\thx\thy$ and
$\thx^{\Nphi}=\thy^{\Nphi}=\one$ on $\h$.
\end{enumerate}
\end{lemma}

\begin{proof}
(i) $t_a$ is multiplication by a unimodular function composed with a translation, so it preserves absolute values
and $L^2$ norms; a direct computation gives $(D_x+\iu D_y)t_af=t_a(D_x+\iu D_y)f$, and $t_at_{-a}=\one$ by (ii).
(ii) Expanding \eqref{eq:ta} twice, $t_at_bf(r)$ has the phase
$-\tfrac12a_xa_y-\tfrac12b_xb_y+(a_x+b_x)y-b_xa_y$ and $t_{a+b}f(r)$ has the phase
$-\tfrac12(a_x+b_x)(a_y+b_y)+(a_x+b_x)y$; the difference is $\tfrac12(a_xb_y-a_yb_x)$. (iii) For $\gamma=(jL_x,kL_y)$
and $\gamma'=(j'L_x,k'L_y)$, $\gamma\wedge\gamma'=2\pi\Nphi(jk'-kj')$, so by (ii)
$\tl_\gamma\tl_{\gamma'}=(-1)^{(jk+j'k'+jk'-kj')\Nphi}t_{\gamma+\gamma'}=\tl_{\gamma+\gamma'}$, since
$jk+j'k'+jk'-kj'\equiv(j+j')(k+k')\pmod 2$. The two conditions \eqref{eq:bc} are $\tl_{(0,L_y)}f=f$ and
$\tl_{(L_x,0)}f=f$. (iv) By (ii), $t_\delta\tl_\gamma=e^{\iu\,\delta\wedge\gamma}\tl_\gamma t_\delta$ with
$\delta\wedge\gamma\in2\pi\Z$ for $\delta\in\Lam/\Nphi$, so $t_\delta$ preserves $\h$; it is isometric for the inner
product of $\h$ by (i), because $\abs{t_\delta f}^2$ is a translate of the $\Lam$-periodic function $\abs f^2$. The
formulas for $\thx$, $\thy$ follow from $\thx\phi_n=\phi_{n+1}$ and $\thy\phi_n=e^{2\pi\iu n/\Nphi}\phi_n$.
\end{proof}

\paragraph{Bargmann form.}
Let $\Hb$ be the Bargmann space of entire functions $g$ with
\[
\norm g_\Hb^2:=\frac1\pi\int_\C\abs{g(z)}^2e^{-\abs z^2}\dd^2z<\infty ;
\]
the monomials $z^m/\sqrt{m!}$ form an orthonormal basis, written $\ket m$, and multiplication by $z$ and $\partial_z$ are
mutually adjoint. For $\beta\in\C$
the \emph{displacement operator} is
\begin{equation}\label{eq:displacement}
(\Dis(\beta)g)(z):=e^{-\abs\beta^2/2}e^{\beta z}g(z-\bar\beta),
\end{equation}
that is $\Dis(\beta)=e^{-\abs\beta^2/2}e^{\beta z}e^{-\bar\beta\partial_z}$, and for $m\ge0$ we write
$\Dis^{(m)}(\beta):=[\bra a\Dis(\beta)\ket b]_{0\le a,b\le m}$ for its truncation to the first $m+1$ levels.
$\mathrm L_n^{(\alpha)}$ denotes the generalized Laguerre polynomial.

\begin{lemma}[symmetric-gauge orbitals and displacements]\label{lem:bargmann}
\begin{enumerate}[(i)]
\item The map $\mathcal U:\Hb\to\hpl$ defined by
\[
(\mathcal U g)(r):=(2\pi)^{-1/2}e^{\iu xy/2}g(w/\sqrt2)e^{-\abs w^2/4}
\]
is unitary. Put $\eta_m:=\mathcal U(z^m/\sqrt{m!})$, so that
\[
\eta_m(r)=\frac{e^{\iu xy/2}\,w^m\,e^{-\abs w^2/4}}{\sqrt{2\pi\,2^m\,m!}},\qquad m\ge0 .
\]
\item $\mathcal U^{-1}t_a\mathcal U=\Dis(\beta(a))$ with $\beta(a):=(a_x-\iu a_y)/\sqrt2$, so that $\abs{\beta(a)}^2=\abs a^2/2$ and
$\ip{\eta_m}{t_a\eta_n}=\bra m\Dis(\beta(a))\ket n$.
\item For $m\ge n$ and $x:=\abs\beta^2$,
\begin{gather*}
\bra m\Dis(\beta)\ket n=\sqrt{\tfrac{n!}{m!}}\,\beta^{m-n}e^{-x/2}\,\mathrm L_n^{(m-n)}(x),\\
\bra n\Dis(\beta)\ket m=\sqrt{\tfrac{n!}{m!}}\,(-\bar\beta)^{m-n}e^{-x/2}\,\mathrm L_n^{(m-n)}(x).
\end{gather*}
In particular $\abs{\bra m\Dis(\beta)\ket n}\le1$, and $\norm{\Dis^{(m)}(\beta)}$ depends only on $\abs\beta$.
\item The orthogonal projection of $L^2(\R^2)$ onto $\hpl$ has the kernel
$\Ppl(r,r')=\sum_m\eta_m(r)\overline{\eta_m(r')}$, with $\abs{\Ppl(r,r')}=(2\pi)^{-1}e^{-\abs{r-r'}^2/4}$, and
$\int \Ppl(r,r')f(r')\dd r'=f(r)$ for every bounded LLL function $f$. Every $f\in\hpl$ satisfies
$\abs{f(r)}\le(2\pi)^{-1/2}\norm f$.
\end{enumerate}
\end{lemma}

\begin{proof}
(i) Write $z=w/\sqrt2$. Then $\abs{\mathcal U g}^2=(2\pi)^{-1}\abs{g(z)}^2e^{-\abs z^2}$ and $\dd r=2\,\dd^2z$, so $\mathcal U$ is
isometric. Since $e^{\iu xy/2-\abs w^2/4}=e^{-x^2/2}e^{w^2/4}$, $\mathcal U g$ is an LLL function; conversely an LLL function
$e^{-x^2/2}\tilde g(w)$ equals $\mathcal U g$ with $g(z)=(2\pi)^{1/2}e^{-z^2/2}\tilde g(\sqrt2z)$, and square integrability is
$g\in\Hb$. (ii) Put $b:=(a_x+\iu a_y)/\sqrt2$. Then $e^{-\abs{z-b}^2/2}=e^{-\abs z^2/2}e^{(\bar zb+z\bar b)/2-\abs b^2/2}$,
and the phases in $t_a\mathcal U g$ combine to $e^{\iu xy/2}e^{\iu(a_xy-a_yx)/2}$ with $\iu(a_xy-a_yx)/2=(z\bar b-\bar zb)/2$.
Hence $\mathcal U^{-1}t_a\mathcal U g(z)=e^{z\bar b-\abs b^2/2}g(z-b)$, which is \eqref{eq:displacement} with $\beta=\bar b$.
(iii) By the Baker--Campbell--Hausdorff formula ($[\beta z,-\bar\beta\partial_z]=\abs\beta^2$),
$\Dis(\beta)=e^{-x/2}e^{\beta z}e^{-\bar\beta\partial_z}$, and
\[
\bra m\Dis(\beta)\ket n=e^{-x/2}\sum_{k\le n}\frac{\sqrt{m!\,n!}\;\beta^{m-k}(-\bar\beta)^{n-k}}{(m-k)!\,(n-k)!\,k!}
=\sqrt{\tfrac{n!}{m!}}\,\beta^{m-n}e^{-x/2}\sum_{i=0}^n\frac{m!\,(-x)^i}{(m-n+i)!\,i!\,(n-i)!},
\]
with $i=n-k$; the last sum is $\mathrm L_n^{(m-n)}(x)$. The second formula follows from
$\bra n\Dis(\beta)\ket m=\overline{\bra m\Dis(-\beta)\ket n}$, since $\Dis(\beta)^*=\Dis(-\beta)$ by (i)--(ii) and
Lemma~\ref{lem:magtrans}(i). The bound $\abs{\bra m\Dis\ket n}\le1$ holds because $\Dis(\beta)$ is unitary. By the two
formulas, $\bra a\Dis(e^{\iu\theta}\beta)\ket b=e^{\iu\theta(a-b)}\bra a\Dis(\beta)\ket b$ for real $\theta$, so the truncations
at $\beta$ and $e^{\iu\theta}\beta$ are conjugate by a diagonal unitary.

(iv) $\sum_m\eta_m(r)\overline{\eta_m(r')}=(2\pi)^{-1}e^{\iu(xy-x'y')/2}e^{z\bar z'-(\abs z^2+\abs{z'}^2)/2}$, of modulus
$(2\pi)^{-1}e^{-\abs{z-z'}^2/2}$, and $\abs{z-z'}^2=\abs{r-r'}^2/2$. For $f=\mathcal U g$ bounded, $\abs{g(z)}\le Ce^{\abs z^2/2}$, and
$\int \Ppl(r,r')f(r')\dd r'=(\mathcal U h)(r)$ with $h(z)=\pi^{-1}\int g(z')e^{z\bar z'}e^{-\abs{z'}^2}\dd^2z'$. The substitution
$z'=z+u$ gives $h(z)=\pi^{-1}\int g(z+u)e^{-\bar zu}e^{-\abs u^2}\dd^2u$, an absolutely convergent integral; in polar
coordinates the mean value of the entire function $u\mapsto g(z+u)e^{-\bar zu}$ over each circle is $g(z)$, so
$h(z)=g(z)$. For $f=\mathcal U g\in\hpl$, that is $g\in\Hb$, both integrals converge absolutely by Cauchy--Schwarz, since
$z'\mapsto g(z')e^{-\abs{z'}^2/2}$ and $z'\mapsto e^{z\bar z'}e^{-\abs{z'}^2/2}$ are square-integrable, and the same computation
proves the reproducing identity on $\hpl$. Finally, for $f\in\hpl$, $\abs{f(r)}=\abs{\ip{\Ppl(\cdot,r)}f}\le \Ppl(r,r)^{1/2}\norm f$.
\end{proof}

For $R=(R_x,R_y)\in\R^2$ and $m\ge0$ we put
\begin{equation}\label{eq:etaR}
\eta_m^R:=e^{-\iu R_xR_y/2}\,t_R\,\eta_m,\qquad
\eta_m^R(r)=\frac{e^{\iu(y-R_y)(x+R_x)/2}\,w_R^m\,e^{-\abs{w_R}^2/4}}{\sqrt{2\pi\,2^m\,m!}},\quad w_R:=w-(R_x+\iu R_y),
\end{equation}
the symmetric-gauge orbitals about the point $R$, written in the gauge $(0,x)$. They form an orthonormal basis of
$\hpl$ for every $R$.

\paragraph{The Landau representation.}
For $s\in\R$ let $\phi^\delta_s(x,y):=(2\pi)^{-1/2}\pi^{-1/4}e^{\iu sy}e^{-(x-s)^2/2}$ ($\delta$-normalized Landau orbitals)
and let $\mathcal H_m(s):=(2^mm!\sqrt\pi)^{-1/2}H_m(s)e^{-s^2/2}$ be the normalized Hermite functions. A direct Gaussian integration with the generating functions of
$\eta_m$ and $\mathcal H_m$ gives
\begin{equation}\label{eq:landaurep}
\ip{\phi^\delta_s}{\eta_m^R}_{\R^2}=e^{-\iu sR_y}\,\mathcal H_m(s-R_x),
\end{equation}
the integral being absolutely convergent. For $R=0$ use the generating function
\[
\sum_m\frac{t^m}{m!}\sqrt{2\pi2^mm!}\,\eta_m=e^{\iu xy/2}e^{tw-\abs w^2/4};
\]
the $y$-integral and then the $x$-integral are Gaussian:
\[
\int_{\R^2}e^{-\iu sy-(x-s)^2/2}\,e^{\iu xy/2}e^{tw-\abs w^2/4}\,\dd r=2\pi e^{-s^2/2+2ts-t^2}
=2\pi e^{-s^2/2}\sum_mH_m(s)\frac{t^m}{m!},
\]
and comparing coefficients gives \eqref{eq:landaurep}; the case of general $R$ follows from
$t_R\phi^\delta_s=e^{-\iu R_y(s+R_x/2)}\phi^\delta_{s+R_x}$. Since $\phi_n=\sqrt\kappa\,\phi^\delta_{\kappa n}$,
\[
\ip{\phi_n}{\eta_m^R}_{\R^2}=\sqrt\kappa\,e^{-\iu\kappa nR_y}\mathcal H_m(\kappa n-R_x).
\]

\subsection{Periodization and unfolding}

For $k\ge1$ we realize $\h^{\otimes k}$ as the space of functions on $(\R^2)^k$ that are LLL functions of each
variable and invariant under $\tl^{(i)}_\gamma$ (the operator $\tl_\gamma$ acting on the $i$-th variable) for every
$i$ and every $\gamma\in\Lam$, with the inner product $\int_{\cell^k}$. This space is spanned by the products
$\psi_{j_1}(r_1)\cdots\psi_{j_k}(r_k)$: expanding such a function $\Psi$ in its first variable in the basis $(\psi_j)$ of
Lemma~\ref{lem:orbitals}, the coefficients $\ip{\psi_j}{\Psi(\cdot,r_2,\dots,r_k)}$ are functions of the same kind in $k-1$ variables
(Morera's theorem), and one inducts. In particular its elements are bounded. Permutations $P_\sigma$ of the
variables act on it unitarily.

\begin{lemma}[periodization and unfolding]\label{lem:unfold}
Let $f$ be a function on $(\R^2)^k$ which is an LLL function of each variable and satisfies
$\abs{f(r)}\le Ce^{-c\sum_i\abs{r_i-r_0}^2}$ for some $C,c>0$ and $r_0\in\R^2$. Then
\[
\Pk^{(k)}f:=\sum_{\gamma\in\Lam^k}\tl^{(1)}_{\gamma_1}\cdots\tl^{(k)}_{\gamma_k}f
\]
converges absolutely and locally uniformly, $\Pk^{(k)}f\in\h^{\otimes k}$, and for every $\Psi\in\h^{\otimes k}$
\begin{equation}\label{eq:unfold}
\ip{\Pk^{(k)}f}{\Psi}=\int_{(\R^2)^k}\overline f\,\Psi .
\end{equation}
$\Pk^{(k)}$ commutes with permutations of the variables. We write $\Pk:=\Pk^{(1)}$.
\end{lemma}

\begin{proof}
$\abs{\tl^{(i)}_{\gamma_i}f(r)}$ is $\abs f$ translated by $\gamma_i$ in the $i$-th variable, so the series converges
absolutely and locally uniformly by the Gaussian bound. Each term is $\prod_ie^{-x_i^2/2}$ times an entire function of
$(w_1,\dots,w_k)$, so the limit is of the same form and is an LLL function of each variable. Reindexing and
Lemma~\ref{lem:magtrans}(iii) give $\tl^{(i)}_\gamma\Pk^{(k)}f=\Pk^{(k)}f$. For \eqref{eq:unfold}, use
$\Psi=\tl^{(1)}_{\gamma_1}\cdots\tl^{(k)}_{\gamma_k}\Psi$ and the fact that for any two functions $g_1,g_2$,
$\overline{\tl_\gamma g_1}\,\tl_\gamma g_2$ is the translate of $\overline{g_1}g_2$ by $\gamma$; thus the term $\gamma$ of
$\int_{\cell^k}\overline{\Pk^{(k)}f}\Psi$ is $\int_{\cell^k-\gamma}\overline f\Psi$, and the translates of $\cell^k$ tile
$(\R^2)^k$. All rearrangements are justified by absolute convergence, since $\Psi$ is bounded. Permutations commute
with the sum because all variables carry the same boundary conditions.
\end{proof}

\subsection{Fock space and lifts}\label{ssec:fock}

Let $\Fock:=\bigoplus_{n=0}^{\Nphi}\bigwedge^n\h$, with the convention that wedges of orthonormal vectors
$e_{i_1}\wedge\cdots\wedge e_{i_n}$, $i_1<\cdots<i_n$, have norm one. For $v\in\bigwedge^k\h$ let $B(v)$ be the adjoint
of left wedge multiplication by $v$; $B(v)$ is conjugate-linear in $v$ and
\begin{equation}\label{eq:Bwedge}
B(v\wedge w)=B(w)B(v).
\end{equation}
For $f\in\h$, $c(f):=B(f)$ is the annihilation operator; it is conjugate-linear in $f$, satisfies the canonical
anticommutation relations $\{c(f),c(g)^*\}=\ip fg$, $\{c(f),c(g)\}=0$, and $\norm{c(f)}=\norm f$. We put
$c_j:=c(\psi_j)$. For a unitary $u$ on $\h$, $\Ff(u)$ is its second quantization. The map
$\iota_n(f_1\wedge\cdots\wedge f_n):=\sqrt{n!}\,A_n(f_1\otimes\cdots\otimes f_n)$, $A_n:=(n!)^{-1}\sum_\sigma\sgn(\sigma)P_\sigma$,
identifies $\bigwedge^n\h$ isometrically with the antisymmetric elements of $\h^{\otimes n}$; we use it without
further mention, so that $\bigwedge^n\h$ is the space of antisymmetric functions in $\h^{\otimes n}$.

For an operator $X$ on $\bigwedge^k\h$ its \emph{lift} is defined by linear extension of
\begin{equation}\label{eq:lift}
\LL_k(\kb vw):=B(v)^*B(w),\qquad v,w\in\textstyle\bigwedge^k\h .
\end{equation}
We extend $X$ to $\h^{\otimes k}$ by $XA_k$; this extension commutes with permutations, and for a $k$-subset $S$ of
$\{1,\dots,n\}$ we write $X_S$ for it acting on the tensor positions in $S$.

\begin{lemma}[lifts]\label{lem:lift}
\begin{enumerate}[(i)]
\item For $v\in\bigwedge^k\h$, $\Phi\in\bigwedge^{n-k}\h$ and $\Psi\in\bigwedge^n\h$:
$v\wedge\Phi=\binom nk^{1/2}A_n(v\otimes\Phi)$ and $B(v)\Psi=\binom nk^{1/2}(\bra v\otimes\one)\Psi$.
\item On $\bigwedge^n\h$, $\LL_k(X)=\sum_{\abs S=k}X_S$ for $n\ge k$, and $\LL_k(X)=0$ for $n<k$; in particular $\LL_k(X)$
restricted to $\bigwedge^k\h$ is $X$.
\item $\LL_k$ is linear, $\LL_k(X)^*=\LL_k(X^*)$, and $X\ge0$ implies $\LL_k(X)\ge0$. Operators $Z_k$ on $\bigwedge^k\h$
with $\sum_{k=0}^{m}\LL_k(Z_k)=0$ vanish.
\item The wedge maps $W_3:\bigwedge^2\h\otimes\h\to\bigwedge^3\h$, $W_3(v\otimes u):=v\wedge u$, and
$W_4:\bigwedge^2\h\otimes\bigwedge^2\h\to\bigwedge^4\h$, $W_4(v\otimes v'):=v\wedge v'$, are $\sqrt3A_3$ and $\sqrt6A_4$
restricted to these subspaces of $\h^{\otimes3}$, $\h^{\otimes4}$.
\end{enumerate}
\end{lemma}

\begin{proof}
(i) For orthonormal basis vectors, $v=e_1\wedge\cdots\wedge e_k$ and $\Phi=e_{k+1}\wedge\cdots\wedge e_n$, the first
identity reads $\sqrt{n!}A_n(e_1\otimes\cdots\otimes e_n)=\binom nk^{1/2}\sqrt{k!(n-k)!}A_n(A_k\otimes A_{n-k})(e_1\otimes\cdots\otimes e_n)$,
which holds since $A_n(A_k\otimes A_{n-k})=A_n$; it extends by multilinearity. Then
$\ip\Phi{B(v)\Psi}=\ip{v\wedge\Phi}\Psi=\binom nk^{1/2}\ip{v\otimes\Phi}{\Psi}$, because $A_n$ is an orthogonal
projection fixing $\Psi$. (ii) By (i), $B(v)^*B(w)\Psi=\binom nkA_n(Y\otimes\one)\Psi$ with $Y=\kb vw$, which commutes with
permutations of its $k$ positions. For antisymmetric $\Psi$,
$A_n(Y\otimes\one)\Psi=(n!)^{-1}\sum_\sigma(P_\sigma(Y\otimes\one)P_\sigma^{-1})\Psi$, and each $k$-subset occurs for
$k!(n-k)!$ permutations; this gives $\binom nk^{-1}\sum_SY_S\Psi$. For $n<k$, $B(w)\Psi=0$. (iii) Linearity and the
adjoint are clear; a positive $X$ is a sum $\sum\kb vv$, and $B(v)^*B(v)\ge0$. If $\sum_k\LL_k(Z_k)=0$, restricting to
$k=0,1,\dots$ particles successively gives $Z_0=0$, $Z_1=0$, and so on, by (ii). (iv) is (i) with $n=3,4$.
\end{proof}

\subsection{Jacobi coordinates and the periodized image}

For $k\ge1$ fix a real orthogonal $k\times k$ matrix $O$ whose first row is $k^{-1/2}(1,\dots,1)$, and let it act on
$(r_1,\dots,r_k)\in(\R^2)^k$ componentwise. The new variables are the \emph{centre-of-mass coordinate}
$\zeta:=k^{-1/2}\sum_ir_i\in\R^2$ and \emph{relative coordinates} $\rho\in(\R^2)^{k-1}$. The specific choices used below
are, for $k=2$, $\rho=(r_1-r_2)/\sqrt2$; for $k=3$, $\rho=(\rho_1,\rho_2)$ with $\rho_1=(r_1-r_2)/\sqrt2$,
$\rho_2=(r_1+r_2-2r_3)/\sqrt6$; for $k=4$, $\rho=(\rho_1,\rho_2,\rho_3)$ with $\rho_1=(r_1-r_2)/\sqrt2$,
$\rho_2=(r_3-r_4)/\sqrt2$, $\rho_3=(r_1+r_2-r_3-r_4)/2$.

\begin{lemma}[Jacobi coordinates]\label{lem:jacobi}
\begin{enumerate}[(i)]
\item A function on $(\R^2)^k$ is an LLL function of each $r_i$ if and only if, written in $(\zeta,\rho)$, it is an LLL
function of each new variable. The change of variables is unitary on $L^2$; it identifies
$\hpl^{\otimes k}$ with $\hpl\otimes\hpl^{\otimes(k-1)}$ (centre of mass $\otimes$ relative).
\item The simultaneous translation $\prod_it^{(i)}_a$ acts as $t_{\sqrt k\,a}$ on $\zeta$ and trivially on $\rho$.
\item A permutation of the particles acts on $(\zeta,\rho)$ by an orthogonal map that fixes $\zeta$.
\end{enumerate}
\end{lemma}

\begin{proof}
(i) With $(\xi,\upsilon)=(Ox,Oy)$, $\partial_{x_i}+\iu\partial_{y_i}+x_i=\sum_jO_{ji}(\partial_{\xi_j}+\iu\partial_{\upsilon_j}+\xi_j)$, and $O$ is
invertible; the Jacobian is one; and $\sum_ix_iy_i$, $\sum_i\abs{w_i}^2$ are invariant, so the Bargmann factorization of
Lemma~\ref{lem:bargmann}(i) is preserved. (ii) $\prod_it_a^{(i)}f(r)=e^{-\iu ka_xa_y/2}e^{\iu a_x\sum_iy_i}f(r-a\one)$; here
$\sum_iy_i=\sqrt k\,\zeta_y$, and $r\mapsto r-a\one$ is $\zeta\mapsto\zeta-\sqrt ka$ with $\rho$ fixed, because the rows of
$O$ other than the first are orthogonal to $(1,\dots,1)$; this is $t_{\sqrt ka}$ in $\zeta$. (iii) A permutation matrix
fixes $(1,\dots,1)$.
\end{proof}

Let $\Tor_k:=\R^2/\sqrt k\Lam$, a rectangular torus with periods $\sqrt kL_x$, $\sqrt kL_y$ and $k\Nphi$ flux quanta, and
let $\U_k$ be its lowest Landau level (Definition~\ref{def:torus} with these periods), the \emph{centre-of-mass
torus space}; by Lemma~\ref{lem:orbitals}, $\dim\U_k=k\Nphi$. Its lattice translations are $\tl'_{\gamma'}$,
$\gamma'\in\sqrt k\Lam$. For $\gamma=(jL_x,lL_y)\in\Lam$, Lemma~\ref{lem:jacobi}(ii) and $(-1)^{jl\Nphi\cdot k}=(-1)^{jl(k\Nphi)}$ show
that the simultaneous lattice translation $\tl^{\otimes k}_\gamma:=\prod_i\tl^{(i)}_\gamma$ acts as $\tl'_{\sqrt k\gamma}$ on $\zeta$ and trivially on
$\rho$.

\begin{definition}[admissible operators, the periodized image]\label{def:admissible}
A \emph{$k$-body admissible operator} is an integral operator $Y$ on functions on $(\R^2)^k$ with a continuous kernel
$Y(r;r')$ such that
\begin{enumerate}[(Y1)]
\item for each $r'$, $Y(\cdot\,;r')$ is an LLL function of each variable, and for each $r$, $\overline{Y(r;\cdot)}$ is;
\item $Y$ commutes with the simultaneous magnetic translations $\prod_it_a^{(i)}$, $a\in\R^2$;
\item $\abs{Y(r;r')}\le C\exp\bigl(-c\sum_i(\abs{r_i-r_1}^2+\abs{r'_i-r_1}^2)\bigr)$ for some $C,c>0$.
\end{enumerate}
For admissible $Y$ the \emph{periodized image} $\Pp_k(Y)$ is the operator on $\h^{\otimes k}$ defined by
\begin{equation}\label{eq:Pk}
\ip\Phi{\Pp_k(Y)\Psi}:=\int_{\mathsf D}\overline{\Phi}\,(Y\Psi),\qquad \Phi,\Psi\in\h^{\otimes k},
\end{equation}
where $\mathsf D\subset(\R^2)^k$ is a fundamental domain for the diagonal action $r\mapsto r+\gamma\one$ of $\Lam$, for
instance $\{r_1\in\cell\}$ or, in Jacobi coordinates, $\{\zeta\in\sqrt k\cell,\ \rho\ \text{arbitrary}\}$.
\end{definition}

\begin{lemma}[the periodized image]\label{lem:Pk}
Let $Y$ be admissible and $\Phi,\Psi\in\h^{\otimes k}$.
\begin{enumerate}[(a)]
\item $Y\Psi(r)=\int Y(r;r')\Psi(r')\dd r'$ converges absolutely, is an LLL function of each variable, satisfies
$\tl^{\otimes k}_\gamma(Y\Psi)=Y\Psi$ for $\gamma\in\Lam$, and $\abs{Y\Psi(r)}\le C'\exp(-c'\sum_i\abs{r_i-r_1}^2)$. The integral
\eqref{eq:Pk} converges absolutely and does not depend on $\mathsf D$.
\item $\Pp_k$ is linear, $\Pp_k(Y)^*=\Pp_k(Y^*)$, and $P_\sigma\Pp_k(Y)P_\sigma^{-1}=\Pp_k(P_\sigma YP_\sigma^{-1})$,
$\Pp_k(Y)P_\sigma=\Pp_k(YP_\sigma)$. In particular $\Pp_k(Y)$ restricted to $\bigwedge^k\h$ equals $\Pp_k(A_kYA_k)$
restricted to $\bigwedge^k\h$.
\item \emph{(Fibering.)} For a Gaussian-decaying relative LLL function $a$ (that is, $\abs{a(\rho)}\le Ce^{-c\abs\rho^2}$) and a
bounded function $\Psi$ on $(\R^2)^k$ put $(\Fib_a\Psi)(\zeta):=\int\overline{a(\rho)}\,\Psi(\zeta,\rho)\,\dd\rho$. Then $\Fib_a$
maps $\h^{\otimes k}$ into $\U_k$; the operator $\kb ab\otimes\one$, with kernel $a(\rho)\overline{b(\rho')}\Ppl(\zeta,\zeta')$,
is admissible; and $\Pp_k(\kb ab\otimes\one)=\Fib_a^*\Fib_b$. Consequently $\Pp_k(Y_{\mathrm{rel}}\otimes\one)\ge0$ for every
positive finite-rank operator $Y_{\mathrm{rel}}$ on the relative variables whose range consists of Gaussian-decaying
relative LLL functions. Moreover, for $g\in\U_k$,
\begin{equation}\label{eq:Fstar}
\Fib_a^*g=\sum_{\gamma\in\Lam^{k-1}}\Bigl(\prod_{i=2}^k\tl^{(i)}_{\gamma_i}\Bigr)\bigl(a\otimes g\bigr),\qquad (a\otimes g)(\zeta,\rho):=a(\rho)g(\zeta),
\end{equation}
the series converging absolutely and locally uniformly.
\end{enumerate}
\end{lemma}

\begin{proof}
(a) The kernel decays in $r'_i-r_1$ for every $i$ and $\Psi$ is bounded, so the integral converges, locally uniformly in
$r$, with the stated Gaussian bound. By (Y1), $e^{x_i^2/2}Y(r;r')$ is an entire function of $w_i$ for fixed values of the other
variables, so by Morera's theorem the locally uniformly convergent integral $e^{x_i^2/2}Y\Psi(r)$ is entire in $w_i$, and $Y\Psi$ is
an LLL function of each variable. By (Y2) and $\tl^{\otimes k}_\gamma=(\pm1)\prod_it^{(i)}_\gamma$,
$\tl^{\otimes k}_\gamma Y\Psi=Y\tl^{\otimes k}_\gamma\Psi=Y\Psi$. The integrand of \eqref{eq:Pk} is therefore invariant under $r\mapsto r+\gamma\one$
(the phases of $\tl^{\otimes k}_\gamma$ cancel in $\overline\Phi\,Y\Psi$) and decays like a Gaussian in the relative coordinates,
so it is integrable over a fundamental domain, and any two fundamental domains give the same integral.
(b) Linearity is clear. For the adjoint, let $G(r,r'):=\overline{\Phi(r)}\,Y(r;r')\,\Psi(r')$. By (Y2), conjugation by
$\tl^{\otimes k}_\gamma$ multiplies the kernel by the phase of $\tl^{\otimes k}_\gamma$ at $r$ and its inverse at $r'$, and $\Phi$, $\Psi$ are
$\tl^{\otimes k}_\gamma$-invariant; hence $G(r+\gamma\one,r'+\gamma\one)=G(r,r')$. By (Y3) and the boundedness of $\Phi,\Psi$, $G$ is
absolutely integrable over $\mathsf D\times(\R^2)^k$ and over $(\R^2)^k\times\mathsf D$, and both integrals equal the
integral of $G$ over a fundamental domain of the joint diagonal action on $(\R^2)^k\times(\R^2)^k$. The first is
$\int_{\mathsf D}\overline\Phi\,Y\Psi$ and the second is $\int_{\mathsf D}\overline{Y^*\Phi}\,\Psi$, where $Y^*$ has the kernel
$\overline{Y(r';r)}$, which is again admissible. For permutations: $P_\sigma$ preserves $\h^{\otimes k}$ and maps
fundamental domains to fundamental domains. The last statement follows from $A_k\Psi=\Psi$ for $\Psi\in\bigwedge^k\h$.
(c) $\Fib_a\Psi$ is an LLL function of $\zeta$ by Lemma~\ref{lem:jacobi}(i) and Morera's theorem as in (a), it is bounded,
and it is $\tl'$-invariant because $\tl^{\otimes k}_\gamma$ acts on $\zeta$ alone; so $\Fib_a\Psi\in\U_k$. The kernel
$a(\rho)\overline{b(\rho')}\Ppl(\zeta,\zeta')$ satisfies (Y1) and (Y2) (Lemma~\ref{lem:jacobi}(ii), and $\Ppl$ commutes with $t_a$),
and (Y3) holds because $\abs\rho^2+\abs{\rho'}^2+\abs{\zeta-\zeta'}^2$ dominates a multiple of
$\sum_i(\abs{r_i-r_1}^2+\abs{r'_i-r_1}^2)$, all these differences being linear combinations of $\rho$, $\rho'$ and
$\zeta-\zeta'$. By Lemma~\ref{lem:bargmann}(iv) applied in $\zeta$ to the bounded LLL function $\Fib_b\Psi$,
$((\kb ab\otimes\one)\Psi)(\zeta,\rho)=a(\rho)(\Fib_b\Psi)(\zeta)$, and with $\mathsf D=\{\zeta\in\sqrt k\cell\}$ the integral
\eqref{eq:Pk} becomes $\int_{\sqrt k\cell}\overline{\Fib_a\Phi}\,\Fib_b\Psi=\ip{\Fib_a\Phi}{\Fib_b\Psi}_{\U_k}$. A positive finite-rank
$Y_{\mathrm{rel}}$ is $\sum_\mu\kb{b_\mu}{b_\mu}$ with $b_\mu$ in its range, and $\Pp_k$ depends only on the kernel of its
argument, so $\Pp_k(Y_{\mathrm{rel}}\otimes\one)=\sum_\mu \Fib_{b_\mu}^*\Fib_{b_\mu}\ge0$. For \eqref{eq:Fstar}: $a\otimes g$ is an LLL function
of each particle, invariant under the diagonal action (because $g\in\U_k$), and decays like a Gaussian in $r_i-r_1$. For
$\Psi\in\h^{\otimes k}$, $\ip\Psi{\Fib_a^*g}=\ip{\Fib_a\Psi}g=\int_{\mathsf D}\overline\Psi\,(a\otimes g)$ with $\mathsf D=\{\zeta\in\sqrt k\cell\}$; the
integrand is invariant under the diagonal action, so we may take $\mathsf D=\{r_1\in\cell\}$ instead, and unfolding the variables
$r_2,\dots,r_k$ into $\cell$ as in Lemma~\ref{lem:unfold} shows that the right side of \eqref{eq:Fstar}, which lies in
$\h^{\otimes k}$ by the argument of Lemma~\ref{lem:unfold}, has the same inner product with every $\Psi$.
\end{proof}

\subsection{The Hamiltonian}

Two-particle LLL functions are, by Lemma~\ref{lem:jacobi}, functions of the relative coordinate $\rho=(r_1-r_2)/\sqrt2$
and of $\zeta=(r_1+r_2)/\sqrt2$, and the relative angular momentum of the pair is the degree in $\rho$. Let
\[
P_1:=\kb{\eta_1}{\eta_1}\otimes\one,
\]
with $\eta_1$ in the relative variable $\rho$: the orthogonal projection of $\hpl^{\otimes2}$ onto the pairs of relative
angular momentum one, with every centre-of-mass state. The planar $N$-particle operator $\sum_{i<j}(P_1)_{ij}$ is the
$V_1$ pseudopotential Hamiltonian in which every pair of relative angular momentum one pays energy one and every other
pair pays nothing, which is the normalization of~\cite{OAI:A-Fock-space-inequality-and-the-Laughlin-spectral-gap-September-24-2026}
and of the planar spectral-gap conjecture~\cite[Appendix~A]{Rougerie}.

\begin{definition}[the torus $V_1$ Hamiltonian]\label{def:HT}
$h:=\Pp_2(P_1)$, an operator on $\h^{\otimes2}$ which preserves $\bigwedge^2\h$ (Lemma~\ref{lem:Pk}(b)); we use the same
letter for its restriction to $\bigwedge^2\h$. The \emph{torus $V_1$ Hamiltonian} is $\HT:=\LL_2(h)$ on $\Fock$. We put
$\VT:=\ran h\subset\bigwedge^2\h$.
\end{definition}

By Lemma~\ref{lem:Pk}(c), $h\ge0$, so $\HT\ge0$; $\HT$ preserves the particle number, and on $n$ particles
$\HT=\sum_{i<j}h_{ij}$ by Lemma~\ref{lem:lift}(ii). The next proposition gives its orbital form; the pair functions
$v_n$ are the image sums of the planar pair states.

\begin{proposition}[orbital form]\label{prop:orbital}
Let $c_0:=\kappa^3/\sqrt{2\pi}$ and $f(d):=d\,e^{-\kappa^2d^2/4}$. For $n\in\Z$ put
\begin{equation}\label{eq:vtau}
v_n:=\sqrt{c_0}\sum_{A+B=n}f(A-B)\,\psi_{[A]}\wedge\psi_{[B]}\in\textstyle\bigwedge^2\h,
\end{equation}
the sum running over $(A,B)\in\Z^2$; it converges absolutely and $v_{n+2\Nphi}=v_n$. Then
\begin{equation}\label{eq:horbital}
h=\sum_{n\in\Z_{2\Nphi}}\kb{v_n}{v_n},\qquad
\HT=c_0\sum_{n\in\Z_{2\Nphi}}\bcen n^*\bcen n,\qquad \bcen n:=\sum_{\substack{d\ge1\\ d\equiv n\ (2)}}2d\,e^{-\kappa^2d^2/4}\,
c_{[(n-d)/2]}\,c_{[(n+d)/2]} .
\end{equation}
Moreover, with $\Fib:=\Fib_{\eta_1}:\h^{\otimes2}\to\U_2$ and $(\psi'_n)_{n\in\Z_{2\Nphi}}$ the orbitals of
$\Tor_2=\R^2/\sqrt2\Lam$ (Lemma~\ref{lem:orbitals} with periods $\sqrt2L_x$, $\sqrt2L_y$), $h=\Fib^*\Fib$ and
\begin{equation}\label{eq:imagesum}
v_n=\Fib^*\psi'_n=\sum_{\gamma\in\Lam}\tl^{(2)}_\gamma\bigl(\eta_1(\rho)\,\psi'_n(\zeta)\bigr),
\end{equation}
the image sum over the second particle of the planar pair state with relative orbital $\eta_1$ and centre-of-mass
state $\psi'_n$ (in the tensor realization of $\bigwedge^2\h$).
\end{proposition}

\begin{proof}
By Lemma~\ref{lem:Pk}(c), $h=\Fib^*\Fib$. For $A,B\in\Z$ write $d=A-B$ and $m=A+B$. In the coordinates
$\rho=(r_1-r_2)/\sqrt2$, $\zeta=(r_1+r_2)/\sqrt2$, orthogonality of the coordinate change gives
\[
\phi_{A}(r_1)\phi_{B}(r_2)=(L_y\sqrt\pi)^{-1}\,e^{\iu\kappa d\rho_y/\sqrt2}e^{-(\rho_x-\kappa d/\sqrt2)^2/2}\;
e^{\iu\kappa m\zeta_y/\sqrt2}e^{-(\zeta_x-\kappa m/\sqrt2)^2/2}.
\]
By \eqref{eq:landaurep} at $R=0$, $\int\overline{\eta_1(\rho)}\,e^{\iu s\rho_y}e^{-(\rho_x-s)^2/2}\dd\rho=(2\pi)^{1/2}\pi^{1/4}\mathcal H_1(s)$,
with $\mathcal H_1(s)=\pi^{-1/4}\sqrt2\,s\,e^{-s^2/2}$. The orbitals $\phi'_m$, $\psi'_n$ of $\Tor_2$ are built with
$\kappa':=\kappa/\sqrt2=2\pi/(\sqrt2L_y)$ and $2\Nphi$ flux quanta, and
$e^{\iu\kappa m\zeta_y/\sqrt2}e^{-(\zeta_x-\kappa m/\sqrt2)^2/2}=(\sqrt2L_y\sqrt\pi)^{1/2}\phi'_m(\zeta)$. Hence
\[
\Fib(\phi_{A}\otimes\phi_{B})=2^{1/4}\sqrt\kappa\;\mathcal H_1(\kappa d/\sqrt2)\,\phi'_m
=2^{1/4}\pi^{-1/4}\kappa^{3/2}f(A-B)\,\phi'_{A+B} .
\]
Since $\sum_{A,B}\abs{\phi_{A}(r_1)\phi_{B}(r_2)}$ is bounded and $\eta_1$ is integrable, dominated convergence gives
$\Fib(\psi_a\otimes\psi_b)=\sum_{A\equiv a,\,B\equiv b}\Fib(\phi_{A}\otimes\phi_{B})$ for $a,b\in\Z_{\Nphi}$. The coefficient of
$\phi'_m$ in this sum depends on $m$ only through its class modulo $2\Nphi$ (shift $A$ and $B$ by $\Nphi$), so, grouping the
$m$ by classes,
\[
\Fib(\psi_a\otimes\psi_b)=2^{1/4}\pi^{-1/4}\kappa^{3/2}\sum_{n\in\Z_{2\Nphi}}C_n(a,b)\,\psi'_n,\qquad
C_n(a,b):=\sum_{\substack{A+B=n\\ A\equiv a,\ B\equiv b}}f(A-B),
\]
for any representative $n$ of the class. As $(\psi'_n)$ is orthonormal and the $C_n$ are real,
$\ip\Phi{h\Psi}=\ip{\Fib\Phi}{\Fib\Psi}=\sum_n\overline{\ip{w_n}\Phi}\,\ip{w_n}\Psi$ with
$w_n:=2^{1/4}\pi^{-1/4}\kappa^{3/2}\sum_{A+B=n}f(A-B)\,\psi_{[A]}\otimes\psi_{[B]}$, and $w_n=\Fib^*\psi'_n$. The tensor
$\sum_{A+B=n}f(A-B)\psi_{[A]}\otimes\psi_{[B]}$ is antisymmetric ($f$ is odd), and it equals $2^{-1/2}$ times the
antisymmetric tensor of $\sum_{A+B=n}f(A-B)\psi_{[A]}\wedge\psi_{[B]}$; since $2^{1/4}\pi^{-1/4}\kappa^{3/2}2^{-1/2}=c_0^{1/2}$,
$w_n=v_n$. This proves $h=\sum_n\kb{v_n}{v_n}$. Then $\HT=\sum_nB(v_n)^*B(v_n)$, and
$B(\psi_{[A]}\wedge\psi_{[B]})=c_{[B]}c_{[A]}$ by \eqref{eq:Bwedge}; pairing $(A,B)$ with $(B,A)$ gives $B(v_n)=\sqrt{c_0}\,\bcen n$.
Finally, \eqref{eq:imagesum} is \eqref{eq:Fstar} with $k=2$, $a=\eta_1$ and $g=\psi'_n$.
\end{proof}

\begin{remark}\label{rem:normalization}
The operator $\sum_n\bcen n^*\bcen n$ differs from $\HT$ by the factor $c_0^{-1}=\sqrt{2\pi}/\kappa^3$, which grows like
$\Nphi^{3/2}$ at fixed aspect ratio; every uniform statement below refers to $\HT$, the normalization in which a pair of
relative angular momentum one has energy one. The kernel of $\HT$ is the joint kernel of the operators $\bcen n$, which
does not depend on any normalization.
\end{remark}

% sections/03-kernel.tex
\section{The kernel}\label{sec:kernel}

The kernel of $\HT$ is the joint kernel of the pair operators $\bcen n$ (Remark~\ref{rem:normalization}). We determine
it for a family of operators that contains the $\bcen n$ and is defined for every modular parameter. The argument is
elementary: a zero mode vanishes to third order at every coincidence of two particles, so it is divisible by the cube
of the theta-function Jastrow factor; at $\Nphi=3N$ the quotient depends on the centre of mass only, and the space
of possible quotients is three-dimensional.

\subsection{Pair operators for every modular parameter}

Let $\tau\in\C$ with $\operatorname{Im}\tau>0$ and $\Nphi\in\N$. For $a\in\R$, $b\in\R$ and $t$ with $\operatorname{Im}t>0$ let
\[
\vartheta\!\begin{bmatrix}a\\b\end{bmatrix}\!(z\,|\,t):=\sum_{n\in\Z}\exp\bigl(\iu\pi t(n+a)^2+2\pi\iu(n+a)(z+b)\bigr),
\qquad \vartheta_1(z):=\vartheta\!\begin{bmatrix}1/2\\1/2\end{bmatrix}\!(z\,|\,\tau).
\]
The series converge absolutely and locally uniformly, and directly from them
\begin{equation}\label{eq:thetashift}
\vartheta\!\begin{bmatrix}a\\b\end{bmatrix}\!(z+1|t)=e^{2\pi\iu a}\vartheta\!\begin{bmatrix}a\\b\end{bmatrix}\!(z|t),\qquad
\vartheta\!\begin{bmatrix}a\\b\end{bmatrix}\!(z+t|t)=e^{-\iu\pi t-2\pi\iu(z+b)}\vartheta\!\begin{bmatrix}a\\b\end{bmatrix}\!(z|t).
\end{equation}
For $j\in\Z_{\Nphi}$ let $\varpi_j(u):=\vartheta\bigl[\begin{smallmatrix}j/\Nphi\\0\end{smallmatrix}\bigr](\Nphi u\,|\,\Nphi\tau)
=\sum_{n\equiv j}\exp(2\pi\iu nu+\iu\pi\tau n^2/\Nphi)$. Let $\Lam_u:=\Z+\tau\Z$.

Let $(\mathsf e_j)_{j\in\Z_{\Nphi}}$ be the standard basis of $\C^{\Nphi}$, $c_j$ the corresponding annihilation operators on
$\bigwedge\C^{\Nphi}$, and for $n\in\Z$
\begin{equation}\label{eq:Bntau}
B_n^{(\tau)}:=\sum_{A+B=n}(A-B)\,\mathsf q^{(A-B)^2}\,c_{[B]}c_{[A]},\qquad \mathsf q^{d^2}:=e^{\iu\pi\tau d^2/(2\Nphi)},
\end{equation}
the sum running over $(A,B)\in\Z^2$; $B^{(\tau)}_{n+2\Nphi}=B^{(\tau)}_n$. Let $T_x$, $T_y$ be the unitaries $\mathsf e_j\mapsto\mathsf e_{j+1}$
and $\mathsf e_j\mapsto e^{2\pi\iu j/\Nphi}\mathsf e_j$.

\emph{Relation to the torus.} For the rectangular torus of Definition~\ref{def:torus} put $\tau:=\iu L_x/L_y$ and
$u:=(y-\iu x)/L_y$. Then $e^{\iu\kappa ny}e^{-(x-\kappa n)^2/2}=e^{-x^2/2}e^{2\pi\iu nu+\iu\pi\tau n^2/\Nphi}$ (use
$\kappa(y-\iu x)=2\pi u$ and $\kappa^2/2=\pi L_x/(L_y\Nphi)$), so
\begin{equation}\label{eq:psiphi}
\psi_j(x,y)=(L_y\sqrt\pi)^{-1/2}\,\varpi_j(u)\,e^{-x^2/2},
\end{equation}
$\mathsf q^{d^2}=e^{-\kappa^2d^2/4}$, and under $\mathsf e_j\mapsto\psi_j$ the operators $B_n^{(\tau)}$ are the $\bcen n$ of
Proposition~\ref{prop:orbital} (pair $(A,B)$ with $(B,A)$), while $T_x$, $T_y$ are $\thx$, $\thy$. In particular the
kernel of $\HT$ on $\bigwedge^N\h$ is the joint kernel of the $B_n^{(\tau)}$.

\emph{Holomorphic parts.} For $S=\{s_1<\dots<s_N\}$ let
$\widehat{\mathsf e_{s_1}\wedge\cdots\wedge\mathsf e_{s_N}}(u_1,\dots,u_N):=\det[\varpi_{s_\alpha}(u_\beta)]_{\alpha,\beta}$, and extend linearly to a map
$\Psi\mapsto\widehat\Psi$ from $\bigwedge^N\C^{\Nphi}$ to functions on $\C^N$ (for $N=0$, $\widehat{1}:=1$). By
\eqref{eq:psiphi}, for a rectangular torus the wave function of $\Psi$ is
$(N!)^{-1/2}(L_y\sqrt\pi)^{-N/2}\prod_ke^{-x_k^2/2}\,\widehat\Psi(u)$. We write $U:=\sum_ku_k$, $e_k$ for the $k$-th unit vector of
$\C^N$, and $\one:=(1,\dots,1)$.

\begin{theorem}[the kernel at $\Nphi=3N$]\label{thm:kernel}
Let $\operatorname{Im}\tau>0$, $N\ge1$ and $\Nphi=3N$. The joint kernel $\mathcal K_N$ of the operators $B_n^{(\tau)}$,
$n\in\Z_{2\Nphi}$, on $\bigwedge^N\C^{\Nphi}$ is three-dimensional. It is spanned by the states $\Psi_0,\Psi_1,\Psi_2$ with
\begin{equation}\label{eq:laughlin}
\widehat{\Psi_j}(u)=A_j(U)\,J(u)^3,\quad J(u):=\prod_{1\le i<k\le N}\vartheta_1(u_i-u_k),\quad
A_j(s):=\vartheta\!\begin{bmatrix}j/3+(N-1)/6\\(N-1)/2\end{bmatrix}\!(3s\,|\,3\tau).
\end{equation}
Moreover $\Ff(T_y)\Psi_j=e^{2\pi\iu m_j/\Nphi}\Psi_j$ with $m_j=N(N-1)/2+jN$, and
$\Ff(T_x)\Psi_j=e^{-\iu\pi(N-1)/3}\Psi_{j+1}$ (indices modulo $3$): the three states are the three centre-of-mass copies
of one state. For a rectangular torus, $\mathcal K_N=\ker(\HT|_{\bigwedge^N\h})$.
\end{theorem}

\begin{proposition}[no zero modes above filling one third]\label{prop:overfill}
If $N'\ge2$ and $\Nphi<3N'$, the joint kernel of the $B^{(\tau)}_n$ on $\bigwedge^{N'}\C^{\Nphi}$ is zero. In particular, at
$\Nphi=3N$ no sector with more than $N$ particles contains a zero mode.
\end{proposition}

For $N\ge2$ the functions \eqref{eq:laughlin} are the torus Laughlin states of~\cite{HaldaneRezayiTorus}. That they are the
only zero modes of the torus $V_1$ Hamiltonian is stated for rectangular tori in~\cite[Section~II.2]{ONDS}, and a proof for
every lattice shape and boundary twist is claimed in~\cite{HouWu}. Theorem~\ref{thm:kernel} proves the count for the operators
$B_n^{(\tau)}$ at every modular parameter, and hence for $\HT$ on every rectangular torus.

\subsection{Theta functions}

For $\ell\in\Z$, $a\in\R$ and $\beta\in\C\setminus\{0\}$ let $\Theta_\ell[a;\beta]$ be the space of entire functions $g$ on
$\C$ with $g(s+1)=e^{2\pi\iu a}g(s)$ and $g(s+\tau)=\beta e^{-2\pi\iu\ell s}g(s)$.

\begin{lemma}\label{lem:theta}
\begin{enumerate}[(a)]
\item $\dim\Theta_\ell[a;\beta]=\ell$ for $\ell\ge1$, and $\Theta_\ell[a;\beta]=0$ for $\ell\le-1$.
\item A nonzero $g\in\Theta_\ell[a;\beta]$, $\ell\ge1$, has exactly $\ell$ zeros, counted with multiplicity, in every period
parallelogram whose boundary avoids its zeros.
\item $\vartheta_1$ is odd, $\vartheta_1(s+1)=-\vartheta_1(s)$, $\vartheta_1(s+\tau)=-e^{-\iu\pi\tau-2\pi\iu s}\vartheta_1(s)$, and its zeros
are exactly the points of $\Lam_u$, all simple.
\item $J(u_\pi)=\sgn(\pi)J(u)$ for permutations $\pi$; $J(u+c\one)=J(u)$ for $c\in\C$; $J(u+e_k)=(-1)^{N-1}J(u)$;
$J(u+\tau e_k)=(-1)^{N-1}e^{-\iu\pi(N-1)\tau-2\pi\iu(Nu_k-U)}J(u)$; and the factor $\vartheta_1(u_i-u_k)$ vanishes, simply, exactly
on the hyperplanes $D_{ik}^\lambda:=\{u_i-u_k=\lambda\}$, $\lambda\in\Lam_u$.
\item If $f$ is a function on $\C^N$ with $f(u+e_k)=f(u)$ and $f(u+\tau e_k)=e^{-\iu\pi\Nphi(2u_k+\tau)}f(u)$ for all $k$, then
$\mathcal A:=f/J^3$ satisfies, where defined, $\mathcal A(u+e_k)=(-1)^{N-1}\mathcal A(u)$ and
\[
\mathcal A(u+\tau e_k)=(-1)^{N-1}\exp\bigl(-2\pi\iu(\Nphi-3N)u_k-6\pi\iu U-\iu\pi(\Nphi-3N+3)\tau\bigr)\,\mathcal A(u).
\]
At $\Nphi=3N$ the factor does not depend on $u_k$: $\mathcal A(u+\tau e_k)=(-1)^{N-1}e^{-3\pi\iu(2U+\tau)}\mathcal A(u)$.
\item The space $\Theta^{(N)}_3:=\{A\ \text{entire}:A(s+1)=(-1)^{N-1}A(s),\ A(s+\tau)=(-1)^{N-1}e^{-3\pi\iu(2s+\tau)}A(s)\}$ is
three-dimensional, with basis $A_0,A_1,A_2$ of \eqref{eq:laughlin}, and $A_j(s+1/3)=e^{2\pi\iu a_j}A_j(s)$ with
$a_j:=j/3+(N-1)/6$.
\end{enumerate}
\end{lemma}

\begin{proof}
(a) $g(s)e^{-2\pi\iu as}$ is entire and $1$-periodic, hence a Laurent series in $e^{2\pi\iu s}$, absolutely and locally
uniformly convergent on $\C^\times$: $g(s)=\sum_nc_ne^{2\pi\iu(n+a)s}$. Comparing coefficients in the $\tau$-relation gives
$c_ne^{2\pi\iu(n+a)\tau}=\beta c_{n+\ell}$, and by iteration, for $k\ge0$,
$\abs{c_{m+k\ell}}=\abs\beta^{-k}\exp(-2\pi\operatorname{Im}\tau[k(m+a)+\ell k(k-1)/2])\abs{c_m}$. If $\ell\ge1$ this is Gaussian decay
in both directions, so any $c_0,\dots,c_{\ell-1}$ determine a convergent series in the space, and conversely; the dimension is
$\ell$. If $\ell\le-1$ and some $c_m\neq0$, then $\abs{c_n}$ grows like $\exp(\pi\operatorname{Im}\tau\abs\ell k^2)$ along
$n=m+k\ell\to-\infty$, so the terms of the series at any point do not tend to zero, a contradiction; hence $g=0$.
(b) $g'/g$ is $1$-periodic and $(g'/g)(s+\tau)=(g'/g)(s)-2\pi\iu\ell$, so the integral of $g'/g$ around the parallelogram with
vertices $s_0,s_0+1,s_0+1+\tau,s_0+\tau$ is $2\pi\iu\ell$; apply the argument principle.
(c) The relations are \eqref{eq:thetashift} with $a=b=1/2$; oddness follows by $n\mapsto-n-1$ in the series. The Fourier
coefficients do not vanish, so $\vartheta_1\not\equiv0$, and $\vartheta_1\in\Theta_1[\tfrac12;-e^{-\iu\pi\tau}]$ has exactly one zero in
every period parallelogram, by (b). Since $\vartheta_1(0)=0$ by oddness, the zeros are the points of $\Lam_u$, each simple.
Consequently $\vartheta_1(s-1)=-\vartheta_1(s)$ and $\vartheta_1(s-\tau)=-e^{-\iu\pi\tau+2\pi\iu s}\vartheta_1(s)$.
(d) A transposition of neighbouring variables changes the sign of one factor (oddness) and permutes the others. The
shift $u\mapsto u+e_k$ changes the sign of the $N-1$ factors containing $u_k$. Under $u\mapsto u+\tau e_k$ the factor with
partner $j$ is multiplied by $-e^{-\iu\pi\tau-2\pi\iu(u_k-u_j)}$ in both orders of the pair (by (c)), and
$\sum_{j\neq k}(u_k-u_j)=Nu_k-U$. The last statement is (c).
(e) Cube the relations in (d), divide, and use $-\iu\pi\Nphi\tau+3\iu\pi(N-1)\tau=-\iu\pi(\Nphi-3N+3)\tau$.
(f) By (a) with $\ell=3$, $e^{2\pi\iu a}=(-1)^{N-1}$ and $\beta=(-1)^{N-1}e^{-3\pi\iu\tau}$, the dimension is $3$. By
\eqref{eq:thetashift} with $t=3\tau$, $A_j(s+1)=e^{6\pi\iu a_j}A_j(s)=(-1)^{N-1}A_j(s)$ and
$A_j(s+\tau)=e^{-2\pi\iu(N-1)/2}e^{-3\pi\iu\tau-6\pi\iu s}A_j(s)$; the $A_j$ have nonzero Fourier coefficients on the disjoint
frequency sets $j+(N-1)/2+3\Z$, so they are independent. The last formula is the first relation of \eqref{eq:thetashift}
at $z=3s$ with $t=3\tau$, shifted by $1$.
\end{proof}

\subsection{Holomorphic parts and pair coefficients}

\begin{lemma}\label{lem:holo}
\begin{enumerate}[(a)]
\item $\Psi\mapsto\widehat\Psi$ is injective, and its image on $\bigwedge^N\C^{\Nphi}$ is the space of antisymmetric entire
functions $f$ on $\C^N$ with $f(u+e_k)=f(u)$ and $f(u+\tau e_k)=e^{-\iu\pi\Nphi(2u_k+\tau)}f(u)$ for every $k$.
\item For $N\ge2$ and $\Psi\in\bigwedge^N\C^{\Nphi}$ write $u=(u_1,u_2,u')$. Then
\[
\widehat\Psi(u_1,u_2,u')=\sum_{a,b\in\Z_{\Nphi}}\varpi_a(u_1)\varpi_b(u_2)\,\widehat{c_bc_a\Psi}(u').
\]
\item $\widehat{\Ff(T_x)\Psi}(u)=e^{2\pi\iu U+\iu\pi N\tau/\Nphi}\,\widehat\Psi(u+\tfrac{\tau}{\Nphi}\one)$ and
$\widehat{\Ff(T_y)\Psi}(u)=\widehat\Psi(u+\tfrac1{\Nphi}\one)$.
\end{enumerate}
\end{lemma}

\begin{proof}
(a) Each $\varpi_j$ lies in $\Theta_{\Nphi}[0;e^{-\iu\pi\Nphi\tau}]$ (from the series, as in Lemma~\ref{lem:theta}(f)), and the
$\varpi_j$ have disjoint Fourier supports $j+\Nphi\Z$, so they form a basis of this $\Nphi$-dimensional space. Slater
determinants therefore have the stated properties. Products $\prod_\beta\varpi_{j_\beta}(u_\beta)$ for distinct index tuples are
linearly independent (fix all variables but one and induct), and distinct sets $S$ give disjoint sets of tuples, so the
map is injective. Conversely let $f$ have the properties. For $N=1$ use the basis. For $N\ge2$, choose points
$s_1,\dots,s_{\Nphi}$ with $M:=[\varpi_j(s_i)]$ invertible (the evaluation functionals separate the span of the $\varpi_j$);
for fixed $u'=(u_2,\dots,u_N)$, $u_1\mapsto f(u_1,u')$ lies in the one-particle space, so $f(u_1,u')=\sum_jc_j(u')\varpi_j(u_1)$
with $c(u')=M^{-1}(f(s_i,u'))_i$, entire and with the one-particle relations in each remaining variable. By induction on the
number of variables (without antisymmetry) $f$ is a finite combination of products of orbitals; the coefficients are unique,
and antisymmetry of $f$ makes them antisymmetric, so $f$ is a combination of Slater determinants.
(b) By linearity take $\Psi=\mathsf e_{s_1}\wedge\cdots\wedge\mathsf e_{s_N}$ and $a=s_p$, $b=s_q$, $p\neq q$. The Laplace expansion of
$\det[\varpi_{s_\alpha}(u_\beta)]$ along the columns $1,2$ gives the coefficient of $\varpi_{s_p}(u_1)\varpi_{s_q}(u_2)$ as
$(-1)^{p+q+1}\widehat{\Psi_{pq}}$ if $p<q$ and $(-1)^{p+q}\widehat{\Psi_{pq}}$ if $p>q$, where $\Psi_{pq}$ is the wedge with
$\mathsf e_{s_p},\mathsf e_{s_q}$ removed. On the Fock side, $c_{s_p}$ acts first with the sign $(-1)^{p-1}$, then $c_{s_q}$ with the sign
$(-1)^{q-2}$ if $q>p$ and $(-1)^{q-1}$ if $q<p$; the totals agree. Terms with $a=b$ or $\{a,b\}\not\subset S$ vanish on both sides.
(c) $T_x$ and $T_y$ act on one-particle holomorphic parts by $f(u)\mapsto e^{\iu\pi(2u+\tau/\Nphi)}f(u+\tau/\Nphi)$ and
$f(u)\mapsto f(u+1/\Nphi)$, as the series of $\varpi_j$ show; apply this to each row of the determinant.
\end{proof}

\subsection{Zero modes vanish to third order}

\begin{lemma}\label{lem:formulation}
Let $N\ge2$, any $\Nphi\ge1$, $\Psi\in\bigwedge^N\C^{\Nphi}$ and $f:=\widehat\Psi$. The following are equivalent:
(i) $B^{(\tau)}_n\Psi=0$ for every $n\in\Z_{2\Nphi}$;
(ii) $[(\partial_{u_1}-\partial_{u_2})f](w,w,u')=0$ for all $w\in\C$, $u'\in\C^{N-2}$;
(iii) $f=(u_1-u_2)^3g$ with $g$ entire;
(iv) for all $i<k$ and $\lambda\in\Lam_u$, $f=(u_i-u_k-\lambda)^3g_{ik}^\lambda$ with $g_{ik}^\lambda$ entire.
\end{lemma}

\begin{proof}
(i)$\Leftrightarrow$(ii). With $\varpi_a'(w)=\sum_{A\equiv a}2\pi\iu A\,e^{2\pi\iu Aw+\iu\pi\tau A^2/\Nphi}$ and Lemma~\ref{lem:holo}(b),
\[
[(\partial_1-\partial_2)f](w,w,u')=\sum_{a,b}\bigl[\varpi_a'\varpi_b-\varpi_a\varpi_b'\bigr](w)\,\widehat{c_bc_a\Psi}(u')
=2\pi\iu\sum_{n\in\Z}e^{2\pi\iu nw+\iu\pi\tau n^2/(2\Nphi)}\,\widehat{B^{(\tau)}_n\Psi}(u'),
\]
using $A^2+B^2=\frac12(A+B)^2+\frac12(A-B)^2$ and grouping by $n=A+B$; all rearrangements are justified by absolute
convergence. At fixed $\operatorname{Im}w$ the right side is an absolutely convergent Fourier series in $\operatorname{Re}w$ in which each
frequency $n\in\Z$ occurs once, with coefficient
$e^{-2\pi n\operatorname{Im}w}e^{\iu\pi\tau n^2/(2\Nphi)}\widehat{B^{(\tau)}_{n}\Psi}(u')$. Hence (ii) holds iff
$\widehat{B^{(\tau)}_n\Psi}\equiv0$ for all $n$, iff (i), by injectivity (Lemma~\ref{lem:holo}(a)).
(ii)$\Leftrightarrow$(iii). Put $\tilde f(R,v,u'):=f(R+v/2,R-v/2,u')$, entire and odd in $v$. Then $\tilde f$ and
$\partial_v^2\tilde f$ vanish at $v=0$, and (ii) says $\partial_v\tilde f|_{v=0}=0$. So the first three Taylor coefficients in $v$
vanish and $\tilde f=v^3\tilde g$ with $\tilde g(R,v,u')=\frac12\int_0^1(1-\theta)^2\partial_v^3\tilde f(R,\theta v,u')\dd\theta$ entire.
Conversely (iii) makes $(\partial_1-\partial_2)f$ vanish on $u_1=u_2$.
(iii)$\Rightarrow$(iv). Antisymmetry transports the factor to every pair $(i,k)$. By the relations of
Lemma~\ref{lem:holo}(a), $f(u+\lambda e_i)=M_\lambda(u)f(u)$ with $M_\lambda$ entire and nowhere zero, so
$f(u)=M_\lambda(u-\lambda e_i)f(u-\lambda e_i)$ inherits the factor $(u_i-u_k-\lambda)^3$. (iv)$\Rightarrow$(iii) is the case
$(1,2,0)$.
\end{proof}

\begin{lemma}[division by $J^3$]\label{lem:division}
If $f$ is entire on $\C^N$ and satisfies Lemma~\ref{lem:formulation}(iv), then $f/J^3$ is entire.
\end{lemma}

\begin{proof}
Let $\ell,\ell'$ be nonconstant affine-linear forms on $\C^N$. (D1) If $q$ is entire and $q=(\ell-\lambda)^3\tilde q_\lambda$ with
$\tilde q_\lambda$ entire for every $\lambda\in\Lam_u$, then $q/\vartheta_1(\ell)^3$ is entire: near a point with $\ell=\lambda$,
$\vartheta_1(\ell)=(\ell-\lambda)g(\ell)$ with $g$ holomorphic and nonzero (Lemma~\ref{lem:theta}(c)), so the quotient is
$\tilde q_\lambda/g(\ell)^3$. (D2) If $q=hg$ with $h,g$ entire, $q=(\ell'-\mu)^3\tilde q$ with $\tilde q$ entire, and $h\not\equiv0$ on
$\mathsf H:=\{\ell'=\mu\}$, then $g=(\ell'-\mu)^3\tilde g$ with $\tilde g$ entire: choose $\xi$ on which the linear part of $\ell'$
equals $1$; an entire function has the factor $(\ell'-\mu)^3$ iff it and its first two derivatives along $\xi$ vanish on $\mathsf H$ (Taylor's formula with integral
remainder along $\xi$); by Leibniz, if $\partial^r_\xi g|_{\mathsf H}=0$ for $r<m\le2$ then
$0=\partial^m_\xi(hg)|_{\mathsf H}=h|_{\mathsf H}\,\partial_\xi^mg|_{\mathsf H}$, so
$\partial^m_\xi g$ vanishes on the dense open subset $\{h\neq0\}$ of the connected hyperplane $\mathsf H$, hence on $\mathsf H$.
Enumerate the pairs $p_1,\dots,p_M$, $M=N(N-1)/2$, with forms $\ell_{p}:=u_i-u_k$ for $p=(i,k)$, and put $q_0:=f$,
$q_r:=q_{r-1}/\vartheta_1(\ell_{p_r})^3$. By induction, $q_r$ is entire and has the factor $(\ell_{p_t}-\mu)^3$ for every $t>r$ and
$\mu\in\Lam_u$: $q_r$ is entire by (D1); and (D2) applies to $q_{r-1}=\vartheta_1(\ell_{p_r})^3q_r$ on $\{\ell_{p_t}=\mu\}$, on which
$\ell_{p_r}$ is nonconstant (distinct pairs give independent forms), so $\vartheta_1(\ell_{p_r})^3\not\equiv0$ there. Then
$f/J^3=q_M$ is entire.
\end{proof}

\subsection{The quotient and the count}

\begin{lemma}\label{lem:quotient}
Let $N\ge2$, $\Nphi=3N$, $\Psi\in\mathcal K_N$, $f=\widehat\Psi$. Then $f/J^3=A(U)$ for some $A\in\Theta_3^{(N)}$.
\end{lemma}

\begin{proof}
$\mathcal A:=f/J^3$ is entire by Lemmas~\ref{lem:formulation} and~\ref{lem:division}, and by Lemma~\ref{lem:theta}(e) (off the
coincidence set, then everywhere by continuity)
\[
\mathcal A(u+e_k)=(-1)^{N-1}\mathcal A(u),\qquad \mathcal A(u+\tau e_k)=(-1)^{N-1}e^{-3\pi\iu(2U+\tau)}\mathcal A(u).
\]
Fix $u$ and $i\neq k$ and put $g(w):=\mathcal A(u+w(e_i-e_k))$; $U$ is constant along this line. Then $g(w+1)=g(w)$, and,
applying the $k$-relation backwards at $v-\tau e_k$ and then the $i$-relation, $\mathcal A(v+\tau e_i-\tau e_k)=\mathcal A(v)$ for every $v$, so
$g(w+\tau)=g(w)$. Thus $g$ is entire and doubly periodic, hence constant (Liouville). Moving along one direction
$e_k-e_1$ at a time, $\mathcal A(u+c)=\mathcal A(u)$ whenever $\sum_kc_k=0$, so $\mathcal A(u)=\mathcal A(Ue_1)=:A(U)$ with $A$ entire; its relations are those of
$\mathcal A$ at $u=se_1$.
\end{proof}

\begin{lemma}\label{lem:converse}
Let $N\ge1$, $\Nphi=3N$ and $A\in\Theta_3^{(N)}$. Then $f_A(u):=A(U)J(u)^3$ is the holomorphic part of a unique
$\Psi_A\in\bigwedge^N\C^{\Nphi}$, $\Psi_A\in\mathcal K_N$, and $A\mapsto\Psi_A$ is linear and injective.
\end{lemma}

\begin{proof}
$f_A$ is entire and antisymmetric. By Lemma~\ref{lem:theta}(d),(f), $f_A(u+e_k)=f_A(u)$ and
$f_A(u+\tau e_k)=(-1)^{N-1}e^{-3\pi\iu(2U+\tau)}(-1)^{N-1}e^{-3\iu\pi(N-1)\tau-6\pi\iu(Nu_k-U)}f_A(u)=e^{-3\pi\iu N(2u_k+\tau)}f_A(u)$,
so $f_A=\widehat{\Psi_A}$ for a unique $\Psi_A$ by Lemma~\ref{lem:holo}(a). For $N\ge2$,
$f_A=(u_1-u_2)^3A(U)[\vartheta_1(u_1-u_2)/(u_1-u_2)]^3\prod_{(i,k)\neq(1,2)}\vartheta_1(u_i-u_k)^3$ with an entire cofactor, so
$\Psi_A\in\mathcal K_N$ by Lemma~\ref{lem:formulation}; for $N=1$ there is nothing to prove. If $A\not\equiv0$ then $f_A\not\equiv0$.
\end{proof}

\begin{proof}[Proof of Theorem~\ref{thm:kernel}]
For $N=1$ every one-particle state is in $\mathcal K_1$, and $\C^{3}$ is three-dimensional; the holomorphic parts are the
space $\Theta_3[0;e^{-3\iu\pi\tau}]=\Theta_3^{(1)}$, and $J=1$. For $N\ge2$, Lemmas~\ref{lem:quotient} and~\ref{lem:converse} give
mutually inverse linear maps between $\mathcal K_N$ and $\Theta_3^{(N)}$ (injectivity of $\Psi\mapsto\widehat\Psi$), so
$\dim\mathcal K_N=3$ by Lemma~\ref{lem:theta}(f), with basis $\Psi_j:=\Psi_{A_j}$. By Lemma~\ref{lem:holo}(c),
$\Ff(T_y)$ shifts $U$ by $N/\Nphi=1/3$ and leaves $J$ invariant, so $\Ff(T_y)\Psi_j=e^{2\pi\iu a_j}\Psi_j$, and
$\Nphi a_j=jN+N(N-1)/2$. $\Ff(T_x)$ multiplies by $e^{2\pi\iu U+\iu\pi\tau/3}$ and shifts $U$ by $\tau/3$; completing the square in
the series,
$\vartheta\bigl[\begin{smallmatrix}a\\b\end{smallmatrix}\bigr](z+\tau|3\tau)=e^{-\iu\pi\tau/3-2\pi\iu(z+b)/3}\vartheta\bigl[\begin{smallmatrix}a+1/3\\b\end{smallmatrix}\bigr](z|3\tau)$,
and $\vartheta[{}^{a+1}_{\;b}]=\vartheta[{}^a_b]$; with $z=3U$ this gives $\Ff(T_x)\Psi_j=e^{-2\pi\iu b/3}\Psi_{j+1}$ with
$b=(N-1)/2$. The last statement is the relation to the torus explained above.
\end{proof}

\begin{proof}[Proof of Proposition~\ref{prop:overfill}]
Let $\Psi$ be in the joint kernel. Lemmas~\ref{lem:formulation} and~\ref{lem:division} do not use $\Nphi=3N'$, so
$\mathcal A:=\widehat\Psi/J^3$ is entire with the relations of Lemma~\ref{lem:theta}(e) (with $N'$ in place of $N$). Put $m:=\Nphi-3N'<0$
and $g(w):=\mathcal A(u+w(e_1-e_2))$ for fixed $u$. As in Lemma~\ref{lem:quotient}, $g(w+1)=g(w)$, and the same two-step
computation with the general relation gives $g(w+\tau)=e^{-2\pi\iu m(u_1-u_2+2w+\tau)}g(w)$. So
$g\in\Theta_{2m}[0;e^{-2\pi\iu m(u_1-u_2+\tau)}]$ with $2m\le-2$, and $g=0$ by Lemma~\ref{lem:theta}(a). Hence $\mathcal A\equiv0$,
$\widehat\Psi\equiv0$ and $\Psi=0$. At $\Nphi=3N$ a sector with $N'>N\ge1$ particles has $N'\ge2$ and $\Nphi<3N'$.
\end{proof}

% sections/04-projector.tex
\section{The pair Gram matrix and the projector form}\label{sec:projector}

On two particles $\HT=h=\sum_n\kb{v_n}{v_n}$ is not a projection: the pair functions $v_n$ of
Proposition~\ref{prop:orbital} are not orthonormal. Their Gram matrix is computed exactly in this section; its
smallest eigenvalue enters the main theorem through the two-body term. We also record the normal-ordered square of a
lifted pair operator, which is the starting point of Section~\ref{sec:square}.

\subsection{The Gram matrix}

Let $G_{nn'}:=\ip{v_n}{v_{n'}}$ for $n,n'\in\Z_{2\Nphi}$, and put
\[
\mathfrak f(s):=(1-s)e^{-s/2},\qquad s_{jk}:=\tfrac14\bigl(j^2L_x^2+k^2L_y^2\bigr)\quad((j,k)\in\Z^2).
\]

\begin{lemma}[the Gram matrix in closed form]\label{lem:gram}
$G_{nn'}=0$ unless $n'\equiv n\pmod{\Nphi}$, and for integers $n$ and $m$
\begin{equation}\label{eq:gram}
G_{n,n+m\Nphi}=\sum_{j\equiv m\ (2)}\ \sum_{k\in\Z}e^{\iu\pi k(n+j\Nphi/2)}\,\mathfrak f(s_{jk}).
\end{equation}
\end{lemma}

\begin{proof}
By the proof of Proposition~\ref{prop:orbital}, $v_n$ is $\sqrt{2c_0}$ times the tensor $\sum_{A+B=n}f(A-B)\psi_{[A]}\otimes\psi_{[B]}$.
Hence $G_{nn'}=2c_0\sum f(A-B)f(A'-B')$ over $A+B=n$, $A'+B'=n'$, $A'\equiv A$, $B'\equiv B\pmod{\Nphi}$. Writing
$A'=A+p\Nphi$, $B'=B+q\Nphi$ forces $(p+q)\Nphi=n'-n$, which proves the first claim. For $n'=n+m\Nphi$ we have $p+q=m$ and
$A'-B'=(A-B)+j\Nphi$ with $j=p-q$, which runs once over the integers $j\equiv m\pmod2$. With $d=A-B$, which runs over
$d\equiv n\pmod 2$,
\[
G_{n,n+m\Nphi}=\sum_{j\equiv m\ (2)}S_j,\qquad S_j:=2c_0\sum_{d\equiv n\ (2)}f(d)f(d+j\Nphi).
\]
Put $M:=j\Nphi$ and $u:=d+M/2$. Then $f(d)f(d+M)=(u^2-M^2/4)e^{-\kappa^2u^2/2}e^{-\kappa^2M^2/8}$, and $u$ runs over the coset
$n+M/2+2\Z$. Poisson summation on a coset, $\sum_{u\in c+2\Z}g(u)=\frac12\sum_{k\in\Z}\hat g(k/2)e^{\iu\pi kc}$ with
$\hat g(\xi)=\int g(u)e^{-2\pi\iu u\xi}\dd u$, and
\[
\int_\R\bigl(u^2-M^2/4\bigr)e^{-\kappa^2u^2/2}e^{-\iu\pi ku}\dd u=\frac{\sqrt{2\pi}}{\kappa^3}\Bigl(1-\frac{\pi^2k^2}{\kappa^2}-\frac{\kappa^2M^2}4\Bigr)e^{-\pi^2k^2/(2\kappa^2)}
\]
give, since $\pi k/\kappa=kL_y/2$ and $\kappa M=jL_x$,
$S_j=\sum_ke^{\iu\pi k(n+j\Nphi/2)}(1-s_{jk})e^{-s_{jk}/2}$.
\end{proof}

Let $\Dlev(\ell):=2(\ell^2/4-1)e^{-\ell^2/8}$ for $\ell>0$, let $\mathcal X:=\{(0,0),(\pm1,0),(0,\pm1)\}$ and
\[
\Ups(L_x,L_y):=\sum_{(j,k)\in\Z^2\setminus\mathcal X}\abs{\mathfrak f(s_{jk})} .
\]

\begin{lemma}[the levels]\label{lem:levels}
The matrix $G$ is block diagonal with $2\times2$ blocks on $\{n,n+\Nphi\}$, and its entries depend on $n$ only through
$n\bmod2$.
\begin{enumerate}[(a)]
\item If $\Nphi$ is even, $G$ has the four eigenvalues $1-(-1)^s\Dlev(L_y)\mp\Dlev(L_x)+r_{s,\pm}$, $s\in\{0,1\}$, each of
multiplicity $\Nphi/2$, with eigenvectors $e_n\pm e_{n+\Nphi}$, $n\equiv s$.
\item If $\Nphi$ is odd, $G$ has the two eigenvalues $1\pm\sqrt{\Dlev(L_x)^2+\Dlev(L_y)^2}+r_\pm$, each of multiplicity $\Nphi$.
\item In both cases $\abs r\le\Ups(L_x,L_y)$, and for $L_x,L_y\ge3$
\[
\Ups(L_x,L_y)\le1.20\,(L_x^2+L_y^2)\,e^{-(L_x^2+L_y^2)/8}+2.02\,\bigl(L_x^2e^{-L_x^2/2}+L_y^2e^{-L_y^2/2}\bigr).
\]
\end{enumerate}
Consequently the extreme eigenvalues $\lambda_{\min},\lambda_{\max}$ of $G$ satisfy
$\dG:=\max\{\lambda_{\max}-1,1-\lambda_{\min}\}\le\Dlev(L_x)+\Dlev(L_y)+\Ups(L_x,L_y)$.
\end{lemma}

\begin{proof}
The block structure and the dependence on $n\bmod2$ are read off \eqref{eq:gram}. Put $d_s:=G_{nn}$ and
$o_s:=G_{n,n+\Nphi}$ for $n\equiv s\pmod2$. Both are real: $d_s$ is a diagonal entry of the Hermitian matrix $G$; for even
$\Nphi$ every phase in $o_s$ is $\pm1$; and for odd $\Nphi$ the terms of $o_s$ with odd $k$ cancel in the pairs $(j,k),(-j,k)$, as
shown in (b), while those with even $k$ have phases $\pm1$.
In $d_s$ the term $(0,0)$ is $1$ and the terms $(0,\pm1)$ give $2(-1)^s\mathfrak f(L_y^2/4)=-(-1)^s\Dlev(L_y)$; in $o_s$ the terms
$(\pm1,0)$ give $2\mathfrak f(L_x^2/4)=-\Dlev(L_x)$. The remaining terms of $d_s$ (even $j$) and of $o_s$ (odd $j$) run over disjoint
subsets of $\Z^2\setminus\mathcal X$, so $\abs{d_s-1+(-1)^s\Dlev(L_y)}+\abs{o_s+\Dlev(L_x)}\le\Ups$.

(a) For even $\Nphi$, $n$ and $n+\Nphi$ have the same parity $s$; the block is
$\bigl[\begin{smallmatrix}d_s&o_s\\o_s&d_s\end{smallmatrix}\bigr]$ with eigenvalues $d_s\pm o_s$ on $e_n\pm e_{n+\Nphi}$, and there are
$\Nphi/2$ blocks of each parity.

(b) For odd $\Nphi$, $n$ and $n+\Nphi$ have opposite parities. Moreover $o_s=:o$ does not depend on $s$: for odd $j$ and odd
$\Nphi$, the terms with odd $k$ cancel in the pairs $(j,k),(-j,k)$ because $\cos(\pi kj\Nphi/2)=0$, and for even $k$ the phase
$e^{\iu\pi kn}$ equals $1$. The block is $\bigl[\begin{smallmatrix}d_s&o\\o&d_{1-s}\end{smallmatrix}\bigr]$, with eigenvalues
$\frac12(d_0+d_1)\pm\sqrt{\frac14(d_0-d_1)^2+o^2}$, the same for all $\Nphi$ blocks. With the leading terms above and
$\abs{\sqrt{a^2+b^2}-\sqrt{a'^2+b'^2}}\le\abs{a-a'}+\abs{b-b'}$ the error is at most $\max_s\abs{d_s-1+(-1)^s\Dlev(L_y)}+\abs{o+\Dlev(L_x)}\le\Ups$.

(c) Put $a:=L_x^2/4$ and $b:=L_y^2/4$, both at least $9/4$. Off $\mathcal X$, $s_{jk}\ge9/2$, so $\abs{\mathfrak f(s)}\le se^{-s/2}$. The
terms with $j=0$, $\abs k\ge2$ sum to at most $2\cdot4be^{-2b}/(1-q)$, where $q\le\frac94e^{-5b/2}\le0.0081$ bounds the ratio of
consecutive terms; this is at most $2.02\,L_y^2e^{-L_y^2/2}$, and likewise for $k=0$, $\abs j\ge2$. The terms with $j,k\neq0$ sum
to $\Sigma_1(a)\Sigma_0(b)+\Sigma_0(a)\Sigma_1(b)$ with $\Sigma_0(a)=\sum_{j\ne0}e^{-j^2a/2}\le2e^{-a/2}/(1-e^{-3a/2})$ and
$\Sigma_1(a)=\sum_{j\neq0}j^2ae^{-j^2a/2}\le2ae^{-a/2}/(1-4e^{-3a/2})$ (use $j^2-1\ge3(j-1)$ and $j^2\le4^{j-1}$); with $e^{-3a/2}\le e^{-27/8}$
this is at most $4(a+b)e^{-(a+b)/2}/((1-e^{-27/8})(1-4e^{-27/8}))\le1.20(L_x^2+L_y^2)e^{-(L_x^2+L_y^2)/8}$.
The bound on $\dG$ follows from (a)--(c) and $\sqrt{a^2+b^2}\le a+b$.
\end{proof}

\subsection{The projector form}

\begin{proposition}[the projector form]\label{prop:projector}
$\spec(h)\subset\{0\}\cup[\lambda_{\min},\lambda_{\max}]$, so $h^2\ge\lambda_{\min}h$. If $\dG<1$, then $\dim\VT=2\Nphi$, and
with $P_V$ the orthogonal projection onto $\VT$ and $H':=\LL_2(P_V)$,
\[
\lambda_{\min}P_V\le h\le\lambda_{\max}P_V,\qquad \lambda_{\min}H'\le\HT\le\lambda_{\max}H',\qquad \ker H'=\ker\HT .
\]
\end{proposition}

\begin{proof}
Let $\mathsf V:\C^{2\Nphi}\to\bigwedge^2\h$ have the columns $v_n$. Then $h=\mathsf V\mathsf V^*$, whose nonzero spectrum is that of
$G=\mathsf V^*\mathsf V$; this gives the first claim, and $h^2-\lambda_{\min}h\ge0$ by the functional calculus. If $\dG<1$ then
$G>0$, $\mathsf V$ is injective, $\dim\VT=2\Nphi$, and with the isometry $W:=\mathsf VG^{-1/2}$ one has $h=WGW^*$ and $P_V=WW^*$,
whence the first sandwich. The lift
preserves order (Lemma~\ref{lem:lift}(iii)), which gives the second, and the kernels coincide because $\lambda_{\min}>0$.
\end{proof}

\subsection{The square of a lifted pair operator}

For an operator $h$ on $\bigwedge^2\h$, extended by $hA_2$ to $\h^{\otimes2}$, and $\{i,j\}\subset\{1,\dots,n\}$, $h_{ij}$ denotes
$h$ acting on the tensor positions $i,j$; $h_{ij}=h_{ji}$.

\begin{lemma}[the normal-ordered square]\label{lem:square}
For every self-adjoint $h$ on $\bigwedge^2\h$, on $\Fock$,
\begin{equation}\label{eq:square}
\LL_2(h)^2=\LL_2(h^2)+\LL_3\bigl(X_3(h)\bigr)+\LL_4\bigl(W_4(h\otimes h)W_4^*\bigr),\qquad
X_3(h):=\sum_{(ij)\neq(kl)}h_{ij}h_{kl}\ \text{ on }\textstyle\bigwedge^3\h,
\end{equation}
the sum running over the six ordered pairs of distinct pairs in $\{1,2,3\}$. Moreover
$X_3(h)=W_3(h\otimes\one)(W_3^*W_3-\one)(h\otimes\one)W_3^*$. For $h=P_V$ a projection this is the identity
$H'^2=H'+\LL_3\bigl(W_3(W_3^*W_3-\one)W_3^*\bigr)+\LL_4(W_4W_4^*)$, with $W_3$, $W_4$ restricted to $\VT\otimes\h$ and
$\VT\otimes\VT$, of~\cite[Lemma~2.1 (Normal-ordered square)]{OAI:A-Fock-space-inequality-and-the-Laughlin-spectral-gap-September-24-2026}.
\end{lemma}

\begin{proof}
On $\bigwedge^n\h$, $\LL_2(h)=\sum_{\abs S=2}h_S$ by Lemma~\ref{lem:lift}(ii), so $\LL_2(h)^2=\sum_{S,S'}h_Sh_{S'}$. The terms with
$S=S'$ give $\sum_S(h^2)_S=\LL_2(h^2)$. The terms with $\abs{S\cap S'}=1$ are grouped by the triple $Q=S\cup S'$ and give
$\sum_{\abs Q=3}(X_3)_Q$, where $X_3$ commutes with permutations of its three positions; this is $\LL_3(X_3)$. The terms with
$S\cap S'=\emptyset$ are grouped by $Q=S\cup S'$, $\abs Q=4$, and give $\LL_4(X_4)$ with $X_4:=\sum h_Sh_{S'}$ over the six ordered
pairs of disjoint pairs in $\{1,2,3,4\}$. On antisymmetric tensors $W_4^*\Psi=\sqrt6\Psi$ and $W_4=\sqrt6A_4$
(Lemma~\ref{lem:lift}(iv)), and $A_4h_{12}h_{34}\Psi=\frac16X_4\Psi$, because each ordered pair of disjoint pairs is the image of
$(\{1,2\},\{3,4\})$ under four permutations; hence $X_4=W_4(h\otimes h)W_4^*$ on $\bigwedge^4\h$. Similarly
$W_3(h\otimes\one)W_3^*=3A_3h_{12}=h_{12}+h_{13}+h_{23}$ and $W_3(h^2\otimes\one)W_3^*=(h^2)_{12}+(h^2)_{13}+(h^2)_{23}$ on
$\bigwedge^3\h$, so $X_3=(W_3(h\otimes\one)W_3^*)^2-W_3(h^2\otimes\one)W_3^*$, which is the second formula. For $h=P_V$,
$h^2=h$ and the restrictions are as stated.
\end{proof}

% sections/05-local-modes.tex
\section{Local modes on the torus}\label{sec:local}

The comparison operator of Section~\ref{sec:moyal} is built, at every base point $R$, from the periodizations of the
planar orbitals and pair states about $R$. This section defines them and bounds their Gram matrix.

\subsection{Periodized orbitals and pair states}

For $p\ge0$ and $R\in\R^2$ let $v_p^{*R}$ be the planar two-particle state with relative orbital $\eta_1$ and centre-of-mass
orbital $\eta_p$ about $\sqrt2R$, that is $v_p^{*R}(r_1,r_2):=\eta_1(\rho)\,\eta_p^{\sqrt2R}(\zeta)$ with $\rho=(r_1-r_2)/\sqrt2$,
$\zeta=(r_1+r_2)/\sqrt2$; it is antisymmetric. Let
\[
v_p^*(i,j):=(i-j)\sqrt{\frac{p!}{2^p\,i!\,j!}}\;\one_{i+j=p+1}\qquad(0\le i<j).
\]

\begin{lemma}[periodized local modes]\label{lem:localmodes}
Let $\chi_m^R:=\Pk\eta_m^R$ and $\ut_p(R):=\Pk^{(2)}v_p^{*R}$.
\begin{enumerate}[(a)]
\item $\chi_m^R\in\h$, $\ut_p(R)\in\bigwedge^2\h$, and
$\ip{\psi_j}{\chi_m^R}=\sqrt\kappa\sum_{n\equiv j}e^{-\iu\kappa nR_y}\mathcal H_m(\kappa n-R_x)$.
\item $v_p^{*R}=\sum_{i<j}v_p^*(i,j)\,\eta_i^R\wedge\eta_j^R$, and hence $\ut_p(R)=\sum_{i<j}v_p^*(i,j)\,\chi_i^R\wedge\chi_j^R$.
\item For $\gamma\in\Lam$ there is a real $\theta=\theta(R,\gamma)$, independent of $m$, with
$\chi_m^{R+\gamma}=e^{\iu\theta}\chi_m^R$; then $\ut_p(R+\gamma)=e^{2\iu\theta}\ut_p(R)$.
\end{enumerate}
\end{lemma}

\begin{proof}
(a) $\eta_m^R$ is an LLL function with Gaussian decay about $R$, so Lemma~\ref{lem:unfold} applies, and
$\ip{\psi_j}{\chi_m^R}=\overline{\ip{\chi_m^R}{\psi_j}}=\int_{\R^2}\overline{\psi_j}\,\eta_m^R=\sum_{n\equiv j}\ip{\phi_n}{\eta_m^R}_{\R^2}$, which
is the stated sum by \eqref{eq:landaurep}. $\Pk^{(2)}$ preserves antisymmetry. (b) At $R=0$, in the Bargmann form
(Lemmas~\ref{lem:bargmann}(i) and~\ref{lem:jacobi}(i)), $v_p^{*0}$ corresponds to $wZ^p/\sqrt{p!}$ with $w=(z_1-z_2)/\sqrt2$,
$Z=(z_1+z_2)/\sqrt2$. In $(z_1-z_2)(z_1+z_2)^p$ the coefficient of $z_1^iz_2^j$, $i+j=p+1$, is
$\binom p{i-1}-\binom pi=p!(i-j)/(i!j!)$; passing to the normalized tensors $z_1^iz_2^j/\sqrt{i!j!}$ and then to the normalized
wedges ($\eta_i\wedge\eta_j$ corresponds to $(\eta_i\otimes\eta_j-\eta_j\otimes\eta_i)/\sqrt2$) gives the coefficient $v_p^*(i,j)$. For
general $R$ apply $t_R\otimes t_R$, which acts as $t_{\sqrt2R}$ on $\zeta$ and trivially on $\rho$ (Lemma~\ref{lem:jacobi}(ii)) and
multiplies both sides by $e^{\iu R_xR_y}$ (by \eqref{eq:etaR}). The second identity follows because $\Pk^{(2)}$ is linear and
$\Pk^{(2)}(\eta_i\wedge\eta_j)=\Pk\eta_i\wedge\Pk\eta_j$.
(c) By \eqref{eq:etaR} and Lemma~\ref{lem:magtrans}(ii), $\eta_m^{R+\gamma}$ is an $m$-independent phase times $t_\gamma\eta_m^R$,
and $\Pk t_\gamma=\pm\Pk\tl_\gamma=\pm\Pk$ by reindexing. The statement for $\ut_p$ follows from (b).
\end{proof}

\subsection{The local Gram matrix}

For $R\in\R^2$ let $\Gamma(R)_{mm'}:=\ip{\chi_m^R}{\chi_{m'}^R}$, $m,m'\ge0$, and let $\Gamma^{(m)}(R)$ be its truncation to
$0\le m,m'\le m$. For $r\ge0$ let
\[
\mu_m(r):=\norm{\Dis^{(m)}(r/\sqrt2)},\qquad \Ec_m(\Tor):=\sum_{\gamma\in\Lam\setminus\{0\}}\mu_m(\abs\gamma);
\]
by Lemma~\ref{lem:bargmann}(iii), $\norm{\Dis^{(m)}(\beta)}=\mu_m(\sqrt2\abs\beta)$ for every $\beta\in\C$.

\begin{lemma}[the local Gram matrix]\label{lem:localgram}
For every $R\in\R^2$ and $m,m'\ge0$,
\begin{equation}\label{eq:localgram}
\Gamma(R)_{mm'}=\sum_{\gamma\in\Lam}\varsigma_\gamma\,e^{\iu\,\gamma\wedge R}\,\bra m\Dis(\beta(\gamma))\ket{m'},\qquad
\varsigma_{(jL_x,kL_y)}:=(-1)^{jk\Nphi},
\end{equation}
with $\beta(\gamma)=(\gamma_x-\iu\gamma_y)/\sqrt2$. Consequently
\[
\norm{\Gamma^{(m)}(R)-\one}\le\Ec_m(\Tor),\qquad \Gamma(R)_{mm}\le1+\sum_{\gamma\in\Lam\setminus0}\abs{\bra m\Dis(\beta(\gamma))\ket m}.
\]
\end{lemma}

\begin{proof}
By Lemma~\ref{lem:unfold}, $\Gamma(R)_{mm'}=\int_{\R^2}\overline{\eta_m^R}\,\Pk\eta_{m'}^R=\sum_\gamma\ip{\eta_m^R}{\tl_\gamma\eta_{m'}^R}_{\R^2}$, absolutely
convergent. By \eqref{eq:etaR} the phases $e^{\mp\iu R_xR_y/2}$ cancel, and
$t_R^*t_\gamma t_R=t_{-R}t_\gamma t_R=e^{\iu\,\gamma\wedge R}t_\gamma$ by Lemma~\ref{lem:magtrans}(ii). Lemma~\ref{lem:bargmann}(ii)
gives $\ip{\eta_m}{t_\gamma\eta_{m'}}=\bra m\Dis(\beta(\gamma))\ket{m'}$. The term $\gamma=0$ is $\delta_{mm'}$, and the
truncation of each other term has norm $\mu_m(\abs\gamma)$.
\end{proof}

\subsection{Bounds on the lattice sums}

\begin{lemma}\label{lem:Ebounds}
For $m\ge0$ and $r\ge0$ put $\mathfrak g_m(r):=(r^2/2+m)^me^{-r^2/4}/m!$.
\begin{enumerate}[(a)]
\item For $m\le23$, $\mu_m$ is nonincreasing on $[16,\infty)$. Consequently, for every rectangular torus whose shorter period
is at least $L$, with $L\ge16$, and every $m\le23$,
\[
\Ec_m(\Tor)\le\Ec_m(L):=\sum_{(j,k)\in\Z^2\setminus0}\mu_m\bigl(L\sqrt{j^2+k^2}\bigr),
\]
and $\Ec_m(L)$ is nonincreasing in $L\ge16$.
\item For $m\ge1$ and $r\ge0$, $\mu_m(r)\le(m+1)\,\mathfrak g_m(r)$, and $\mathfrak g_m$ is decreasing
for $r^2>2m$. Hence, for every rectangular torus whose shorter period is at least $L$ with $L^2\ge2m$,
\[
\Ec_m(\Tor)\le\widehat{\Ec}_m(L):=4(m+1)\,\mathfrak g_m(L)\Bigl(1+\sum_{n\ge2}(n+1)\,n^m\,e^{-(n-1)L^2/4}\Bigr)\xrightarrow[L\to\infty]{}0 ,
\]
and the same bound holds for $\Ec_m(L)$.
\item For every $m\ge0$ and $\abs\beta=r/\sqrt2$, $\abs{\bra m\Dis(\beta)\ket m}\le\min\{1,\mathfrak g_m(r)\}$;
$\mathfrak g_m(r)$ is nondecreasing in $m$, and $\min\{1,\mathfrak g_z(r)\}$ is
nonincreasing in $r$. Hence, for every rectangular torus whose shorter period is at least $L>0$,
\[
\sigma_z(\Tor):=\sup_R\max_{m\le z}\bigl(\Gamma(R)_{mm}-1\bigr)\le\Xi_z(L):=\sum_{(j,k)\in\Z^2\setminus0}\min\bigl\{1,\mathfrak g_z\bigl(L\sqrt{j^2+k^2}\bigr)\bigr\}.
\]
\end{enumerate}
\end{lemma}

\begin{proof}
(a) Take $\beta=r/\sqrt2>0$ and put $x:=\beta^2=r^2/2$, so that $r\ge16$ means $x\ge128$. By Lemma~\ref{lem:bargmann}(iii), for
$a,b\le m$ the entry $\bra a\Dis(\beta)\ket b$ equals $\pm\sqrt{n!/(n+\alpha)!}\,x^{\alpha/2}e^{-x/2}\mathrm L_n^{(\alpha)}(x)$ with
$n=\min(a,b)$, $\alpha=\abs{a-b}$, and the sign $(-1)^{b-a}$ if $a<b$. The monic polynomials $p_i:=(-1)^ii!\,\mathrm L_i^{(\alpha)}$
satisfy $p_{i+1}=(x-(2i+\alpha+1))p_i-i(i+\alpha)p_{i-1}$, so, expanding a determinant along its last row, $p_n(x)=\det(x-M_n)$ for
the symmetric tridiagonal matrix $M_n$ with diagonal entries $2i+\alpha+1$ ($0\le i<n$) and off-diagonal entries $\sqrt{i(i+\alpha)}$
($1\le i<n$). Hence the zeros $x_1,\dots,x_n$ of $\mathrm L_n^{(\alpha)}$ are real; they are positive (for $x\le0$ all terms of
$\mathrm L_n^{(\alpha)}(x)$ have the same sign), and by Gershgorin's theorem and $\sqrt{i(i+\alpha)}\le i+\alpha/2$ they are at most
$\max_i(4i+2\alpha+2)\le4n+2\alpha-2=2a+2b-2\le90$. So for $x\ge128$ the polynomial $\mathrm L_n^{(\alpha)}(x)$ has the sign
$(-1)^n$ of its leading coefficient, and every entry of the column $b$ of $\Dis^{(m)}(\beta)$ has the sign $(-1)^b$; thus
$\Dis^{(m)}(\beta)=\abs{\Dis^{(m)}(\beta)}\,\mathrm{diag}((-1)^b)$ with the entrywise modulus $\abs{\Dis^{(m)}(\beta)}$, and
$\mu_m(r)=\norm{\abs{\Dis^{(m)}(\beta)}}$. The logarithmic derivative of the modulus of an entry is
$\frac\alpha{2x}-\frac12+\sum_i\frac1{x-x_i}$. By convexity of $t\mapsto1/(x-t)$ on $[0,90]$, $\frac1{x-t}\le\frac1x+\frac t{x(x-90)}$,
and $\sum_ix_i=\tr M_n=n(n+\alpha)$; hence for $x\ge128$
\[
\frac\alpha{2x}+\sum_i\frac1{x-x_i}\le\frac{\alpha+2n}{2x}+\frac{n(n+\alpha)}{x(x-90)}\le\frac{46}{256}+\frac{529}{128\cdot38}<0.29<\frac12 .
\]
So every entry of $\abs{\Dis^{(m)}(\beta)}$ is nonincreasing in $r\ge16$. For entrywise nonnegative matrices the spectral norm is
$\max\{v^{\mathsf T}Au:u,v\ge0,\ \norm u=\norm v=1\}$, which is monotone in the entries. For a torus with periods $L_x,L_y\ge L\ge16$,
$\abs{(jL_x,kL_y)}\ge L\sqrt{j^2+k^2}\ge16$ for $(j,k)\neq0$, which gives the bound and the monotonicity in $L$.

(b) For $a\ge b$ write $n=b$, $\alpha=a-b$. The coefficients of $\mathrm L_n^{(\alpha)}(-x)=\sum_i\binom{n+\alpha}{n-i}x^i/i!$ are
dominated by those of $(x+n+\alpha)^n/n!$, since $(n+\alpha)!/(\alpha+i)!\le(n+\alpha)^{n-i}$; and $\abs{\mathrm L_n^{(\alpha)}(x)}\le\mathrm L_n^{(\alpha)}(-x)$
for $x\ge0$ termwise. Hence the entry is at most $\sqrt{b!/a!}\,x^{(a-b)/2}e^{-x/2}(x+a)^b/b!\le(x+m)^{(a+b)/2}e^{-x/2}/\sqrt{a!b!}$,
which is nondecreasing in $a$ and in $b$ as long as $x+m\ge a+1$, so at most $(x+m)^me^{-x/2}/m!=\mathfrak g_m(r)$. The spectral norm of
an $(m+1)\times(m+1)$ matrix is at most $m+1$ times its largest entry. $\frac{\dd}{\dd r}\log\mathfrak g_m(r)=r\bigl(\frac m{r^2/2+m}-\frac12\bigr)<0$ for
$r^2>2m$; so $\mu_m(\abs{(jL_x,kL_y)})\le(m+1)\,\mathfrak g_m(L\sqrt{j^2+k^2})$ when $L^2\ge2m$. Group the lattice points by
$n=j^2+k^2$: there are $4$ with $n=1$ and at most $2(2\sqrt n+1)\le4(n+1)$ for each $n\ge2$, and
$\mathfrak g_m(L\sqrt n)\le n^me^{-(n-1)L^2/4}\mathfrak g_m(L)$.

(c) $\abs{\bra m\Dis(\beta)\ket m}\le1$ by unitarity, and $=e^{-x/2}\abs{\mathrm L_m^{(0)}(x)}\le e^{-x/2}\mathrm L_m^{(0)}(-x)\le\mathfrak g_m(r)$ by
(b) with $\alpha=0$. Next $\mathfrak g_{m+1}/\mathfrak g_m\ge(x+m+1)/(m+1)\ge1$. Finally $\frac{\dd}{\dd r}\log\mathfrak g_z>0$ only for
$r^2<2z$, and there $\mathfrak g_z(r)\ge\mathfrak g_z(0)=z^z/z!\ge1$; so $\min\{1,\mathfrak g_z\}$ equals $1$ on $[0,\sqrt{2z}]$ and is
decreasing beyond. With Lemma~\ref{lem:localgram} and $\abs{(jL_x,kL_y)}\ge L\sqrt{j^2+k^2}$ this gives the bound on $\sigma_z$.
\end{proof}

\begin{remark}\label{rem:closedbound}
The bound (b) is far from sharp near the thresholds of the main theorem. For $m=22$ and $m=23$, the indices of the binding
four-body term, it gives $\widehat\Ec_{22}(L)\approx9.82$, $0.0255$ and $3.6\cdot10^{-5}$ and $\widehat\Ec_{23}(L)\approx77.9$,
$0.221$ and $3.4\cdot10^{-4}$ at $L=16$, $17$ and $18$, against the values $\Ec_{22}(17)=1.13\cdot10^{-6}$ and
$\Ec_{23}(17)=4.73\cdot10^{-6}$ computed from the exact norms (Table~\ref{tab:constants}). It is used only for the limit
$L\to\infty$ and its rate; the numerical values in Section~\ref{sec:theorem} use the exact norms $\mu_m$ and part (a).
\end{remark}

% sections/06-moyal.tex
\section{The torus Moyal identities and the comparison operator}\label{sec:moyal}

On the sphere the comparison operator is averaged over rotations, and Schur's lemma for $SU(2)$ reduces it to fixed
spin blocks. On the torus we average the base point $R$ of the local modes over the torus. The magnetic translations
of the centre-of-mass torus then play the role of the rotations, and the resulting identities are exact.

\subsection{Magnetic translations of the centre-of-mass torus}

\begin{lemma}\label{lem:irred}
Let $\Tor'$ be a rectangular torus with periods $L'_x,L'_y$, lattice $\Lam'$, $M=L'_xL'_y/2\pi\in\N$ flux quanta and lowest
Landau level $\U$, and let $\Pk'$ be its periodization. For $\delta\in\Lam'/M$, $t_\delta$ is unitary on $\U$ and
$t_\delta\Pk'f=\Pk't_\delta f$ for every Gaussian-decaying LLL function $f$. An operator on $\U$ that commutes with all $t_\delta$,
$\delta\in\Lam'/M$, is a multiple of the identity.
\end{lemma}

\begin{proof}
The first statement is Lemma~\ref{lem:magtrans}(iv) for $\Tor'$, and $t_\delta\tl'_{\gamma'}=e^{\iu\,\delta\wedge\gamma'}\tl'_{\gamma'}t_\delta$
with $\delta\wedge\gamma'\in2\pi\Z$. Let $C$ commute with $\theta'_y:=t_{(0,-L'_y/M)}$ and $\theta'_x:=t_{(L'_x/M,0)}$. By
Lemma~\ref{lem:magtrans}(iv), $\theta'_y$ has the $M$ distinct eigenvalues $e^{2\pi\iu j/M}$ on the orbitals $\psi'_j$, so
$C\psi'_j=c_j\psi'_j$, and $\theta'_x\psi'_j=\psi'_{j+1}$ forces $c_{j+1}=c_j$.
\end{proof}

\subsection{The Moyal identity}

For $k\ge1$ a \emph{relative state} is a function $a$ of the relative Jacobi coordinates $\rho\in(\R^2)^{k-1}$ which is an LLL
function of each of them and decays like a Gaussian. For $R\in\R^2$ and $p\ge0$ we write $a\otimes\varphi_p^R$ for the planar
$k$-particle function $(\zeta,\rho)\mapsto a(\rho)\,\eta_p^{\sqrt kR}(\zeta)$: relative state $a$, centre-of-mass orbital of
degree $p$ about the point $R$ (the centre of mass $k^{-1}\sum r_i$ equals $\zeta/\sqrt k$). By Lemma~\ref{lem:jacobi}(i) it is an
LLL function of each particle and decays like a Gaussian about $R$ in all variables, so $\Pk^{(k)}(a\otimes\varphi_p^R)\in\h^{\otimes k}$.

\begin{theorem}[torus Moyal identity]\label{thm:moyal}
Let $k\ge1$, let $a,b$ be relative states and $p,q\ge0$. The operator $\kb{\Pk^{(k)}(a\otimes\varphi_p^R)}{\Pk^{(k)}(b\otimes\varphi_q^R)}$ on
$\h^{\otimes k}$ is a $\Lam$-periodic continuous function of $R$, and
\begin{equation}\label{eq:moyal}
\frac k{2\pi}\int_\cell\kb{\Pk^{(k)}(a\otimes\varphi_p^R)}{\Pk^{(k)}(b\otimes\varphi_q^R)}\,\dd R=\delta_{pq}\,\Pp_k\bigl(\kb ab\otimes\one\bigr)=\delta_{pq}\,\Fib_a^*\Fib_b .
\end{equation}
\end{theorem}

\begin{proof}
Let $\Pk'$ be the periodization of the centre-of-mass torus $\Tor_k=\R^2/\sqrt k\Lam$. For $\Psi\in\h^{\otimes k}$, by
Lemma~\ref{lem:unfold} on $\Tor$, Fubini in Jacobi coordinates, and Lemma~\ref{lem:unfold} on $\Tor_k$ (recall
$\Fib_b\Psi\in\U_k$),
\[
\ip{\Pk^{(k)}(b\otimes\varphi_q^R)}\Psi=\int_{(\R^2)^k}\overline{b(\rho)\,\eta_q^{\sqrt kR}(\zeta)}\,\Psi=\int_{\R^2}\overline{\eta_q^{\sqrt kR}}\,\Fib_b\Psi
=\ip{\Pk'\eta_q^{\sqrt kR}}{\Fib_b\Psi}_{\U_k}.
\]
Hence, as a form in $(\Phi,\Psi)$, the left side of \eqref{eq:moyal} is $\ip{\Fib_a\Phi}{\mathcal Q_{pq}\Fib_b\Psi}$ with
\[
\mathcal Q_{pq}:=\frac1{2\pi}\int_{\sqrt k\cell}\kb{\Pk'\eta_p^{\zeta_0}}{\Pk'\eta_q^{\zeta_0}}\,\dd\zeta_0
\]
(substitute $\zeta_0=\sqrt kR$). \emph{Periodicity.} For $\gamma'\in\sqrt k\Lam$, $\eta_p^{\zeta_0+\gamma'}$ is a $p$-independent phase times
$t_{\gamma'}\eta_p^{\zeta_0}$ (by \eqref{eq:etaR} and Lemma~\ref{lem:magtrans}(ii)), and $\Pk't_{\gamma'}=\pm\Pk'$; so the integrand is
periodic, which also proves the first assertion. \emph{Schur.} For $\delta\in\sqrt k\Lam/(k\Nphi)$, $t_\delta\eta_p^{\zeta_0}$ is a
$p$-independent phase times $\eta_p^{\zeta_0+\delta}$ and $t_\delta$ commutes with $\Pk'$ (Lemma~\ref{lem:irred}); so
$t_\delta\mathcal Q_{pq}t_\delta^*$ is the integral of the translated periodic integrand over a period cell, which is
$\mathcal Q_{pq}$. By Lemma~\ref{lem:irred}, $\mathcal Q_{pq}=c_{pq}\one$ on $\U_k$. \emph{The scalar.} $c_{pq}=\tr\mathcal Q_{pq}/(k\Nphi)$ and
$\tr\kb{\Pk'\eta_p^{\zeta_0}}{\Pk'\eta_q^{\zeta_0}}=\ip{\Pk'\eta_q^{\zeta_0}}{\Pk'\eta_p^{\zeta_0}}
=\sum_{\gamma'\in\sqrt k\Lam}\varsigma'_{\gamma'}e^{\iu\,\gamma'\wedge\zeta_0}\bra q\Dis(\beta(\gamma'))\ket p$, exactly as in
Lemma~\ref{lem:localgram}, the series converging absolutely. For $\gamma'=(j\sqrt kL_x,l\sqrt kL_y)\neq0$ the integral of
$e^{\iu\,\gamma'\wedge\zeta_0}=e^{\iu j\sqrt kL_x\zeta_{0,y}}e^{-\iu l\sqrt kL_y\zeta_{0,x}}$ over $\sqrt k\cell$ vanishes, because
$\sqrt kL_x\cdot\sqrt kL_y=2\pi k\Nphi$; the term $\gamma'=0$ gives $\abs{\sqrt k\cell}\,\delta_{pq}=2\pi k\Nphi\,\delta_{pq}$. Hence
$c_{pq}=\delta_{pq}$, and Lemma~\ref{lem:Pk}(c) gives the second equality.
\end{proof}

The identity holds for every $k$ with the normalization built into the definition of $\Pp_k$. For $k=2$ the normalization
of $h$ itself was identified in Proposition~\ref{prop:orbital}, so no further normalization step is needed.

\begin{corollary}[pair level]\label{cor:pairmoyal}
For $p,q\ge0$, $\frac1\pi\int_\cell\kb{\ut_p(R)}{\ut_q(R)}\dd R=\delta_{pq}\,h$. Consequently $\ut_p(R)\in\VT$ for every $p$ and $R$, and
\[
\frac1\pi\int_\cell B(\ut_p(R))^*B(\ut_p(R))\,\dd R=\HT\qquad\text{for every }p\ge0 .
\]
\end{corollary}

\begin{proof}
Take $k=2$ and $a=b=\eta_1$; then $\Pk^{(2)}(\eta_1\otimes\varphi_p^R)=\ut_p(R)$ and $\Pp_2(\kb{\eta_1}{\eta_1}\otimes\one)=h$. If $w\perp\VT$,
then $\int_\cell\abs{\ip w{\ut_p(R)}}^2\dd R=\pi\ip w{hw}=0$, and $R\mapsto\ip w{\ut_p(R)}$ is continuous, so it vanishes. The last
identity follows by applying the linear map $\LL_2$ (Lemma~\ref{lem:lift}) to the first, the integrand being continuous with
values in a finite-dimensional space.
\end{proof}

\subsection{Planar coupled bases}

The certificate of~\cite{OAI:A-Fock-space-inequality-and-the-Laughlin-spectral-gap-September-24-2026} is written in terms of
coupling coefficients. For $0<u<1$, $v:=\sqrt{1-u^2}$ and integers $0\le z\le T$, $0\le p\le T$, the coefficients
$s^{(u)}_{z,T,p}$ are defined by the polynomial identity
\begin{equation}\label{eq:coupling}
\frac{(uX-vY)^z(vX+uY)^{T-z}}{\sqrt{z!(T-z)!}}=\sum_{p=0}^Ts^{(u)}_{z,T,p}\,\frac{X^pY^{T-p}}{\sqrt{p!(T-p)!}}
\end{equation}
\cite[Lemma~3.3 (Fixed-deficit coupling limit)]{OAI:A-Fock-space-inequality-and-the-Laughlin-spectral-gap-September-24-2026}, and are zero for indices outside these
ranges. The substitution $(X,Y)\mapsto(uX-vY,vX+uY)$ is orthogonal, hence unitary for the Bargmann inner product, so for
fixed $T$ the real matrix $(s^{(u)}_{z,T,p})_{z,p}$ is orthogonal.

\paragraph{Three particles.} Use the Jacobi coordinates $\rho_1=(r_1-r_2)/\sqrt2$, $\rho_2=(r_1+r_2-2r_3)/\sqrt6$ and
$\zeta=(r_1+r_2+r_3)/\sqrt3$, and put $s_{z,T,p}:=s^{(1/\sqrt3)}_{z,T,p}$,
\[
\xi_z(\rho):=\eta_1(\rho_1)\,\eta_z(\rho_2),\qquad \ket{z,T}^R:=\xi_z\otimes\varphi^R_{T-z}\quad(0\le z\le T).
\]
Permutations of the particles act on $\rho$ only (Lemma~\ref{lem:jacobi}(iii)); we write $W_3\xi_z:=\sqrt3A_3\xi_z$, so that
$W_3\ket{z,T}^R=(W_3\xi_z)\otimes\varphi^R_{T-z}$.

\begin{lemma}\label{lem:coupled3}
For $p,j\ge0$ and $R\in\R^2$, $v_p^{*R}\otimes\eta_j^R=\sum_{z=0}^{p+j}s_{z,p+j,p}\ket{z,p+j}^R$, where $v_p^{*R}$ acts on the
particles $1,2$ and $\eta_j^R$ on particle $3$.
\end{lemma}

\begin{proof}
At $R=0$, in the Bargmann form, $v_p^{*0}\otimes\eta_j$ is $w\,X^pY^j/\sqrt{p!j!}$ with $w=(z_1-z_2)/\sqrt2$, $X=(z_1+z_2)/\sqrt2$,
$Y=z_3$. For $u=1/\sqrt3$, $uX-vY=(z_1+z_2-2z_3)/\sqrt6$ and $vX+uY=(z_1+z_2+z_3)/\sqrt3$ are the Bargmann forms of $\rho_2$ and
$\zeta$, so $\ket{z,T}^0$ is $w$ times the left side of \eqref{eq:coupling}; invert the orthogonal matrix. For general $R$ apply
$t_R^{\otimes3}$, which acts as $t_{\sqrt3R}$ on $\zeta$ and trivially on $\rho$ and multiplies both sides by $e^{3\iu R_xR_y/2}$.
\end{proof}

\paragraph{Four particles.} Use $\rho_1=(r_1-r_2)/\sqrt2$, $\rho_2=(r_3-r_4)/\sqrt2$, $\rho_3=(r_1+r_2-r_3-r_4)/2$ and
$\zeta=(r_1+r_2+r_3+r_4)/2$. For odd $r$ and $D\ge r$ let
\[
\omega_{D,r}(\rho):=\eta_1(\rho_1)\,\eta_r(\rho_2)\,\eta_{D-r}(\rho_3),\qquad \ket{D,T;r}^R:=\omega_{D,r}\otimes\varphi^R_{T-D}\quad(T\ge D),
\]
and, following~\cite[Section~4.2, (23)]{OAI:A-Fock-space-inequality-and-the-Laughlin-spectral-gap-September-24-2026},
$S_r(D,T;p,j,k):=\sqrt2\,s^{(1/\sqrt2)}_{r,j+k,j}\,s^{(1/\sqrt2)}_{D-r,T-r,p}$, interpreted as zero if $r>j+k$, for $p+j+k=T$. Let
$\mathcal R_D$ be the set of odd integers in $[1,D]$.

\begin{lemma}\label{lem:coupled4}
For $p\ge0$ and $j\ne k$, with $T=p+j+k$, $S_r(D,T;p,j,k)=-S_r(D,T;p,k,j)$ and
$v_p^{*R}\otimes(\eta_j^R\wedge\eta_k^R)=\sum_{D\le T}\sum_{r\in\mathcal R_D}S_r(D,T;p,j,k)\ket{D,T;r}^R$, with $v_p^{*R}$ on the
particles $1,2$ and $\eta_j^R\wedge\eta_k^R$ on $3,4$. The vectors $\ket{D,T;r}^R$, $r\in\mathcal R_D$, $D\le T$, are orthonormal;
in particular, for $b=(p,j,k)$ with $p+j+k=D$ and $j<k$, the matrix $S_{rb}:=S_r(D,D;b)$ satisfies $SS^{\mathsf T}=\one$.
\end{lemma}

\begin{proof}
At $R=0$, in the Bargmann form, $\eta_j\wedge\eta_k$ on the particles $3,4$ is $(z_3^jz_4^k-z_3^kz_4^j)/\sqrt{2j!k!}$. By
\eqref{eq:coupling} with $u=v=1/\sqrt2$, $X=z_3$, $Y=z_4$ and $m=j+k$, the normalized monomial $z_3^jz_4^{m-j}/\sqrt{j!(m-j)!}$ is
$\sum_rs^{(1/\sqrt2)}_{r,m,j}\,\omega_{34}^r\Omega_{34}^{m-r}/\sqrt{r!(m-r)!}$ with $\omega_{34}=(z_3-z_4)/\sqrt2$, $\Omega_{34}=(z_3+z_4)/\sqrt2$;
exchanging $X$ and $Y$ multiplies the left side of \eqref{eq:coupling} by $(-1)^r$, so $s^{(1/\sqrt2)}_{r,m,m-j}=(-1)^rs^{(1/\sqrt2)}_{r,m,j}$,
and the antisymmetric combination keeps the odd $r$ with the coefficient $\sqrt2\,s^{(1/\sqrt2)}_{r,m,j}$. Next couple the pair
centres $X'=(z_1+z_2)/\sqrt2$ (degree $p$) and $Y'=\Omega_{34}$ (degree $m-r$) by \eqref{eq:coupling} with $u=1/\sqrt2$: the
relative variable $(X'-Y')/\sqrt2$ is the Bargmann form of $\rho_3$ and $(X'+Y')/\sqrt2$ that of $\zeta$, and the relative degree
$d=D-r$ comes with the coefficient $s^{(1/\sqrt2)}_{D-r,T-r,p}$. With the factor $w=(z_1-z_2)/\sqrt2$ of $v_p^{*0}$ this is the
expansion. The vectors $\ket{D,T;r}^0$ are products of normalized monomials in the orthonormal Jacobi coordinates, hence
orthonormal; translate by $t_R^{\otimes4}$ for general $R$. The antisymmetry in $(j,k)$ was shown above. Finally, at fixed $T$
both $\{v_p^{*R}\otimes(\eta_j^R\wedge\eta_k^R):p+j+k=T,\ j<k\}$ and $\{\ket{D,T;r}^R:r\in\mathcal R_D,\ D\le T\}$ are orthonormal bases of the
same space (in the Bargmann form: $w$ times the polynomials of degree $T$ in the remaining variables that are odd under
$z_3\leftrightarrow z_4$), so the change of basis is orthogonal; its rows with $D=T$ are the rows of $S$, whence $SS^{\mathsf T}=\one$.
\end{proof}

We write $W_4\omega_{D,r}:=\sqrt6A_4\,\omega_{D,r}$; then $W_4\ket{D,T;r}^R=(W_4\omega_{D,r})\otimes\varphi^R_{T-D}$, and
$W_4\ket{D,D;r}^0$ is the vector $w^*_{Dr}$ of~\cite[Section~4.2, (27)]{OAI:A-Fock-space-inequality-and-the-Laughlin-spectral-gap-September-24-2026}.

\subsection{The certificate}

The following data and inequalities are taken from~\cite[Section~4 and Appendix~A]{OAI:A-Fock-space-inequality-and-the-Laughlin-spectral-gap-September-24-2026};
the exact rows are reproduced in Appendix~\ref{app:certificate}. There are seven \emph{rows} $\nu=1,\dots,7$, each with an
integer $t_\nu\in\{0,\dots,7\}$, a real $\lambda_\nu$ and real coefficients $\alpha^\nu_{pj}$ supported in
\[
\mathcal S_\nu:=\{(p,j)\in\Z^2:0\le p\le7,\ 0\le j\le8,\ 0\le p+j-t_\nu\le8\}.
\]
Put $\etas:=\sum_\nu\lambda_\nu^2=93527408868499/10^{14}$ and $\eps:=3\cdot10^{-6}$. For $z\ge0$ define
\begin{equation}\label{eq:ez}
e_z:=\sum_\nu\sum_{T=\max(z,t_\nu)}^{15}\Bigl[2\lambda_\nu\,s_{z,T,t_\nu}\sum_{p+j=T}\alpha^\nu_{pj}s_{z,T,p}
+\Bigl(\sum_{p+j=T}\alpha^\nu_{pj}s_{z,T,p}\Bigr)^2\Bigr],
\end{equation}
so that $e_z=0$ for $z\ge16$, and for $1\le D\le23$ and $r,s\in\mathcal R_D$
\begin{equation}\label{eq:ED}
(E_D)_{rs}:=-\sum_\nu\sum_{\substack{(p,j),(q,l)\in\mathcal S_\nu\\ T\ge D}}\alpha^\nu_{pj}\alpha^\nu_{ql}\,S_r(D,T;p,j,k)\,S_s(D,T;q,l,i),
\end{equation}
where $i=p+j-t_\nu$, $k=q+l-t_\nu$ and $T=p+j+k$. Let $(e_r)_{r\in\mathcal R_D}$ be the standard basis of $\C^{\mathcal R_D}$, so
that $\kb{e_1}{e_1}$ is the matrix with $(1,1)$ entry $1$ and all other entries zero, and let $G_D$ be the Gram matrix of the
vectors $W_4\omega_{D,r}$, $r\in\mathcal R_D$. The certificate consists of the inequalities
\begin{align}
&e_z<\tfrac32(3z-1)(-\tfrac12)^z\qquad\text{for } z\in\{2,4,5,\dots,15\},\label{eq:C3}\\
&G_D\,(2\kb{e_1}{e_1}-E_D+\eps\one)\,G_D\ge0\qquad\text{for }1\le D\le23,\label{eq:C4}
\end{align}
which are proved in~\cite[Appendix~A]{OAI:A-Fock-space-inequality-and-the-Laughlin-spectral-gap-September-24-2026} by exact rational
arithmetic. We use them as data.

\begin{lemma}[compression; {\cite[Lemma~5.1 (Compression and annihilator bounds)]{OAI:A-Fock-space-inequality-and-the-Laughlin-spectral-gap-September-24-2026}}]\label{lem:compression}
Let $W$ be a linear map between finite-dimensional Hilbert spaces, $G=W^*W$, and $C$ self-adjoint on the domain of $W$. Then
$GCG\ge0$ if and only if $WCW^*\ge0$. In particular \eqref{eq:C4} is equivalent to
$\sum_{r,s\in\mathcal R_D}(2\kb{e_1}{e_1}-E_D+\eps\one)_{rs}\kb{W_4\omega_{D,r}}{W_4\omega_{D,s}}\ge0$.
\end{lemma}

\begin{proof}
Both conditions say that the quadratic form of $C$ is nonnegative on $\ran G=\ran W^*$. Apply this to the map
$e_r\mapsto W_4\omega_{D,r}$.
\end{proof}

\subsection{The comparison operator and its parts}

\begin{definition}[the comparison operator]\label{def:KT}
For $R\in\R^2$ and each row $\nu$ let
\[
F_\nu(R):=\lambda_\nu\,B(\ut_{t_\nu}(R))+\sum_{(p,j)\in\mathcal S_\nu}\alpha^\nu_{pj}\,c(\chi^R_{p+j-t_\nu})^*\,c(\chi_j^R)\,B(\ut_p(R)),
\]
\[
\KT:=\frac1\pi\int_\cell\sum_\nu F_\nu(R)^*F_\nu(R)\,\dd R .
\]
By Lemma~\ref{lem:localmodes}(c), $F_\nu(R+\gamma)=e^{-2\iu\theta}F_\nu(R)$ for $\gamma\in\Lam$, so the integrand is a function on
$\Tor$. Clearly $\KT\ge0$; by Corollary~\ref{cor:pairmoyal}, every summand of $F_\nu(R)$ annihilates $\ker\HT$.
\end{definition}

These are the rows of~\cite[Section~4, (14)]{OAI:A-Fock-space-inequality-and-the-Laughlin-spectral-gap-September-24-2026} with every planar
mode replaced by its periodization; each summand removes two particles and lowers the total orbital degree by $t_\nu+1$.

\begin{proposition}[the parts of $\KT$]\label{prop:parts}
Let $Y_z:=\kb{W_3\xi_z}{W_3\xi_z}\otimes\one$ and $\Cth:=72\sum_\nu\norm{\alpha^\nu}^2$, where $\norm{\alpha^\nu}^2=\sum_{p,j}(\alpha^\nu_{pj})^2$.
Then $\KT=\KT^{(2)}+\KT^{(3)}+\KT^{(4)}$ with
\begin{enumerate}[(a)]
\item $\KT^{(2)}=\etas\,\HT$;
\item $\KT^{(3)}=\frac23\LL_3\Pp_3\bigl(\sum_{z\le15}e_zY_z\bigr)+\mathcal C_3$ with $\pm\mathcal C_3\le\Cth\,\Ec_8(\Tor)\bigl(1+\Ec_8(\Tor)\bigr)\HT$;
\item $\KT^{(4)}=\frac12\LL_4\Pp_4\bigl(\sum_{D=1}^{23}\sum_{r,s\in\mathcal R_D}(E_D)_{rs}\kb{W_4\omega_{D,r}}{W_4\omega_{D,s}}\otimes\one\bigr)$;
\item with $\mathcal T_4:=\LL_4\Pp_4\bigl(\sum_{D=1}^{23}\kb{W_4\omega_{D,1}}{W_4\omega_{D,1}}\otimes\one\bigr)$,
\[
\tfrac12\LL_4\Pp_4\Bigl(\sum_{D=1}^{23}\sum_{r\in\mathcal R_D}\kb{W_4\omega_{D,r}}{W_4\omega_{D,r}}\otimes\one\Bigr)\le1222\,(1+\Ec_{23})^2\HT,\qquad
\mathcal T_4\le2444\,(1+\Ec_{23})^2\HT,
\]
with $\Ec_{23}=\Ec_{23}(\Tor)$, and $\KT^{(4)}\le\mathcal T_4+1222\,\eps\,(1+\Ec_{23})^2\HT$.
\end{enumerate}
\end{proposition}

\begin{proof}
Fix $R$ and a row; write $t=t_\nu$, $\alpha_b=\alpha^\nu_{pj}$ for $b=(p,j)$, $i=i(b):=p+j-t$, $c_m:=c(\chi^R_m)$, $B_p:=B(\ut_p(R))$ and
$\Gamma=\Gamma(R)$. Since $\{c(f),c(g)^*\}=\ip fg$, $c_ic_k^*=\Gamma_{ik}-c_k^*c_i$, and normal ordering gives
\[
F^*F=\lambda^2B_t^*B_t+\sum_b\lambda\alpha_b\bigl[(c_iB_t)^*c_jB_p+\mathrm{h.c.}\bigr]+\sum_{b,b'}\alpha_b\alpha_{b'}\Gamma_{ik}(c_jB_p)^*c_lB_q
-\sum_{b,b'}\alpha_b\alpha_{b'}O_{p,j,k}^*O_{q,l,i},
\]
with $b'=(q,l)$, $k=i(b')$ and $O_{p,j,k}:=c_kc_jB_p$. By \eqref{eq:Bwedge}, $c_jB_p=B(\ut_p\wedge\chi_j)$ and $O_{p,j,k}=B(\ut_p\wedge\chi_j\wedge\chi_k)$,
and by Lemma~\ref{lem:localmodes}(b), $\ut_p(R)\wedge\chi^R_j=\Pk^{(3)}W_3(v_p^{*R}\otimes\eta_j^R)$ and
$\ut_p\wedge\chi_j\wedge\chi_k=\Pk^{(4)}W_4(v_p^{*R}\otimes(\eta^R_j\wedge\eta^R_k))$. Let $\KT^{(2)}$, $\KT^{(3)}$, $\KT^{(4)}$ be $\frac1\pi\int_\cell$ of the
sums over the rows of the two-, three- and four-body terms, and $\mathcal C_3$ the part of $\KT^{(3)}$ with $\Gamma_{ik}-\delta_{ik}$ in place
of $\Gamma_{ik}$.

(a) is Corollary~\ref{cor:pairmoyal} with $p=t$.

(b) For a finite-rank planar operator $X=\sum_\mu\kb{x_\mu}{y_\mu}$ with Gaussian-decaying $x_\mu,y_\mu$ write
$\Pk^{(3)}[X]:=\sum_\mu\kb{\Pk^{(3)}x_\mu}{\Pk^{(3)}y_\mu}$. With $\delta_{ik}$ in place of $\Gamma_{ik}$, the three-body terms summed over
the rows are $\LL_3\bigl(\Pk^{(3)}[W_3M_3(R)W_3^*]\bigr)$, where
\[
M_3(R):=\sum_\nu\Bigl[\lambda_\nu\sum_b\alpha_b\bigl(\kb{v_t^{*R}\otimes\eta_i^R}{v_p^{*R}\otimes\eta_j^R}+\mathrm{h.c.}\bigr)
+\sum_{b,b':\,i=k}\alpha_b\alpha_{b'}\kb{v_p^{*R}\otimes\eta_j^R}{v_q^{*R}\otimes\eta_l^R}\Bigr];
\]
we expand $M_3(R)$ with Lemma~\ref{lem:coupled3}. Every term has equal total degree on both sides ($t+i=p+j$, and $p+j=q+l$ when $i=k$),
so $M_3(R)=\sum m_{(z,c),(z',c')}\kb{\xi_z\otimes\varphi_c^R}{\xi_{z'}\otimes\varphi_{c'}^R}$ with real coefficients independent of $R$ and
supported on $z+c=z'+c'\le15$. By Theorem~\ref{thm:moyal} with $k=3$ and $\frac1\pi=\frac23\cdot\frac3{2\pi}$, only $c=c'$, hence $z=z'$,
survives, and the result is $\frac23\LL_3\Pp_3(\sum_z\bigl[\sum_cm_{(z,c),(z,c)}\bigr]Y_z)$. Finally
$\sum_cm_{(z,c),(z,c)}=\sum_{T\ge z}\ip{\ket{z,T}^0}{M_3(0)\ket{z,T}^0}$, and since $\ip{\ket{z,T}^0}{v_p^{*0}\otimes\eta_j}=s_{z,T,p}$ for
$p+j=T$ (Lemma~\ref{lem:coupled3}), this is \eqref{eq:ez}. For $\mathcal C_3$, let $x_b:=c_jB_p\Psi$ and $M_{bb'}:=\alpha_b\alpha_{b'}(\Gamma-\one)_{ik}$.
Only the modes $0,\dots,8$ occur in $(i,k)$, so $\norm M\le\norm M_{\mathrm{HS}}\le\Ec_8\norm{\alpha^\nu}^2$ by Lemma~\ref{lem:localgram}, and
since $[\ip{x_b}{x_{b'}}]\ge0$, $\abs{\sum M_{bb'}\ip{x_b}{x_{b'}}}\le\norm M\sum_b\norm{x_b}^2$. Next
$\norm{c_jB_p\Psi}^2\le\Gamma_{jj}\norm{B_p\Psi}^2\le(1+\Ec_8)\norm{B_p\Psi}^2$, each $p\le7$ occurs with at most nine values of $j$, and
$\frac1\pi\int_\cell\sum_{p\le7}\norm{B(\ut_p(R))\Psi}^2\dd R=8\ip\Psi{\HT\Psi}$ by Corollary~\ref{cor:pairmoyal}. Summing over the rows gives
$\abs{\ip\Psi{\mathcal C_3\Psi}}\le72\,\Ec_8(1+\Ec_8)\sum_\nu\norm{\alpha^\nu}^2\ip\Psi{\HT\Psi}$.

(c) The four-body terms summed over the rows are $\LL_4\bigl(\Pk^{(4)}[W_4M_4(R)W_4^*]\bigr)$, with the notation of (b) and
\[
M_4(R):=-\sum_\nu\sum_{b,b'}\alpha_b\alpha_{b'}\kb{v_p^{*R}\otimes(\eta_j^R\wedge\eta_k^R)}{v_q^{*R}\otimes(\eta_l^R\wedge\eta_i^R)}.
\]
Expanding
with Lemma~\ref{lem:coupled4}, both sides have total degree $T=p+j+k=q+l+i\le23$; Theorem~\ref{thm:moyal} with $k=4$ and
$\frac1\pi=\frac12\cdot\frac4{2\pi}$ keeps the equal centre-of-mass degrees, hence equal $D$, with the copy matrix
$\sum_{T\ge D}\ip{\ket{D,T;r}^0}{M_4(0)\ket{D,T;s}^0}$, which is \eqref{eq:ED} because
$\ip{\ket{D,T;r}^0}{v_p^{*0}\otimes(\eta_j\wedge\eta_k)}=S_r(D,T;p,j,k)$.

(d) By Theorem~\ref{thm:moyal} with $k=4$ and $p=q=0$,
\begin{gather*}
\LL_4\Pp_4\bigl(\kb{W_4\omega_{D,r}}{W_4\omega_{D,r}}\otimes\one\bigr)=\frac2\pi\int_\cell B(x_{D,r}(R))^*B(x_{D,r}(R))\,\dd R,\\
x_{D,r}(R):=\Pk^{(4)}W_4\ket{D,D;r}^R=\sum_bS_{rb}\,\ut_p(R)\wedge\chi_j^R\wedge\chi_k^R,
\end{gather*}
with $b=(p,j,k)$, $p+j+k=D$, $j<k$ (Lemma~\ref{lem:coupled4}; $S$ is real). Hence $B(x_{D,r})=\sum_bS_{rb}O_b$ with
$O_b:=c_kc_jB_p$, and since $SS^{\mathsf T}=\one$,
\[
\sum_r\norm{B(x_{D,r})\Psi}^2\le\sum_b\norm{O_b\Psi}^2\le(1+\Ec_{23})^2\sum_b\norm{B_p\Psi}^2,
\]
the mode indices $j<k\le D$ being at most $23$.
For the single vector $r=1$, Cauchy--Schwarz with $\sum_bS_{1b}^2=1$ gives the same bound for $\norm{B(x_{D,1})\Psi}^2$. There are
$n_D=\sum_{m=1}^D\lceil m/2\rceil=\lfloor(D+1)^2/4\rfloor$ triples $b$, and $\sum_{D=1}^{23}n_D=1222$. Averaging with
Corollary~\ref{cor:pairmoyal} gives the two bounds. Finally, by Lemma~\ref{lem:compression} and \eqref{eq:C4}, for each $D$,
\[
\sum_{r,s}(E_D)_{rs}\kb{W_4\omega_{D,r}}{W_4\omega_{D,s}}\le2\kb{W_4\omega_{D,1}}{W_4\omega_{D,1}}+\eps\sum_r\kb{W_4\omega_{D,r}}{W_4\omega_{D,r}}
\]
as relative operators, and $\frac12\LL_4\Pp_4(\,\cdot\otimes\one)$ preserves order (Lemmas~\ref{lem:Pk}(c) and~\ref{lem:lift}(iii)), which gives the last inequality.
\end{proof}

% sections/07-square.tex
\section{The square of \texorpdfstring{$\HT$}{H\_T}}\label{sec:square}

By Lemma~\ref{lem:square}, $\HT^2=\LL_2(h^2)+\LL_3(X_3(h))+\LL_4(W_4(h\otimes h)W_4^*)$. In this section we show that the
three-body coefficient $X_3(h)$ is \emph{exactly} the periodized image of the planar three-body coefficient. The reason is
structural: $X_3(h)$ is a sum of products of two pair operators sharing exactly one particle, and periodization is
multiplicative on such products. Consequently the three-body target transfers to the torus without error, and besides the
non-orthonormality term of Proposition~\ref{prop:parts}(b) only the planar tail has to be controlled.

\subsection{The image-sum calculus}

For a set $I$ of particles, $\Pp_I(Y)$ denotes the periodized image of an operator $Y$ admissible in the variables of $I$, and
$\Pp_I(Y)\otimes\one$ its action on $\h^{\otimes n}$, $I\subset\{1,\dots,n\}$; on functions, $Y$ acts in the variables of $I$ with the
other variables as parameters.

\begin{lemma}[image-sum representation]\label{lem:imagesum}
Let $Y$ be admissible on $I\subset\{1,\dots,n\}$, $s\in I$ and $\Psi\in\h^{\otimes n}$. Then
\[
\bigl((\Pp_I(Y)\otimes\one)\Psi\bigr)(r)=\sum_{\gamma\in\Lam^{I\setminus\{s\}}}\Bigl(\prod_{i\in I\setminus\{s\}}\tl^{(i)}_{\gamma_i}\Bigr)(Y\Psi)(r),
\]
the series converging absolutely and locally uniformly.
\end{lemma}

\begin{proof}
Expanding $\Psi$ in a product basis of $\h^{\otimes I}\otimes\h^{\otimes I^c}$ reduces the claim to $n=\abs I=k$; relabel so that $s=1$.
By Lemma~\ref{lem:Pk}(a), $Y\Psi$ is an LLL function of each variable decaying like a Gaussian in $r_i-r_1$, so the series $S$
converges absolutely and locally uniformly, and $S$ is an LLL function of each variable. Reindexing shows
$\tl^{(i)}_\gamma S=S$ for $i\ge2$. For $i=1$ write $\tl^{(1)}_\gamma=\tl^{\otimes k}_\gamma\prod_{i\ge2}(\tl^{(i)}_\gamma)^{-1}$ with $\tl^{\otimes k}_\gamma=\prod_i\tl^{(i)}_\gamma$; all
these operators commute, $\tl^{\otimes k}_\gamma$ commutes with $Y$ by (Y2) and fixes $\Psi$, so $\tl^{(1)}_\gamma S=S$ by reindexing. Thus
$S\in\h^{\otimes k}$, and for $\Phi\in\h^{\otimes k}$, unfolding the variables $r_2,\dots,r_k$ as in Lemma~\ref{lem:unfold} gives
$\ip\Phi S=\int_{\{r_1\in\cell\}}\overline\Phi\,Y\Psi=\ip\Phi{\Pp_k(Y)\Psi}$.
\end{proof}

\begin{lemma}[products sharing one particle]\label{lem:product}
Let $I_1=\{s\}\cup I'_1$ and $I_2=\{s\}\cup I'_2$ with $I'_1\cap I'_2=\emptyset$, let $Y_1$ be admissible on $I_1$ and $Y_2$ on $I_2$, and
$I:=I_1\cup I_2$. Then $Y_1Y_2$ is admissible on $I$, and on $\h^{\otimes n}$
\[
\bigl(\Pp_{I_1}(Y_1)\otimes\one\bigr)\bigl(\Pp_{I_2}(Y_2)\otimes\one\bigr)=\Pp_I(Y_1Y_2)\otimes\one .
\]
\end{lemma}

\begin{proof}
On LLL functions the identity in a variable is the integral operator with the kernel $\Ppl$ (Lemma~\ref{lem:bargmann}(iv)), so the
kernel of $Y_1Y_2$ is $\int Y_1(r_s,r_{I'_1};u,r'_{I'_1})\,Y_2(u,r_{I'_2};r'_s,r'_{I'_2})\,\dd u$. It satisfies (Y1) and (Y2), and the Gaussian
bounds (Y3) of the two factors, centred at $r_s$ and at $u$, combine under the integral over $u$ into a Gaussian bound in all
$r_i-r_s$ and $r'_i-r_s$, $i\in I$. By Lemma~\ref{lem:imagesum} with the anchor $s$,
$(\Pp_{I_2}(Y_2)\otimes\one)\Psi=\sum_{\gamma\in\Lam^{I'_2}}\tl_\gamma(Y_2\Psi)$, and applying Lemma~\ref{lem:imagesum} again,
\[
(\Pp_{I_1}(Y_1)\otimes\one)(\Pp_{I_2}(Y_2)\otimes\one)\Psi=\sum_{\gamma_1\in\Lam^{I'_1}}\tl_{\gamma_1}Y_1\sum_{\gamma_2\in\Lam^{I'_2}}\tl_{\gamma_2}Y_2\Psi
=\sum_{\gamma_1,\gamma_2}\tl_{\gamma_1}\tl_{\gamma_2}\,Y_1Y_2\Psi,
\]
because $Y_1$ acts on the variables of $I_1$, which are disjoint from $I'_2$, and therefore commutes with $\tl_{\gamma_2}$; the
interchange of $Y_1$ with the sum is justified by dominated convergence, the partial sums being bounded uniformly. The right
side is $(\Pp_I(Y_1Y_2)\otimes\one)\Psi$ by Lemma~\ref{lem:imagesum} with the anchor $s$ and the images of $I'_1\cup I'_2=I\setminus\{s\}$.
\end{proof}

The lemma fails for products sharing two particles: then the image sums of the second factor do not commute with the first,
and cross-image terms appear. This is why $h^2\neq h$: by Lemma~\ref{lem:gram}, $h$ is not a projection.

\subsection{The planar three-body Gram operator}

Let $Q_{12}:=(P_1)_{12}$ on planar three-particle functions; its range consists of the functions $\xi_z(\rho)\,g(\zeta)$, $z\ge0$,
and their closed span (Section~\ref{sec:moyal}). On this range $W_3=\sqrt3A_3$.

\begin{lemma}[the planar coefficients]\label{lem:qz}
On $\ran Q_{12}$, $W_3^*W_3-\one=\sum_{z\ge0}q_z\,\kb{\xi_z}{\xi_z}\otimes\one$ with
\[
q_z=(3z-1)\bigl(-\tfrac12\bigr)^z .
\]
In particular $\norm{W_3\xi_z}^2=1+q_z$, so $W_3\xi_z=0$ for $z\in\{0,1,3\}$ and $W_3\xi_z\neq0$ otherwise.
\end{lemma}

\begin{proof}
$W_3^*W_3=3Q_{12}A_3Q_{12}$. Since $P_{12}Q_{12}=Q_{12}P_{12}=-Q_{12}$, grouping the permutations in pairs $\sigma$, $\sigma P_{12}$ gives
$3Q_{12}A_3Q_{12}=Q_{12}(\one-P_{13}-P_{23})Q_{12}$, and $P_{12}P_{13}P_{12}=P_{23}$ gives $Q_{12}P_{13}Q_{12}=Q_{12}P_{23}Q_{12}$; so
$W_3^*W_3-\one=-2\,Q_{12}P_{23}Q_{12}$. In the Bargmann form with $w_1=(z_1-z_2)/\sqrt2$ and $w_2=(z_1+z_2-2z_3)/\sqrt6$ (the forms of $\rho_1$,
$\rho_2$), the exchange $z_2\leftrightarrow z_3$ acts by the substitution
\[
w_1\mapsto\tfrac12w_1+\tfrac{\sqrt3}2w_2,\qquad w_2\mapsto\tfrac{\sqrt3}2w_1-\tfrac12w_2,
\]
and fixes $(z_1+z_2+z_3)/\sqrt3$. It preserves the degree, and $\ran Q_{12}$ contains one relative state of each degree $1+z$, so
$Q_{12}P_{23}Q_{12}=\sum_zb_z\kb{\xi_z}{\xi_z}\otimes\one$, where $b_z$ is the coefficient of $w_1w_2^z$ in
$(\frac12w_1+\frac{\sqrt3}2w_2)(\frac{\sqrt3}2w_1-\frac12w_2)^z$, the normalizations $\sqrt{z!}$ being equal on both sides. The two
contributions are $\frac12(-\frac12)^z$ and $\frac{\sqrt3}2\cdot z\frac{\sqrt3}2(-\frac12)^{z-1}$, so $b_z=-\frac12(3z-1)(-\frac12)^z$ and $q_z=-2b_z$.
\end{proof}

\subsection{The three-body coefficient}

Recall $Y_z=\kb{W_3\xi_z}{W_3\xi_z}\otimes\one$, $s_{z,T,p}=s^{(1/\sqrt3)}_{z,T,p}$, and $\abs{q_z}\le3z2^{-z}$ for $z\ge1$.

\begin{theorem}[the square of $\HT$]\label{thm:square}
On $\Fock$, for every rectangular torus,
\begin{equation}\label{eq:squareHT}
\HT^2=\LL_2(h^2)+\sum_{z\ge0}q_z\,\LL_3\Pp_3(Y_z)+\LL_4\bigl(W_4(h\otimes h)W_4^*\bigr),
\end{equation}
the series converging absolutely in norm. Moreover, for every $z\ge0$,
\begin{equation}\label{eq:yz}
\LL_3\Pp_3(Y_z)=\frac3{2\pi}\int_\cell B(y_z(R))^*B(y_z(R))\,\dd R,\qquad y_z(R):=\sum_{p=0}^zs_{z,z,p}\,\ut_p(R)\wedge\chi_{z-p}^R .
\end{equation}
\end{theorem}

\begin{proof}
By Lemma~\ref{lem:square} it suffices to show $X_3(h)=\sum_zq_z\Pp_3(Y_z)$ on $\bigwedge^3\h$. The operator $P_1$ vanishes on functions that
are even in the relative coordinate, so $h=\Pp_2(P_1)$ vanishes on symmetric tensors and $h_{ij}=\Pp_{\{i,j\}}(P_1)\otimes\one$. Each term
$h_{ij}h_{kl}$ of $X_3(h)$ has $\abs{\{i,j\}\cap\{k,l\}}=1$, so by Lemma~\ref{lem:product}
\[
X_3(h)=\Pp_3(Y_{\mathrm{pl}}),\qquad Y_{\mathrm{pl}}:=\sum_{(ij)\neq(kl)}(P_1)_{ij}(P_1)_{kl},
\]
and by Lemma~\ref{lem:Pk}(b), on $\bigwedge^3\h$ we may replace $Y_{\mathrm{pl}}$ by $A_3Y_{\mathrm{pl}}A_3$. On planar antisymmetric
three-particle functions the computation in the proof of Lemma~\ref{lem:square}, with $P_1^2=P_1$, gives
$A_3Y_{\mathrm{pl}}A_3=W_3(W_3^*W_3-\one)W_3^*$ with $W_3$ restricted to $\ran Q_{12}$, and Lemma~\ref{lem:qz} turns this into
$\sum_zq_zY_z$, because $W_3$ acts on the relative variables only. The kernels of the $Y_z$ are
$(W_3\xi_z)(\rho)\overline{(W_3\xi_z)(\rho')}\Ppl(\zeta,\zeta')$. Since $\abs{\eta_z(x)}=\abs x^ze^{-\abs x^2/4}/\sqrt{2\pi2^zz!}$ and permutations
preserve $\abs\rho$, $\abs{(W_3\xi_z)(\rho)}\le C\abs\rho^{z+1}e^{-\abs\rho^2/4}/\sqrt{2^zz!}$; with $\abs{q_z}\le3z2^{-z}$ for $z\ge1$
(and $Y_0=0$) and
$\sum_zz\,x^z/z!=xe^x$, the series $\sum_z\abs{q_z}\abs{Y_z(r;r')}$ is bounded by
$C(1+\abs\rho\abs{\rho'})^2e^{-(\abs\rho^2+\abs{\rho'}^2)/4+\abs\rho\abs{\rho'}/4}\abs{\Ppl(\zeta,\zeta')}\le C'e^{-(\abs\rho^2+\abs{\rho'}^2)/9}\abs{\Ppl(\zeta,\zeta')}$. So the kernel series
converges locally uniformly to the (continuous) kernel of $A_3Y_{\mathrm{pl}}A_3$, under an admissible Gaussian bound, and dominated
convergence in \eqref{eq:Pk} gives $\Pp_3(\sum_zq_zY_z)=\sum_zq_z\Pp_3(Y_z)$, absolutely convergent on the finite-dimensional space
$\bigwedge^3\h$. For \eqref{eq:yz}, apply Theorem~\ref{thm:moyal} with $k=3$, $a=b=W_3\xi_z$ and $p=q=0$:
$\Pk^{(3)}(W_3\xi_z\otimes\varphi_0^R)=\Pk^{(3)}W_3\ket{z,z}^R=\sum_ps_{z,z,p}\,\ut_p(R)\wedge\chi^R_{z-p}$ by \eqref{eq:coupling} and
Lemma~\ref{lem:localmodes}(b), and then lift with $\LL_3$.
\end{proof}

In words: the three-body coefficient of $\HT^2$ is the periodized planar three-body target on \emph{every} block $z$, with
the planar coefficients $q_z$ and no correction depending on the torus.

\subsection{The three-body comparison and the tail}

\begin{proposition}[three-body comparison]\label{prop:three}
On $\Fock$, for every rectangular torus,
\[
\LL_3\bigl(X_3(h)\bigr)-\KT^{(3)}\ge-\bigl(\dtail(\Tor)+\Cth\,\Ec_8(\Tor)(1+\Ec_8(\Tor))\bigr)\HT,
\]
where
\[
\dtail(\Tor):=\frac32\sum_{\substack{z\ge17\\ z\ \mathrm{odd}}}(3z-1)(z+1)2^{-z}\bigl(1+\sigma_z(\Tor)\bigr).
\]
Moreover $\frac32\sum_{z\ge17,\,z\ \mathrm{odd}}(3z-1)(z+1)2^{-z}=\frac{61}{4096}$, and for every rectangular torus whose shorter period
is at least $L>0$,
\[
\dtail(\Tor)-\frac{61}{4096}\le\frac32\sum_{\substack{z\ge17\\ z\ \mathrm{odd}}}(3z-1)(z+1)2^{-z}\,\Xi_z(L)
\le\sum_{(j,k)\in\Z^2\setminus0}\Omega\bigl(L\sqrt{j^2+k^2}\bigr),
\]
where $\Omega(r):=36\,e^{-r^2/16}+9\,(y^2+5y+8)\,2^{-y}$ with $y:=r^2/20$. The middle term is nonincreasing in $L$, and the right
side is $O(L^4\,2^{-L^2/20})$ as $L\to\infty$.
\end{proposition}

\begin{proof}
By Theorem~\ref{thm:square} and Proposition~\ref{prop:parts}(b), with $e_z=0$ for $z\ge16$,
\[
\LL_3(X_3(h))-\KT^{(3)}=\sum_{z\le15}\bigl(q_z-\tfrac23e_z\bigr)\LL_3\Pp_3(Y_z)+\sum_{z\ge16}q_z\,\LL_3\Pp_3(Y_z)-\mathcal C_3 .
\]
\emph{Near part.} $Y_z=0$ for $z\in\{0,1,3\}$ (Lemma~\ref{lem:qz}), and $q_z-\frac23e_z>0$ for $z\in\{2,4,5,\dots,15\}$ by \eqref{eq:C3}; so
the first sum is $\LL_3\Pp_3$ of a positive finite-rank relative operator tensored with $\one$, which is positive by
Lemmas~\ref{lem:Pk}(c) and~\ref{lem:lift}(iii). \emph{Tail.} For even $z\ge16$, $q_z>0$ and the term is positive. For odd $z\ge17$,
$q_z=-(3z-1)2^{-z}$, and by \eqref{eq:yz}, $B(y_z(R))=\sum_ps_{z,z,p}\,c(\chi^R_{z-p})B(\ut_p(R))$ (by \eqref{eq:Bwedge}; $s$ is real); by
Cauchy--Schwarz with $\sum_ps_{z,z,p}^2=1$ and $\norm{c(\chi_m^R)}^2=\Gamma(R)_{mm}\le1+\sigma_z(\Tor)$ for $m\le z$,
\[
0\le\ip\Psi{\LL_3\Pp_3(Y_z)\Psi}\le\frac3{2\pi}\int_\cell(1+\sigma_z)\sum_{p\le z}\norm{B(\ut_p(R))\Psi}^2\dd R=\frac32(1+\sigma_z)(z+1)\ip\Psi{\HT\Psi},
\]
the last step by Corollary~\ref{cor:pairmoyal}. Summing over odd $z\ge17$ gives $-\dtail\HT$, and $\mathcal C_3$ is bounded in
Proposition~\ref{prop:parts}(b). \emph{The sum.} With $z=17+2m$, $(3z-1)(z+1)=900+208m+12m^2$, and $\sum_m4^{-m}=\frac43$,
$\sum_mm4^{-m}=\frac49$, $\sum_mm^24^{-m}=\frac{20}{27}$ give $\frac32\cdot2^{-17}\cdot\frac{3904}3=\frac{61}{4096}$. \emph{The excess.}
$\sigma_z(\Tor)\le\Xi_z(L)$ by Lemma~\ref{lem:Ebounds}(c), and $\Xi_z(L)$ is nonincreasing in $L$. For $z\ge1$ and $r^2\ge20z$,
$\mathfrak g_z(r)\le e^{-r^2/16}$: with $r^2=2zt$, $t\ge10$, and $z!\ge(z/e)^z$, $\log\mathfrak g_z(r)\le z(1+\log(t+1)-t/2)\le-zt/8$, since
$1+\log(t+1)\le3t/8$ for $t\ge10$. Exchange the sums over $z$ and over the lattice points, bound the odd $z\ge17$ by all $z\ge1$
and $\frac32(3z-1)(z+1)$ by $\frac92z(z+1)$, and fix a lattice point with $r=L\sqrt{j^2+k^2}$. The terms with $z<r^2/20$ have
$\min\{1,\mathfrak g_z(r)\}\le e^{-r^2/16}$ and contribute at most $\frac92\sum_{z\ge1}z(z+1)2^{-z}e^{-r^2/16}=36e^{-r^2/16}$. The terms
with $z\ge z_0:=\lceil r^2/20\rceil$ have $\min\{1,\mathfrak g_z\}\le1$ and contribute at most
$\frac92\sum_{z\ge z_0}z(z+1)2^{-z}=9(z_0^2+3z_0+4)2^{-z_0}\le9(y^2+5y+8)2^{-y}$, using $y\le z_0\le y+1$. This is the bound by
$\sum\Omega$. For the asymptotics, $\Omega$ is decreasing for $r^2\ge20$, at most $4(n+1)$ lattice points have $j^2+k^2=n\ge2$, and
$\Omega(L\sqrt n)\le n^2e^{-(n-1)(\log2)L^2/20}\,\Omega(L)$ for $L^2\ge20$, so the sum is at most a constant times $\Omega(L)=O(L^42^{-L^2/20})$.
\end{proof}

The near part is the planar three-body inequality \eqref{eq:C3} transported to the torus by the order-preserving map
$\LL_3\Pp_3$; there is no error term between the torus and the plane on any block, frustrated or not.

% sections/08-bessel.tex
\section{The four-body Bessel bound}\label{sec:bessel}

The four-body coefficient of $\HT^2$ is $\LL_4(W_4(h\otimes h)W_4^*)$, a product of two pair operators on disjoint pairs, and its
comparison with $\KT^{(4)}$ is not exact: the periodized planar target $\mathcal T_4$ of Proposition~\ref{prop:parts}(d) involves only
the channels $D\le23$ of the relative coordinate of the two pair centres, and these channels are localized. The bound below
compares $\mathcal T_4$ with $\LL_4(W_4(h\otimes h)W_4^*)$ up to a factor $1+\Ec_{22}(\Tor)$; the proof reduces it to a Gram matrix of
the same kind as the local Gram matrix of Section~\ref{sec:local}, now for the pair centres.

\begin{lemma}[four-body Bessel bound]\label{lem:bessel}
Let $\mathcal Z_4:=\Pp_4\bigl(\sum_{D=1}^{23}\kb{\omega_{D,1}}{\omega_{D,1}}\otimes\one\bigr)$ on $\h^{\otimes4}$, and let $h\otimes h$ act on the
particles $(1,2)$ and $(3,4)$. Then
\[
\mathcal Z_4\le\bigl(1+\Ec_{22}(\Tor)\bigr)\,h\otimes h .
\]
\end{lemma}

\begin{proof}
Let $\Fib=\Fib_{\eta_1}:\h^{\otimes2}\to\U_2$ be the fibering map of Proposition~\ref{prop:orbital}, applied to the particles $(1,2)$ and to the
particles $(3,4)$; it maps $\h^{\otimes4}$ to $\U_2\otimes\U_2$, the two-particle space of the torus $\Tor_2=\R^2/\sqrt2\Lam$, whose two
particles are the pair centres $\zeta_{12}=(r_1+r_2)/\sqrt2$ and $\zeta_{34}=(r_3+r_4)/\sqrt2$. Their Jacobi coordinates are
$\rho_3=(\zeta_{12}-\zeta_{34})/\sqrt2$ and $\zeta=(\zeta_{12}+\zeta_{34})/\sqrt2$. For $d\ge0$ let $\Fib'_d$ be the fibering map of the torus
$\Tor_2$ with the relative state $\eta_d(\rho_3)$; it maps $\U_2\otimes\U_2$ to the centre-of-mass space $\U''$ of $\Tor_2$, the lowest
Landau level of $\R^2/2\Lam$. Since $\omega_{D,1}=\eta_1(\rho_1)\eta_1(\rho_2)\eta_{D-1}(\rho_3)$, Fubini gives
$\Fib_{\omega_{D,1}}=\Fib'_{D-1}(\Fib\otimes\Fib)$, and Lemma~\ref{lem:Pk}(c) gives
\[
\mathcal Z_4=\sum_{d=0}^{22}(\Fib\otimes\Fib)^*(\Fib'_d)^*\Fib'_d(\Fib\otimes\Fib)=(\Fib\otimes\Fib)^*\mathcal G^*\mathcal G(\Fib\otimes\Fib),\qquad h\otimes h=(\Fib\otimes\Fib)^*(\Fib\otimes\Fib),
\]
with $\mathcal G:=(\Fib'_d)_{d\le22}:\U_2\otimes\U_2\to\bigoplus_{d\le22}\U''$. Hence $\mathcal Z_4\le\norm{\mathcal G\mathcal G^*}\,h\otimes h$.

By \eqref{eq:Fstar}, applied on $\Tor_2$,
$(\Fib'_{d'})^*g=\sum_{\gamma'\in\sqrt2\Lam}\tl^{(2)}_{\gamma'}\bigl(\eta_{d'}(\rho_3)g(\zeta)\bigr)$ for $g\in\U''$, and here
$\tl^{(2)}_{\gamma'}=t^{(2)}_{\gamma'}$: $\Tor_2$ carries $2\Nphi$ flux quanta, so its lattice translation by $\gamma'=\sqrt2(jL_x,kL_y)$
has the sign $(-1)^{jk\cdot2\Nphi}=1$. A magnetic translation of the second
particle by $a$ acts as $t_{-a/\sqrt2}$ on $\rho_3$ and as $t_{a/\sqrt2}$ on $\zeta$ (as in Lemma~\ref{lem:jacobi}(ii)). Hence, writing
$\gamma'=\sqrt2\gamma$ with $\gamma\in\Lam$,
\[
\Fib'_d(\Fib'_{d'})^*=\sum_{\gamma\in\Lam}\ip{\eta_d}{t_{-\gamma}\eta_{d'}}\;t_\gamma\big|_{\U''},\qquad
\mathcal G\mathcal G^*=\sum_{\gamma\in\Lam}\Dis^{(22)}(\beta(-\gamma))\otimes t_\gamma\big|_{\U''},
\]
the series converging absolutely. Each $t_\gamma$, $\gamma\in\Lam$, is unitary on $\U''$ by
Lemma~\ref{lem:irred}, since $\Lam$ is contained in the lattice $2\Lam/(4\Nphi)$ of $\R^2/2\Lam$. The term $\gamma=0$ is the identity,
and $\norm{\Dis^{(22)}(\beta(-\gamma))}=\mu_{22}(\abs\gamma)$. So $\norm{\mathcal G\mathcal G^*}\le1+\Ec_{22}(\Tor)$.
\end{proof}

In these coordinates the relative coordinate $\rho_3$ of the two pair centres carries unit charge, and it equals the difference
of the two pair centres of mass; so the image displacements in the proof are those of a single electron, and $\Ec_{22}$ is the
same lattice sum as for the local Gram matrix of the orbitals $\chi_0^R,\dots,\chi_{22}^R$.

\begin{proposition}[four-body comparison]\label{prop:four}
On $\Fock$, for every rectangular torus, with $\Ec_m=\Ec_m(\Tor)$,
\begin{gather*}
\mathcal T_4\le(1+\Ec_{22})\,\LL_4\bigl(W_4(h\otimes h)W_4^*\bigr)\quad\text{and}\\
\LL_4\bigl(W_4(h\otimes h)W_4^*\bigr)-\KT^{(4)}\ge-(1+\Ec_{23})^2\bigl(1222\,\eps+2444\,\Ec_{22}\bigr)\HT .
\end{gather*}
\end{proposition}

\begin{proof}
$W_4\omega_{D,1}=\sqrt6A_4\omega_{D,1}$ with $A_4$ acting on the relative variables, so
$\kb{W_4\omega_{D,1}}{W_4\omega_{D,1}}\otimes\one=6A_4(\kb{\omega_{D,1}}{\omega_{D,1}}\otimes\one)A_4$, and by Lemma~\ref{lem:Pk}(b) its periodized
image is $6A_4\Pp_4(\kb{\omega_{D,1}}{\omega_{D,1}}\otimes\one)A_4$. The operator $\Pp_4(\kb{\omega_{D,1}}{\omega_{D,1}}\otimes\one)$ maps into
$\bigwedge^2\h\otimes\bigwedge^2\h$, and on that space $6A_4XA_4=W_4XW_4^*$ on $\bigwedge^4\h$ (Lemma~\ref{lem:lift}(iv)). Hence
$\mathcal T_4=\LL_4(W_4\mathcal Z_4W_4^*)$, and Lemma~\ref{lem:bessel}, conjugation by $W_4$ and the lift preserve order. For the second
inequality, Proposition~\ref{prop:parts}(d) gives $\KT^{(4)}\le\mathcal T_4+1222\eps(1+\Ec_{23})^2\HT$ and $\mathcal T_4\le2444(1+\Ec_{23})^2\HT$; so
\[
\LL_4(W_4(h\otimes h)W_4^*)-\KT^{(4)}\ge\frac{\mathcal T_4}{1+\Ec_{22}}-\mathcal T_4-1222\eps(1+\Ec_{23})^2\HT\ge-\Ec_{22}\,\mathcal T_4-1222\eps(1+\Ec_{23})^2\HT,
\]
using $\mathcal T_4\ge0$, and the bound on $\mathcal T_4$ finishes the proof.
\end{proof}

% sections/09-theorem.tex
\section{The theorem}\label{sec:theorem}

This section proves Theorems~\ref{thm:transferintro} and~\ref{thm:openai}. Theorem~\ref{thm:main} and Corollary~\ref{cor:uniform}, with
the rate of Remark~\ref{rem:rate} and the thresholds of Section~\ref{ssec:constants}, are the proof of Theorem~\ref{thm:openai};
Section~\ref{ssec:transferproof} proves Theorem~\ref{thm:transferintro} by the same argument with the data of the certificate left
general.

\subsection{The budget}

Recall $\gs=1-\etas-\frac{61}{4096}-1222\,\eps=4616733319001/10^{14}$, where the last equality is exact arithmetic with
$\etas=93527408868499/10^{14}$ and $\eps=3/10^6$, and $\Cth=72\sum_\nu\norm{\alpha^\nu}^2$ (Appendix~\ref{app:certificate}).

\begin{theorem}[uniform gap on rectangular tori]\label{thm:main}
For every rectangular torus $\Tor$, on the whole Fock space $\Fock$,
\[
\HT^2\ge\KT+\gT\,\HT\ge\gT\,\HT,
\]
where, with $\Ec_m=\Ec_m(\Tor)$, $\lambda_{\min}$ the smallest eigenvalue of the pair Gram matrix and $\dtail=\dtail(\Tor)$,
\begin{equation}\label{eq:gT}
\begin{aligned}
\gT&:=\lambda_{\min}-\etas-\dtail-\Cth\,\Ec_8(1+\Ec_8)-(1+\Ec_{23})^2\bigl(1222\,\eps+2444\,\Ec_{22}\bigr)\\
&\phantom:=\gs-(1-\lambda_{\min})-\Bigl(\dtail-\frac{61}{4096}\Bigr)-\Cth\,\Ec_8(1+\Ec_8)\\
&\phantom{:=\gs}-1222\,\eps\bigl((1+\Ec_{23})^2-1\bigr)-2444\,\Ec_{22}(1+\Ec_{23})^2 .
\end{aligned}
\end{equation}
\end{theorem}

\begin{proof}
By Theorem~\ref{thm:square} and Proposition~\ref{prop:parts},
\[
\HT^2-\KT=\bigl[\LL_2(h^2)-\etas\HT\bigr]+\bigl[\LL_3(X_3(h))-\KT^{(3)}\bigr]+\bigl[\LL_4(W_4(h\otimes h)W_4^*)-\KT^{(4)}\bigr].
\]
By Proposition~\ref{prop:projector}, $h^2\ge\lambda_{\min}h$, and $\LL_2$ preserves order with $\LL_2(h)=\HT$; so the first bracket is at
least $(\lambda_{\min}-\etas)\HT$. The second is bounded by Proposition~\ref{prop:three} and the third by Proposition~\ref{prop:four}.
Finally $\KT\ge0$ and $\HT\ge0$.
\end{proof}

All the constants in \eqref{eq:gT} depend on the torus only through its lattice of periods, and none depends on the
particle number. They are small when the shorter period is large:

\begin{corollary}[uniform bound]\label{cor:uniform}
For $L\ge16$ let
\begin{multline*}
\gl(L):=\gs-\bigl[2\Dlev(L)+\Ups(L,L)\bigr]-\frac32\sum_{\substack{z\ge17\\ z\ \mathrm{odd}}}\frac{(3z-1)(z+1)}{2^z}\,\Xi_z(L)
-\Cth\,\Ec_8(L)(1+\Ec_8(L))\\
-1222\,\eps\bigl((1+\Ec_{23}(L))^2-1\bigr)-2444\,\Ec_{22}(L)(1+\Ec_{23}(L))^2,
\end{multline*}
with $\Ec_m(L)$ and $\Xi_z(L)$ as in Lemma~\ref{lem:Ebounds}. Then $\gT\ge\gl(L)$ for every rectangular torus whose shorter period
is at least $L$, the function $\gl$ is nondecreasing on $[16,\infty)$, and $\gl(L)\to\gs$ as $L\to\infty$.
\end{corollary}

\begin{proof}
By Lemma~\ref{lem:levels}, $1-\lambda_{\min}\le\dG\le\Dlev(L_x)+\Dlev(L_y)+\Ups(L_x,L_y)$. The function $\Dlev$ is decreasing on
$(\sqrt{12},\infty)$, and every term $\abs{\mathfrak f(s_{jk})}=(s_{jk}-1)e^{-s_{jk}/2}$ of $\Ups$ is decreasing in $L_x$ and $L_y$, since off
$\mathcal X$ it has $s_{jk}\ge9/2>3$; so the bound is at most $2\Dlev(L)+\Ups(L,L)$ and is nonincreasing in $L$. The tail term is
bounded by Proposition~\ref{prop:three}, and $\Ec_m(\Tor)\le\Ec_m(L)$ for $m\in\{8,22,23\}$ by Lemma~\ref{lem:Ebounds}(a); all
subtracted terms are nonincreasing in $L$. As $L\to\infty$, $\Dlev(L)\to0$ and $\Ups(L,L)\to0$ (Lemma~\ref{lem:levels}(c)), the tail
excess tends to zero (Proposition~\ref{prop:three}), and $\Ec_m(L)\to0$ (Lemma~\ref{lem:Ebounds}(b)).
\end{proof}

\begin{remark}[the rate of convergence]\label{rem:rate}
As $L\to\infty$,
\[
0\le\gs-\gl(L)=O\bigl(L^4\,2^{-L^2/20}\bigr).
\]
The rate is that of the majorant $\Omega$ of Proposition~\ref{prop:three} for the tail term, and it is not claimed to be sharp.
The same rate holds for the transfer principle, Theorem~\ref{thm:transferintro}: for every planar certificate of the same form,
the bound $\etr(L)$ on $\etr(\Tor)$ obtained in the proof of Theorem~\ref{thm:transfer} is $O(L^4\,2^{-L^2/20})$, with a
constant that depends on the certificate.
\end{remark}

\begin{proof}
$\gs-\gl(L)$ is the sum of the five nonnegative terms subtracted in the definition of $\gl$. By the definition of $\Dlev$ and
Lemma~\ref{lem:levels}(c), $2\Dlev(L)+\Ups(L,L)=O(L^2e^{-L^2/8})$. By Lemma~\ref{lem:Ebounds}(b),
$\Ec_m(L)\le\widehat\Ec_m(L)=O(L^{2m}e^{-L^2/4})$, so the three terms containing $\Ec_8$, $\Ec_{22}$ and $\Ec_{23}$ are
$O(L^{46}e^{-L^2/4})$. By Proposition~\ref{prop:three}, the tail term
$\frac32\sum_{z\ge17,\,z\ \mathrm{odd}}(3z-1)(z+1)2^{-z}\,\Xi_z(L)$ is at most $\sum_{(j,k)\neq0}\Omega(L\sqrt{j^2+k^2})$, which is
$O(L^4\,2^{-L^2/20})$. Since $2^{-L^2/20}=e^{-(\log2)L^2/20}$ and $(\log2)/20<1/8<1/4$, the other four terms are
$o(L^4\,2^{-L^2/20})$. The proof of Theorem~\ref{thm:transfer} bounds the terms of $\etr(\Tor)$ by the same three results, with
$\Ec_{M_c}$, $\Ec_{D_c-1}$ and $\Ec_{D_c}$ in place of $\Ec_8$, $\Ec_{22}$ and $\Ec_{23}$, and the tail by the same majorant (whose
derivation bounds the sum over all $z\ge1$); this gives the second statement.
\end{proof}

\subsection{The constants}\label{ssec:constants}

Table~\ref{tab:constants} lists the terms of $\gl(L)$. The first column after $L$ is the bound $2\Dlev(L)+\Ups(L,L)$ on
$1-\lambda_{\min}$; the binding term is the four-body Bessel term, followed by the tail excess, while the pair-level and
non-orthonormality terms are negligible.

\begin{table}[ht]
\centering\small
\begin{tabular}{@{}rccccc@{}}
\toprule
$L$ & $2\Dlev(L)+\Ups(L,L)$ & $\Ec_8(L)$ & $\Ec_{22}(L)$ & $\Ec_{23}(L)$ & tail excess\\
\midrule
16.5 & $4.5\cdot10^{-13}$ & $2.1\cdot10^{-17}$ & $1.50\cdot10^{-5}$ & $5.73\cdot10^{-5}$ & $5.42\cdot10^{-4}$\\
17   & $5.8\cdot10^{-14}$ & $5.3\cdot10^{-19}$ & $1.13\cdot10^{-6}$ & $4.73\cdot10^{-6}$ & $1.93\cdot10^{-4}$\\
17.5 & $7.2\cdot10^{-15}$ & $1.2\cdot10^{-20}$ & $6.95\cdot10^{-8}$ & $3.19\cdot10^{-7}$ & $8.38\cdot10^{-5}$\\
18   & $8.2\cdot10^{-16}$ & $2.2\cdot10^{-22}$ & $3.56\cdot10^{-9}$ & $1.77\cdot10^{-8}$ & $4.22\cdot10^{-5}$\\
19   & $9.0\cdot10^{-18}$ & $5.2\cdot10^{-26}$ & $5.45\cdot10^{-12}$ & $3.16\cdot10^{-11}$ & $4.05\cdot10^{-6}$\\
20   & $7.6\cdot10^{-20}$ & $7.1\cdot10^{-30}$ & $4.18\cdot10^{-15}$ & $2.79\cdot10^{-14}$ & $4.95\cdot10^{-7}$\\
\bottomrule
\end{tabular}

\medskip
\begin{tabular}{@{}rcccc@{}}
\toprule
$L$ & $\Cth\Ec_8(1+\Ec_8)$ & $1222\,\eps\bigl((1+\Ec_{23})^2-1\bigr)$ & $2444\,\Ec_{22}(1+\Ec_{23})^2$ & $\gl(L)$\\
\midrule
16.5 & $6.0\cdot10^{-15}$ & $4.2\cdot10^{-7}$ & $3.67\cdot10^{-2}$ & $0.0089238941$\\
17   & $1.5\cdot10^{-16}$ & $3.5\cdot10^{-8}$ & $2.75\cdot10^{-3}$ & $0.0432229058$\\
17.5 & $3.3\cdot10^{-18}$ & $2.3\cdot10^{-9}$ & $1.70\cdot10^{-4}$ & $0.0459136282$\\
18   & $6.2\cdot10^{-20}$ & $1.3\cdot10^{-10}$ & $8.70\cdot10^{-6}$ & $0.0461164613$\\
19   & $1.5\cdot10^{-23}$ & $2.3\cdot10^{-13}$ & $1.33\cdot10^{-8}$ & $0.0461632655$\\
20   & $2.0\cdot10^{-27}$ & $2.0\cdot10^{-16}$ & $1.02\cdot10^{-11}$ & $0.0461668382$\\
\bottomrule
\end{tabular}
\caption{The terms of the uniform bound $\gl(L)$ of Corollary~\ref{cor:uniform}. The tail excess is
$\frac32\sum_{z\ge17\ \mathrm{odd}}(3z-1)(z+1)2^{-z}\,\Xi_z(L)$, and $\gs=0.04616733319001$.}
\label{tab:constants}
\end{table}

\paragraph{Thresholds.} The computed function $\gl$ crosses $0$ at $L=16.4562$, $1/25$ at $L=16.8566$, and $\gs-10^{-5}$ at
$L=18.6843$ (magnetic lengths, to the precision shown). Hence $\gl(L)>0$ for $L\ge16.457$, $\gl(L)>1/25$ for $L\ge16.857$, and
$\gl(L)\ge\gs-10^{-5}$ for $L\ge18.685$; by Corollary~\ref{cor:uniform} the same holds for $\gT$ on every rectangular torus whose
shorter period is at least these values.

\begin{remark}[on the evaluation of the constants]\label{rem:evaluation}
Every quantity in $\gl(L)$ is given by an explicit expression: $\Dlev$ is elementary, $\Ups(L,L)$ and $\Xi_z(L)$ are lattice sums of
elementary functions with Gaussian decay, the sum over odd $z$ converges geometrically (Proposition~\ref{prop:three}), and
$\mu_m(r)$ is the spectral norm of an explicit $(m+1)\times(m+1)$ matrix whose entries are Laguerre polynomials times Gaussians
(Lemma~\ref{lem:bargmann}(iii)). The reduction from the supremum over all rectangular tori to the evaluation at the radii $L\sqrt{j^2+k^2}$ is
rigorous by Lemma~\ref{lem:Ebounds}(a), whose proof shows that for $r\ge16$ every entry of $\Dis^{(m)}(r/\sqrt2)$, $m\le23$, has a fixed
sign and a modulus decreasing in $r$, so that the spectral norm is a monotone function of the entries. What remains is a
numerical evaluation. The entries of Table~\ref{tab:constants} were computed in $60$-digit floating-point arithmetic from these
expressions, with $\gs$, $61/4096$ and $\etas$ as exact rationals, the displacement entries from the Laguerre representation, and
the spectral norms in double precision cross-checked against $60$-digit eigenvalue computations (relative agreement
$3\cdot10^{-16}$); the lattice sums and the series over $z$ were truncated where the omitted terms are far below the displayed
digits. A second, independent implementation with different truncations reproduces every entry to the digits it reports (four
to seven significant figures for $\gl$). The values are therefore
floating-point evaluations of rigorous closed-form bounds, not interval-arithmetic certificates, and the thresholds were located
by bisection of the computed function.
\end{remark}

\subsection{Consequences}

\begin{corollary}[spectrum and kernel]\label{cor:spectrum}
Let $\Tor$ be a rectangular torus with $\gT>0$; by Corollary~\ref{cor:uniform} this holds whenever the shorter period is at least
$16.457$.
\begin{enumerate}[(a)]
\item For every $N$, $\spec\bigl(\HT|_{\bigwedge^N\h}\bigr)\subset\{0\}\cup[\gT,\infty)$.
\item At $\Nphi=3N$, $\HT|_{\bigwedge^N\h}\ge\gT(\one-P_N)$, where $P_N$ is the orthogonal projection onto the three-dimensional span of the
Laughlin states $\Psi_0,\Psi_1,\Psi_2$ of Theorem~\ref{thm:kernel}, and $\HT|_{\bigwedge^{N'}\h}\ge\gT$ for every $N'>N$.
\item If $\dG<1$, the projector form satisfies $H'^2\ge(\gT/\lambda_{\max})H'$.
\end{enumerate}
In particular, on every rectangular torus whose shorter period is at least $16.857$, the gap above the three-dimensional kernel at
$\Nphi=3N$ is larger than $1/25$, uniformly in $N$ and in the aspect ratio.
\end{corollary}

\begin{proof}
(a) $\HT\ge0$ preserves the particle number, and an eigenvalue $\lambda$ satisfies $\lambda^2\ge\gT\lambda$ by Theorem~\ref{thm:main}.
(b) By (a) and the spectral theorem, $\HT\ge\gT(\one-P_{\ker})$ on each sector; Theorem~\ref{thm:kernel} identifies the kernel at
$N$ particles and Proposition~\ref{prop:overfill} shows that it is zero for $N'>N$. (c) By Proposition~\ref{prop:projector},
$H'\ge\lambda_{\max}^{-1}\HT\ge(\gT/\lambda_{\max})(\one-P_{\ker})$ with $\ker H'=\ker\HT$.
\end{proof}

At the square aspect ratio the shorter period is $L=\sqrt{2\pi\Nphi}$, so $\gl>0$ for $\Nphi\ge44$ and $\gl>1/25$ for $\Nphi\ge46$.
At filling one third, $\Nphi=3N$: for $N=15$ ($L=16.815$) the bound is $\gl\ge0.0385$, and for $N=16$ ($L=17.366$) it is
$\gl\ge0.0456>1/25$. On elongated tori the shorter period is smaller at the same $N$, and correspondingly larger systems are
needed.

\subsection{The transfer principle}\label{ssec:transferproof}

It remains to prove Theorem~\ref{thm:transferintro}, for a planar certificate of the same form, with the inequalities
\eqref{eq:certgeneral} and the constant $\gamma$ of the definition in Section~\ref{ssec:transfer}. For the certificate
of~\cite{OAI:A-Fock-space-inequality-and-the-Laughlin-spectral-gap-September-24-2026}, $P_c=7$, $M_c=8$, $Z_c=15$, $D_c=23$,
$\delta_\infty(15)=61/4096$, $N_4=1222$ and $\gamma=\gs$. We restate the theorem in its full form.

\begin{theorem}[the transfer principle; Theorem~\ref{thm:transferintro} restated]\label{thm:transfer}
Given a planar certificate of the same form with constant $\gamma$, for every rectangular torus, on $\Fock$,
\[
\HT^2\ge(\gamma-\etr(\Tor))\,\HT,
\]
where, with $\Ec_m=\Ec_m(\Tor)$ and $\sigma_z=\sigma_z(\Tor)$,
\begin{multline*}
\etr(\Tor):=(1-\lambda_{\min})+\frac32\sum_{\substack{z>Z_c\\ z\ \mathrm{odd}}}\frac{(3z-1)(z+1)}{2^z}\sigma_z
+(M_c+1)(P_c+1)\sum_\nu\norm{\alpha^\nu}^2\,\Ec_{M_c}(1+\Ec_{M_c})\\
+\eps N_4\bigl((1+\Ec_{D_c})^2-1\bigr)+2N_4\,\Ec_{D_c-1}(1+\Ec_{D_c})^2 .
\end{multline*}
For every rectangular torus whose shorter period is at least $L$, $\etr(\Tor)\le\etr(L)$ with $\etr(L)\to0$ as $L\to\infty$.
In particular, if $\gamma>0$, then for every $\gamma'<\gamma$ there is $L_0$ such that $\HT^2\ge\gamma'\HT$ on the whole Fock space of every
rectangular torus whose shorter period is at least $L_0$, uniformly in the particle number and the aspect ratio.
\end{theorem}

\begin{proof}
The proofs of Proposition~\ref{prop:parts}, Proposition~\ref{prop:three}, Lemma~\ref{lem:bessel}, Proposition~\ref{prop:four} and
Theorem~\ref{thm:main} use the specific certificate only through the data of the definition in Section~\ref{ssec:transfer}, and go
through verbatim with the following changes. In Proposition~\ref{prop:parts}, the modes in $(i,k)$ and $j$ are at most $M_c$, each pair label $p\le P_c$
occurs with at most $M_c+1$ values of $j$, and $\frac1\pi\int\sum_{p\le P_c}\norm{B(\ut_p)\Psi}^2=(P_c+1)\ip\Psi{\HT\Psi}$, which gives the
constant $(M_c+1)(P_c+1)\sum_\nu\norm{\alpha^\nu}^2$ with $\Ec_{M_c}$; the three-body blocks are $z\le Z_c$, the four-body blocks
$D\le D_c$, with mode indices at most $D_c$ (whence $\Ec_{D_c}$) and $\sum_{D\le D_c}n_D=N_4$. In Proposition~\ref{prop:three} the near
part uses the first family in \eqref{eq:certgeneral} (non-strict inequalities suffice, and the blocks $z\in\{0,1,3\}$ vanish), the
even blocks $z>Z_c$ are positive and the odd ones are bounded as before. In Lemma~\ref{lem:bessel} the pair-centre channels are
$d\le D_c-1$, whence $\Ec_{D_c-1}$, and Proposition~\ref{prop:four} gives $-(1+\Ec_{D_c})^2(\eps N_4+2N_4\Ec_{D_c-1})\HT$. Collecting the
terms as in \eqref{eq:gT} gives the first claim. For the second, $1-\lambda_{\min}\le2\Dlev(L)+\Ups(L,L)$ (proof of
Corollary~\ref{cor:uniform}), the tail is bounded by Proposition~\ref{prop:three}, and $\Ec_m(\Tor)\le\widehat\Ec_m(L)$ for
$L^2\ge2m$ by Lemma~\ref{lem:Ebounds}(b); each bound tends to zero.
\end{proof}

% sections/10-outlook.tex
\section{What is not proved, and what remains}\label{sec:outlook}

\paragraph{The transfer principle.} Theorem~\ref{thm:transferintro} separates the two halves of the argument: the planar certificate,
a finite computation in the plane, and its passage to rectangular tori, which the theorem carries out for every planar
certificate of the form of~\cite{OAI:A-Fock-space-inequality-and-the-Laughlin-spectral-gap-September-24-2026}, with errors
bounded by explicit functions of the shorter period that decay at a Gaussian rate. It is the part of this paper on which
further work can build, in particular on the spinful hollow-core model discussed at the end of this section. For the
certificate of~\cite{OAI:A-Fock-space-inequality-and-the-Laughlin-spectral-gap-September-24-2026} the bounds are explicit:
evaluated numerically (Remark~\ref{rem:evaluation}), they give the thresholds $16.457$, $16.857$ and $18.685$ magnetic lengths
for a positive constant, a constant above $1/25$ and a constant within $10^{-5}$ of $\gs$ (Theorem~\ref{thm:openai}), and
the lower bound $\gl(L)$ tends to $\gs$ at the Gaussian rate $O(L^4\,2^{-L^2/20})$ of Remark~\ref{rem:rate}. The candidate proof of~\cite{HouWu} claims a uniform
bound of broader scope, on flat tori of every lattice shape and with every boundary twist, above an existential threshold;
the bounds here are restricted to rectangular tori and are explicit.

\paragraph{Oblique tori.} Theorem~\ref{thm:main} is proved for rectangular tori, uniformly in the aspect ratio; it says nothing
about tori with an oblique lattice of periods. The kernel theorem (Theorem~\ref{thm:kernel} and Proposition~\ref{prop:overfill})
holds for every modular parameter, but for non-rectangular $\tau$ it is a statement about the operators $B^{(\tau)}_n$ of
\eqref{eq:Bntau}, and we have not identified $\sum_n(B_n^{(\tau)})^*B_n^{(\tau)}$ with the periodized planar pair projector on an oblique
torus. The definition $\HT:=\LL_2(\Pp_2(P_1))$ makes sense for every lattice, and we expect most of the argument to carry over: the
proof of the Moyal identity (Theorem~\ref{thm:moyal}) uses only the irreducibility of the fine-lattice magnetic translations of
the centre-of-mass torus and the vanishing of the averaged image terms, and the image-sum calculus of Section~\ref{sec:square}
never uses the shape of the lattice. What would have to be redone is the part tied to the rectangular Landau gauge: the orbital
form (Proposition~\ref{prop:orbital}), the pair Gram matrix and its levels (Lemmas~\ref{lem:gram} and~\ref{lem:levels}), the
irreducibility argument of Lemma~\ref{lem:irred} written in an oblique orbital basis, and the lattice sums $\Ec_m$ and $\Xi_z$, whose
first shell can contain six vectors instead of four. A bound uniform in the modular parameter would allow the modular parameter to
be varied within the gapped phase, which is the setting in which the Hall viscosity appears as an adiabatic curvature over the
space of flat metrics~\cite{ASZ,Read09}; this is a separate problem, not addressed here.

\paragraph{Sharper constants.} The threshold is set almost entirely by the four-body term $2444\,\Ec_{22}(1+\Ec_{23})^2$ and, after
it, by the tail excess (Table~\ref{tab:constants}). The first comes from bounding the periodized four-body target $\mathcal T_4$ by
the pair energy through Cauchy--Schwarz over all $n_D$ triples of each block (Proposition~\ref{prop:parts}(d)); the second from the
crude bound $\abs{\bra m\Dis\ket m}\le\min\{1,\mathfrak g_m\}$ on the diagonal entries of the local Gram matrix. A sharper estimate of
$\mathcal T_4$, or the exact Laguerre values in $\Xi_z$, would lower the threshold; we have not attempted either. The bound on the
pair-level term is already far below the others, and so is the non-orthonormality term $\Cth\Ec_8(1+\Ec_8)$.

\paragraph{Finite tori below the threshold.} Below the threshold the method gives nothing. For instance at the square aspect ratio
and filling one third the bound is positive only for $N\ge15$; smaller systems are accessible to exact diagonalization, which is a
different kind of statement.

\paragraph{The spinful case.} For spin-$\frac12$ fermions with the hollow-core (pure $V_1$) interaction, at $\Nphi=2N$ flux quanta,
the zero modes are the Haldane--Rezayi states. We expect a spinful transfer principle of the form of Theorem~\ref{thm:transferintro}: a
planar certificate of the same form for the spinful model, with constant $\gamma>0$, should give a uniform gap above the zero
modes on rectangular tori. What it would need is the spinful version of each step: the spinful pair Gram matrix and its levels; the torus Moyal
identities for pair states with spin labels, which should follow from the same Schur argument, since the magnetic translations
act on the orbital factor only; and the three-body coefficient of the square of the spinful Hamiltonian, whose exact
periodization should follow from Lemma~\ref{lem:product}, which does not depend on internal labels, but whose planar coefficients
replace the $q_z$ by spin-dependent blocks that have to be computed. None of this is done here. In the contrapositive, such a
principle would show that no planar certificate of this form can exist for the spinful hollow-core model if a uniform gap above
its zero modes on rectangular tori is excluded.


\appendix
% appendices/A-certificate.tex
\section{The certificate data}\label{app:certificate}

This appendix records the data of the planar certificate
of~\cite{OAI:A-Fock-space-inequality-and-the-Laughlin-spectral-gap-September-24-2026} that the proof consumes, in the form given
there (its Appendix~A and the file \texttt{verification/rows.txt} of its source directory). Nothing in this appendix is re-derived;
the inequalities \eqref{eq:C3} and \eqref{eq:C4} are proved in~\cite[Appendix~A]{OAI:A-Fock-space-inequality-and-the-Laughlin-spectral-gap-September-24-2026}
by exact rational arithmetic.

\subsection{The rows}

Let $P:=10^7$. Each row below starts with the integers $t_\nu$ and $\ell_\nu=P\lambda_\nu$; an entry \texttt{pj:a} specifies the
integer $a^\nu_{pj}$, with $p$ and $j$ read as single digits; omitted entries are zero, and indented lines continue the
preceding row. The coefficients are
\begin{equation}\label{eq:alphadata}
\alpha^\nu_{pj}=\frac{a^\nu_{pj}}P\sqrt{\frac{2^{p-t_\nu}\,t_\nu!\,(p+j-t_\nu)!}{p!\,j!}} ,
\end{equation}
\cite[Appendix~A.1, (39)]{OAI:A-Fock-space-inequality-and-the-Laughlin-spectral-gap-September-24-2026}; every nonzero entry has
$0\le p+j-t_\nu\le8$, so the support lies in $\mathcal S_\nu$.
{\footnotesize
\begin{verbatim}
0 7181518 05:-849651 07:-2958744 15:751659 16:-43084 17:-146216 32:720949 35:-127514
  42:222437 43:-262822 53:-117700 61:-55024 71:-166049
1 -112991 05:142772 06:-142342 07:527055 08:50742 16:197917 18:-18343 24:5764 27:-21183
  32:75832 35:-6115 45:2660 53:-9455 61:-27091 63:-3494 71:479 72:-15331
1 -4247311 05:-12207250 06:-14064191 07:-14063663 08:-15662679 16:483831 18:3987214
  24:-814558 27:888887 32:966462 35:145978 45:-171722 53:34350 61:389319 63:97603
  71:412447 72:-53786
2 3438746 14:6689266 15:6620205 16:6625065 17:4969272 18:4878238 27:-1387380 36:-586188
  37:344087 42:-1120634 45:-318530 46:559169 55:157455 61:1207103 63:211617 64:33240
  71:577001 73:-142578
5 -98645 05:-2719776 08:-8519913 14:-949523 17:-1952105 18:-1562606 28:-8919966 32:248826
  35:-1017623 37:3656005 38:-2081871 42:-33163 43:-305905 46:35677 47:-1320749
  48:-1040378 53:502094 55:1060633 58:-998003 61:-113617 63:184919 64:209239 65:-68198
  67:-266297 72:-606696 73:-563620 74:437345 76:-432189
5 -3160276 05:-1698129 08:10575875 14:-252576 17:-8002616 18:22592844 28:7370640
  32:289294 35:1400647 37:1987740 38:4030699 42:-1731401 43:-1034062 46:-1106371
  47:-1257051 48:1461735 53:1408149 55:1136847 58:1878501 61:-2838516 63:-734099
  64:1062416 65:11745 67:611312 72:-480776 73:20459 74:316925 76:37092
7 1441784 07:3441286 08:-6097155 37:-1002399 47:-251321 58:1143225 66:-17397 71:-1007464
  73:-662286 76:-344621 78:-1115284
\end{verbatim}
}

\subsection{The constants}

\begin{itemize}
\item $P^2\etas=\sum_\nu\ell_\nu^2=93527408868499$, so $\etas=93527408868499/10^{14}=0.93527408868499$.
\item $\eps=3/10^6$, and $\sum_{D=1}^{23}n_D=\sum_{m=1}^{12}m^2+\sum_{m=1}^{11}m(m+1)=1222$ with $n_D=\lfloor(D+1)^2/4\rfloor$.
\item $\frac32\sum_{z\ge17,\,z\ \mathrm{odd}}(3z-1)(z+1)2^{-z}=61/4096$ (Proposition~\ref{prop:three}).
\item $\gs=1-\etas-61/4096-1222\eps=4616733319001/10^{14}=0.04616733319001>1/25$.
\item $\sum_\nu\norm{\alpha^\nu}^2$ is a rational number determined by \eqref{eq:alphadata}; its value is $3.9764514174$ to the digits
shown, so $\Cth=72\sum_\nu\norm{\alpha^\nu}^2=286.3045\ldots$; only the bound $\Cth\le286.31$ is used.
\item Only the orbital modes $0,\dots,8$ and the pair labels $0,\dots,7$ occur in the rows; the three-body blocks are $z\le15$ and the
four-body blocks $D\le23$.
\end{itemize}

\subsection{The three-body inequalities}

For $z\in\{2,4,5,\dots,15\}$ the margins $m_z:=\frac32(3z-1)(-\frac12)^z-e_z$ are rational, with $e_z$ given by \eqref{eq:ez}.
In~\cite[Appendix~A.2]{OAI:A-Fock-space-inequality-and-the-Laughlin-spectral-gap-September-24-2026} they are computed exactly,
by a formula in which all radicals cancel; the integer parts $\lfloor10^6m_z\rfloor$ recorded there are
\begin{center}
\setlength{\tabcolsep}{4pt}\begin{tabular}{rrrrrrrrrrrrrr}
\toprule
$z$&2&4&5&6&7&8&9&10&11&12&13&14&15\\
\midrule
$\lfloor10^6m_z\rfloor$&1298313&211317&249&250&249&250&250&250&249&250&250&6470&250\\
\bottomrule
\end{tabular}
\end{center}
All are positive, which is \eqref{eq:C3}.

\subsection{The four-body inequalities}

For $1\le D\le23$ and $r,s\in\mathcal R_D$ put $u_r:=1/(r!(D-r)!)$. Then $(E_D)_{rs}=\sqrt{u_ru_s}\,Y_{rs}$ and
$(G_D)_{rs}=\sqrt{u_ru_s}\,Z_{rs}$ with rational matrices $Y$, $Z$ given by explicit finite sums over the rows,
and the matrix in \eqref{eq:C4} equals $D_{\sqrt u}MD_{\sqrt u}$ with $M=ZBZ$, $B=\mathrm{diag}((2\one_{r=1}+\eps)u_r)-D_uYD_u$
\cite[Appendix~A.3, (42)--(46)]{OAI:A-Fock-space-inequality-and-the-Laughlin-spectral-gap-September-24-2026}. Positivity of the rational symmetric
matrix $M$ is verified there through the signs of the coefficients of $\det(x\one+M)$, computed exactly by the Leverrier--Faddeev
recurrence: all coefficients after the leading one are nonnegative, so $M$ has no negative eigenvalue. The recorded sign patterns
(\texttt{+} positive, \texttt{0} zero, in descending degree) are
\begin{center}
\begin{tabular}{rl@{\qquad}rl@{\qquad}rl}
\toprule
$D$&signs&$D$&signs&$D$&signs\\
\midrule
1&\texttt{0}&9&\texttt{++000}&17&\texttt{++++++000}\\
2&\texttt{0}&10&\texttt{+0000}&18&\texttt{+++++0000}\\
3&\texttt{00}&11&\texttt{+++000}&19&\texttt{+++++++000}\\
4&\texttt{00}&12&\texttt{++0000}&20&\texttt{++++++0000}\\
5&\texttt{+00}&13&\texttt{++++000}&21&\texttt{++++++++000}\\
6&\texttt{000}&14&\texttt{+++0000}&22&\texttt{+++++++0000}\\
7&\texttt{+000}&15&\texttt{+++++000}&23&\texttt{+++++++++000}\\
8&\texttt{+000}&16&\texttt{++++0000}&&\\
\bottomrule
\end{tabular}
\end{center}
which is \eqref{eq:C4}.

% appendices/B-notation.tex
\section{Notation}\label{app:notation}

The table lists the symbols that recur across sections, with the place where each is defined. Symbols used in one section
only, and symbols local to a proof, are defined where they are used.

\begingroup\small
\begin{longtable}{@{}>{\raggedright\arraybackslash}p{0.25\textwidth}>{\raggedright\arraybackslash}p{0.52\textwidth}>{\raggedright\arraybackslash}p{0.16\textwidth}@{}}
\toprule
symbol & meaning & defined in\\
\midrule
\endfirsthead
\toprule
symbol & meaning & defined in\\
\midrule
\endhead
\bottomrule
\endfoot
$\Tor$, $\Lam$, $\cell$; $L_x$, $L_y$, $L$ & the torus $\R^2/\Lam$, its lattice of periods, a period cell; the periods; the shorter period
& Def.~\ref{def:torus}\\
$\Nphi$, $\kappa$ & number of flux quanta, $L_xL_y=2\pi\Nphi$; $\kappa=2\pi/L_y$ & Def.~\ref{def:torus}\\
$\h$, $\hpl$ & lowest Landau level of the torus; planar lowest Landau level & Def.~\ref{def:torus}, \eqref{eq:ta}\\
$w$, $a\wedge b$ & complex coordinate $x+\iu y$; $a_xb_y-a_yb_x$ & Section~\ref{sec:setting}\\
$\phi_n$, $\psi_j$, $[n]$ & Landau-gauge orbitals ($n\in\Z$); torus orbitals ($j\in\Z_{\Nphi}$); class of $n$ in $\Z_{\Nphi}$
& Lem.~\ref{lem:orbitals}\\
$\mathcal H_m$ & normalized Hermite functions & \eqref{eq:landaurep}\\
$t_a$, $\tl_\gamma$, $t_\delta$; $\thx$, $\thy$ & magnetic translations; lattice translations, $\gamma\in\Lam$; fine-lattice translations,
$\delta\in\Lam/\Nphi$; the translations $\psi_j\mapsto\psi_{j+1}$, $\psi_j\mapsto e^{2\pi\iu j/\Nphi}\psi_j$ & \eqref{eq:ta}, Lem.~\ref{lem:magtrans}\\
$\Hb$, $\ket m$, $\mathcal U$ & Bargmann space and its monomial basis; the unitary $\Hb\to\hpl$ & Lem.~\ref{lem:bargmann}\\
$\Dis(\beta)$, $\Dis^{(m)}(\beta)$, $\beta(a)$ & displacement operator; its truncation to the levels $0,\dots,m$;
$(a_x-\iu a_y)/\sqrt2$ & \eqref{eq:displacement}, Lem.~\ref{lem:bargmann}\\
$\eta_m$, $\eta_m^R$ & symmetric-gauge orbitals about the origin and about $R$ & Lem.~\ref{lem:bargmann}, \eqref{eq:etaR}\\
$\Ppl(r,r')$ & kernel of the orthogonal projection onto $\hpl$ & Lem.~\ref{lem:bargmann}(iv)\\
$\Pk^{(k)}$, $\Pk$ & periodization of planar $k$-particle functions; $\Pk=\Pk^{(1)}$ & Lem.~\ref{lem:unfold}\\
$\Fock$, $\Ff(u)$ & fermionic Fock space over $\h$; second quantization of a unitary $u$ & \S\ref{ssec:fock}\\
$B(v)$, $c(f)$, $c_j$ & adjoint of wedge multiplication by $v$; annihilation operators, $c_j=c(\psi_j)$ & \eqref{eq:Bwedge}\\
$A_n$; $P_\sigma$ & antisymmetrizer on $\h^{\otimes n}$; permutation of the variables & \S\ref{ssec:fock}; before Lem.~\ref{lem:unfold}\\
$\LL_k$, $X_S$ & lift of an operator on $\bigwedge^k\h$ to $\Fock$; $X$ acting on the tensor positions in $S$ & \eqref{eq:lift}\\
$W_3$, $W_4$ & the wedge maps $\bigwedge^2\h\otimes\h\to\bigwedge^3\h$ and $\bigwedge^2\h\otimes\bigwedge^2\h\to\bigwedge^4\h$
& Lem.~\ref{lem:lift}(iv)\\
$\zeta$, $\rho$; $\Tor_k$, $\U_k$; $\tl^{\otimes k}_\gamma$ & centre-of-mass and relative Jacobi coordinates; the centre-of-mass torus
$\R^2/\sqrt k\Lam$ and its lowest Landau level; the simultaneous lattice translation & Lem.~\ref{lem:jacobi}\\
$\Pp_k(Y)$, $\mathsf D$ & periodized image of an admissible operator $Y$; a fundamental domain of the diagonal action
& Def.~\ref{def:admissible}\\
$\Fib_a$ & fibering map, the partial inner product with the relative state $a$ & Lem.~\ref{lem:Pk}(c)\\
$P_1$, $h$, $\HT$, $\VT$ & planar projector onto relative angular momentum one; $h=\Pp_2(P_1)$; the torus $V_1$ Hamiltonian
$\HT=\LL_2(h)$; $\VT=\ran h$ & Def.~\ref{def:HT}\\
$c_0$, $f(d)$, $v_n$, $\bcen n$ & $\kappa^3/\sqrt{2\pi}$; $d\,e^{-\kappa^2d^2/4}$; pair functions; pair operators
& Prop.~\ref{prop:orbital}\\
$\tau$, $B_n^{(\tau)}$ & modular parameter; theta-function pair operators & \eqref{eq:Bntau}\\
$\Psi_j$ & the three Laughlin states at $\Nphi=3N$ & Thm.~\ref{thm:kernel}\\
$G$, $\mathfrak f$, $s_{jk}$ & pair Gram matrix; $\mathfrak f(s)=(1-s)e^{-s/2}$; $\frac14(j^2L_x^2+k^2L_y^2)$ & Lem.~\ref{lem:gram}\\
$\Dlev$, $\Ups$; $\lambda_{\min}$, $\lambda_{\max}$, $\dG$ & $\Dlev(\ell)=2(\ell^2/4-1)e^{-\ell^2/8}$; remainder of the
pair Gram levels; extreme eigenvalues of $G$ and their deviation from $1$ & Lem.~\ref{lem:levels}\\
$P_V$, $H'$ & projection onto $\VT$; $\LL_2(P_V)$ & Prop.~\ref{prop:projector}\\
$X_3(h)$ & three-body coefficient of $\LL_2(h)^2$ & Lem.~\ref{lem:square}\\
$\chi^R_m$, $v_p^{*R}$, $\ut_p(R)$ & periodized orbitals about $R$; planar pair states; periodized pair states
& Lem.~\ref{lem:localmodes}\\
$\Gamma(R)$, $\Gamma^{(m)}(R)$; $\mu_m$, $\Ec_m(\Tor)$ & local Gram matrix and its truncation; norms of truncated displacement
matrices and their lattice sum & Lem.~\ref{lem:localgram}\\
$\Ec_m(L)$, $\widehat\Ec_m(L)$, $\mathfrak g_m$; $\sigma_z$, $\Xi_z$ & bounds on $\Ec_m(\Tor)$; the majorant
$(r^2/2+m)^me^{-r^2/4}/m!$; excess of the diagonal of $\Gamma(R)$ and its bound & Lem.~\ref{lem:Ebounds}\\
$\varphi_p^R$ & centre-of-mass orbital of degree $p$ about $R$ & before Thm.~\ref{thm:moyal}\\
$s^{(u)}_{z,T,p}$, $s_{z,T,p}$; $\xi_z$, $\ket{z,T}^R$ & coupling coefficients, $s_{z,T,p}=s^{(1/\sqrt3)}_{z,T,p}$; three-body
relative states and coupled basis & \eqref{eq:coupling}, Lem.~\ref{lem:coupled3}\\
$\omega_{D,r}$, $\ket{D,T;r}^R$, $S_r$, $\mathcal R_D$ & four-body relative states, coupled basis and coefficients; the odd
integers in $[1,D]$ & Lem.~\ref{lem:coupled4}\\
$\nu$, $t_\nu$, $\lambda_\nu$, $\alpha^\nu_{pj}$, $\mathcal S_\nu$ & the rows of the certificate and their data
& \S\ref{ssec:transfer}, before \eqref{eq:ez}, App.~\ref{app:certificate}\\
$\etas$, $\eps$; $e_z$, $E_D$, $G_D$ & $\sum_\nu\lambda_\nu^2$; the allowance $3/10^6$; the three-body coefficients, four-body
matrices and four-body Gram matrices of the certificate & \eqref{eq:ez}, \eqref{eq:ED}\\
$F_\nu(R)$, $\KT$ & the rows at the base point $R$; the comparison operator & Def.~\ref{def:KT}\\
$\KT^{(k)}$, $Y_z$, $\Cth$, $\mathcal C_3$, $\mathcal T_4$ & $k$-body part of $\KT$; $\kb{W_3\xi_z}{W_3\xi_z}\otimes\one$;
$72\sum_\nu\norm{\alpha^\nu}^2$; non-orthonormality term; periodized four-body target & Prop.~\ref{prop:parts}\\
$q_z$ & planar three-body coefficients $(3z-1)(-\frac12)^z$ & Lem.~\ref{lem:qz}\\
$\dtail$, $\Omega$ & tail constant of the three-body comparison; majorant of its excess & Prop.~\ref{prop:three}\\
$\gT$, $\gs$, $\gl$ & the torus constant; the planar constant; the uniform lower bound
& Thm.~\ref{thm:main}, Cor.~\ref{cor:uniform}\\
$\etr(\Tor)$, $\etr(L)$; $P_c$, $M_c$, $Z_c$, $D_c$, $N_4$ & the error of the transfer principle and its bound; the
cutoffs and the four-body count of a certificate of the same form & \S\ref{ssec:transfer}, Thm.~\ref{thm:transferintro}\\
\end{longtable}
\endgroup


\begin{thebibliography}{99}

\bibitem{OAI:A-Fock-space-inequality-and-the-Laughlin-spectral-gap-September-24-2026}
OpenAI.
\newblock A {F}ock-space inequality and the {L}aughlin spectral gap.
\newblock OpenAI Math Release preprint
  \href{https://github.com/openai/math/blob/main/preprints/A-Fock-space-inequality-and-the-Laughlin-spectral-gap-September-24-2026/A-Fock-space-inequality-and-the-Laughlin-spectral-gap-September-24-2026.pdf}{OAI:A-Fock-space-inequality-and-the-Laughlin-spectral-gap-September-24-2026},
  2026.

\bibitem{Laughlin}
R.~B. Laughlin.
\newblock Anomalous quantum {H}all effect: an incompressible quantum fluid with fractionally charged excitations.
\newblock \emph{Phys. Rev. Lett.} 50, 1395--1398 (1983).

\bibitem{Haldane}
F.~D.~M. Haldane.
\newblock Fractional quantization of the {H}all effect: a hierarchy of incompressible quantum fluid states.
\newblock \emph{Phys. Rev. Lett.} 51, 605--608 (1983).

\bibitem{HaldaneRezayiTorus}
F.~D.~M. Haldane and E.~H. Rezayi.
\newblock Periodic {L}aughlin--{J}astrow wave functions for the fractional quantized {H}all effect.
\newblock \emph{Phys. Rev. B} 31, 2529--2531(R) (1985).

\bibitem{ONDS}
G.~Ortiz, Z.~Nussinov, J.~Dukelsky, and A.~Seidel.
\newblock Repulsive interactions in quantum {H}all systems as a pairing problem.
\newblock \emph{Phys. Rev. B} 88, 165303 (2013).

\bibitem{NWY}
B.~Nachtergaele, S.~Warzel, and A.~Young.
\newblock Spectral gaps and incompressibility in a $\nu=1/3$ fractional quantum {H}all system.
\newblock \emph{Commun. Math. Phys.} 383, 1093--1149 (2021).

\bibitem{WY}
S.~Warzel and A.~Young.
\newblock The spectral gap of a fractional quantum {H}all system on a thin torus.
\newblock \emph{J. Math. Phys.} 63, 041901 (2022).

\bibitem{Gao}
Q.~Gao.
\newblock Improving the lower bound on the {L}aughlin gap: {A}n {SDP} search.
\newblock Preprint, arXiv:2610.12700 (2026).

\bibitem{ScheerKhalaf}
M.~G. Scheer and E.~Khalaf.
\newblock A hierarchy of exact lower bounds on the {L}aughlin spectral gap.
\newblock Preprint, arXiv:2610.18854 (2026).

\bibitem{HouWu}
Uniform {L}aughlin gaps on large flat tori from finite {F}ock-space certificates.
\newblock Candidate proof for review, 16~October 2026, prepared by X.-J.~Hou with advice from Y.-H.~Wu.
\newblock \url{https://github.com/certainpineapple/Laughlin-torus-gap}, commit 22144aa.

\bibitem{Rougerie}
N.~Rougerie.
\newblock On the {L}aughlin function and its perturbations.
\newblock \emph{S\'eminaire Laurent Schwartz --- EDP et applications} (2018--2019), Exp.~II, 1--17.

\bibitem{ASZ}
J.~E. Avron, R.~Seiler, and P.~G. Zograf.
\newblock Viscosity of quantum {H}all fluids.
\newblock \emph{Phys. Rev. Lett.} 75, 697--700 (1995).

\bibitem{Read09}
N.~Read.
\newblock Non-{A}belian adiabatic statistics and {H}all viscosity in quantum {H}all states and $p_x+ip_y$ paired
superfluids.
\newblock \emph{Phys. Rev. B} 79, 045308 (2009).

\end{thebibliography}

\end{document}
